%%
% Copyright (c) 2017 - 2026, Pascal Wagler;
% Copyright (c) 2014 - 2026, John MacFarlane
%
% All rights reserved.
%
% Redistribution and use in source and binary forms, with or without
% modification, are permitted provided that the following conditions
% are met:
%
% - Redistributions of source code must retain the above copyright
% notice, this list of conditions and the following disclaimer.
%
% - Redistributions in binary form must reproduce the above copyright
% notice, this list of conditions and the following disclaimer in the
% documentation and/or other materials provided with the distribution.
%
% - Neither the name of John MacFarlane nor the names of other
% contributors may be used to endorse or promote products derived
% from this software without specific prior written permission.
%
% THIS SOFTWARE IS PROVIDED BY THE COPYRIGHT HOLDERS AND CONTRIBUTORS
% "AS IS" AND ANY EXPRESS OR IMPLIED WARRANTIES, INCLUDING, BUT NOT
% LIMITED TO, THE IMPLIED WARRANTIES OF MERCHANTABILITY AND FITNESS
% FOR A PARTICULAR PURPOSE ARE DISCLAIMED. IN NO EVENT SHALL THE
% COPYRIGHT OWNER OR CONTRIBUTORS BE LIABLE FOR ANY DIRECT, INDIRECT,
% INCIDENTAL, SPECIAL, EXEMPLARY, OR CONSEQUENTIAL DAMAGES (INCLUDING,
% BUT NOT LIMITED TO, PROCUREMENT OF SUBSTITUTE GOODS OR SERVICES;
% LOSS OF USE, DATA, OR PROFITS; OR BUSINESS INTERRUPTION) HOWEVER
% CAUSED AND ON ANY THEORY OF LIABILITY, WHETHER IN CONTRACT, STRICT
% LIABILITY, OR TORT (INCLUDING NEGLIGENCE OR OTHERWISE) ARISING IN
% ANY WAY OUT OF THE USE OF THIS SOFTWARE, EVEN IF ADVISED OF THE
% POSSIBILITY OF SUCH DAMAGE.
%%

%%
% This is the Eisvogel pandoc LaTeX template.
%
% For usage information and examples visit the official GitHub page:
% https://github.com/Wandmalfarbe/pandoc-latex-template
%%
% Options for packages loaded elsewhere
\PassOptionsToPackage{unicode}{hyperref}
\PassOptionsToPackage{hyphens}{url}
\PassOptionsToPackage{dvipsnames,svgnames,x11names,table}{xcolor}
\PassOptionsToPackage{space}{xeCJK}
\documentclass[
  paper=a4,
  ,captions=tableheading
]{scrartcl}
\usepackage{xcolor}
\usepackage[margin=2.5cm,includehead=true,includefoot=true,centering,]{geometry}
\usepackage{amsmath,amssymb}

% add backlinks to footnote references, cf. https://tex.stackexchange.com/questions/302266/make-footnote-clickable-both-ways
\usepackage{footnotebackref}
\setcounter{secnumdepth}{5}
\usepackage{iftex}
\ifPDFTeX
  \usepackage[T1]{fontenc}
  \usepackage[utf8]{inputenc}
  \usepackage{textcomp} % provide euro and other symbols
\else % if luatex or xetex
  \usepackage{unicode-math} % this also loads fontspec
  \defaultfontfeatures{Scale=MatchLowercase}
  \defaultfontfeatures[\rmfamily]{Ligatures=TeX,Scale=1}
\fi
\usepackage{lmodern}
\ifPDFTeX\else
  % xetex/luatex font selection
  \setmainfont[]{DejaVu Sans}
  \setmonofont[]{DejaVu Sans Mono}
  \ifXeTeX
    \usepackage{xeCJK}
    \setCJKmainfont[]{Noto Serif CJK SC}
  \fi
  \ifLuaTeX
    \usepackage[]{luatexja-fontspec}
    \setmainjfont[]{Noto Serif CJK SC}
  \fi
\fi
% Use upquote if available, for straight quotes in verbatim environments
\IfFileExists{upquote.sty}{\usepackage{upquote}}{}
\IfFileExists{microtype.sty}{% use microtype if available
  \usepackage[]{microtype}
  \UseMicrotypeSet[protrusion]{basicmath} % disable protrusion for tt fonts
}{}

% Use setspace anyway because we change the default line spacing.
% The spacing is changed early to affect the titlepage and the TOC.
\usepackage{setspace}
\setstretch{1.2}
\makeatletter
\@ifundefined{KOMAClassName}{% if non-KOMA class
  \IfFileExists{parskip.sty}{%
    \usepackage{parskip}
  }{% else
    \setlength{\parindent}{0pt}
    \setlength{\parskip}{6pt plus 2pt minus 1pt}}
}{% if KOMA class
  \KOMAoptions{parskip=half}}
\makeatother
\usepackage{color}
\usepackage{fancyvrb}
\newcommand{\VerbBar}{|}
\newcommand{\VERB}{\Verb[commandchars=\\\{\}]}
\DefineVerbatimEnvironment{Highlighting}{Verbatim}{commandchars=\\\{\}}
% Add ',fontsize=\small' for more characters per line
\newenvironment{Shaded}{}{}
\newcommand{\AlertTok}[1]{\textcolor[rgb]{1.00,0.00,0.00}{\textbf{#1}}}
\newcommand{\AnnotationTok}[1]{\textcolor[rgb]{0.38,0.63,0.69}{\textbf{\textit{#1}}}}
\newcommand{\AttributeTok}[1]{\textcolor[rgb]{0.49,0.56,0.16}{#1}}
\newcommand{\BaseNTok}[1]{\textcolor[rgb]{0.25,0.63,0.44}{#1}}
\newcommand{\BuiltInTok}[1]{\textcolor[rgb]{0.00,0.50,0.00}{#1}}
\newcommand{\CharTok}[1]{\textcolor[rgb]{0.25,0.44,0.63}{#1}}
\newcommand{\CommentTok}[1]{\textcolor[rgb]{0.38,0.63,0.69}{\textit{#1}}}
\newcommand{\CommentVarTok}[1]{\textcolor[rgb]{0.38,0.63,0.69}{\textbf{\textit{#1}}}}
\newcommand{\ConstantTok}[1]{\textcolor[rgb]{0.53,0.00,0.00}{#1}}
\newcommand{\ControlFlowTok}[1]{\textcolor[rgb]{0.00,0.44,0.13}{\textbf{#1}}}
\newcommand{\DataTypeTok}[1]{\textcolor[rgb]{0.56,0.13,0.00}{#1}}
\newcommand{\DecValTok}[1]{\textcolor[rgb]{0.25,0.63,0.44}{#1}}
\newcommand{\DocumentationTok}[1]{\textcolor[rgb]{0.73,0.13,0.13}{\textit{#1}}}
\newcommand{\ErrorTok}[1]{\textcolor[rgb]{1.00,0.00,0.00}{\textbf{#1}}}
\newcommand{\ExtensionTok}[1]{#1}
\newcommand{\FloatTok}[1]{\textcolor[rgb]{0.25,0.63,0.44}{#1}}
\newcommand{\FunctionTok}[1]{\textcolor[rgb]{0.02,0.16,0.49}{#1}}
\newcommand{\ImportTok}[1]{\textcolor[rgb]{0.00,0.50,0.00}{\textbf{#1}}}
\newcommand{\InformationTok}[1]{\textcolor[rgb]{0.38,0.63,0.69}{\textbf{\textit{#1}}}}
\newcommand{\KeywordTok}[1]{\textcolor[rgb]{0.00,0.44,0.13}{\textbf{#1}}}
\newcommand{\NormalTok}[1]{#1}
\newcommand{\OperatorTok}[1]{\textcolor[rgb]{0.40,0.40,0.40}{#1}}
\newcommand{\OtherTok}[1]{\textcolor[rgb]{0.00,0.44,0.13}{#1}}
\newcommand{\PreprocessorTok}[1]{\textcolor[rgb]{0.74,0.48,0.00}{#1}}
\newcommand{\RegionMarkerTok}[1]{#1}
\newcommand{\SpecialCharTok}[1]{\textcolor[rgb]{0.25,0.44,0.63}{#1}}
\newcommand{\SpecialStringTok}[1]{\textcolor[rgb]{0.73,0.40,0.53}{#1}}
\newcommand{\StringTok}[1]{\textcolor[rgb]{0.25,0.44,0.63}{#1}}
\newcommand{\VariableTok}[1]{\textcolor[rgb]{0.10,0.09,0.49}{#1}}
\newcommand{\VerbatimStringTok}[1]{\textcolor[rgb]{0.25,0.44,0.63}{#1}}
\newcommand{\WarningTok}[1]{\textcolor[rgb]{0.38,0.63,0.69}{\textbf{\textit{#1}}}}

% Workaround/bugfix from jannick0.
% See https://github.com/jgm/pandoc/issues/4302#issuecomment-360669013)
% or https://github.com/Wandmalfarbe/pandoc-latex-template/issues/2
%
% Redefine the verbatim environment 'Highlighting' to break long lines (with
% the help of fvextra). Redefinition is necessary because it is unlikely that
% pandoc includes fvextra in the default template.
\usepackage{fvextra}
\DefineVerbatimEnvironment{Highlighting}{Verbatim}{breaklines,fontsize=\footnotesize,commandchars=\\\{\}}

\usepackage{longtable,booktabs,array}
\newcounter{none} % for unnumbered tables
\usepackage{calc} % for calculating minipage widths
% Correct order of tables after \paragraph or \subparagraph
\usepackage{etoolbox}
\makeatletter
\patchcmd\longtable{\par}{\if@noskipsec\mbox{}\fi\par}{}{}
\makeatother
% Allow footnotes in longtable head/foot
\IfFileExists{footnotehyper.sty}{\usepackage{footnotehyper}}{\usepackage{footnote}}
\makesavenoteenv{longtable}
\setlength{\emergencystretch}{3em} % prevent overfull lines
\providecommand{\tightlist}{%
  \setlength{\itemsep}{0pt}\setlength{\parskip}{0pt}}
\usepackage{tikz}
\usetikzlibrary{arrows.meta,positioning,fit,backgrounds}
\usepackage{bookmark}
\IfFileExists{xurl.sty}{\usepackage{xurl}}{} % add URL line breaks if available
\urlstyle{same}
\definecolor{default-linkcolor}{HTML}{A50000}
\definecolor{default-filecolor}{HTML}{A50000}
\definecolor{default-citecolor}{HTML}{4077C0}
\definecolor{default-urlcolor}{HTML}{4077C0}

\hypersetup{
  pdftitle={Agentic RL on Terminal-Bench with slime and AgentCore Runtime},
  pdfauthor={Kaige Yang --- Senior Applied Scientist, AWS},
  colorlinks=true,
  linkcolor={default-linkcolor},
  filecolor={default-filecolor},
  citecolor={default-citecolor},
  urlcolor={default-urlcolor},
  breaklinks=true,
  pdfcreator={LaTeX via pandoc with the Eisvogel template}}

\title{Agentic RL on Terminal-Bench with slime and AgentCore Runtime}
\usepackage{etoolbox}
\makeatletter
\providecommand{\subtitle}[1]{% add subtitle to \maketitle
  \apptocmd{\@title}{\par {\large #1 \par}}{}{}
}
\makeatother
\subtitle{Hands-on guidance, grounded in a working Terminal-Bench RL
demo}
\author{Kaige Yang --- Senior Applied Scientist, AWS}
\date{September 30, 2026}


%
% for the background color of the title page
%
\usepackage{pagecolor}
\usepackage{afterpage}
\usepackage[margin=2.5cm,includehead=true,includefoot=true,centering]{geometry}

%
% break urls
%
\PassOptionsToPackage{hyphens}{url}

%
% When using babel or polyglossia with biblatex, loading csquotes is recommended
% to ensure that quoted texts are typeset according to the rules of your main language.
%
\usepackage{csquotes}

%
% captions
%
\definecolor{caption-color}{HTML}{777777}
\usepackage[font={stretch=1.2}, textfont={color=caption-color}, position=top, skip=4mm, labelfont=bf, singlelinecheck=false, justification=raggedright]{caption}
\setcapindent{0em}

%
% blockquote
%
\definecolor{blockquote-border}{RGB}{221,221,221}
\definecolor{blockquote-text}{RGB}{119,119,119}
\usepackage{mdframed}
\newmdenv[rightline=false,bottomline=false,topline=false,linewidth=3pt,linecolor=blockquote-border,skipabove=\parskip]{customblockquote}
\renewenvironment{quote}{\begin{customblockquote}\list{}{\rightmargin=0em\leftmargin=0em}%
\item\relax\color{blockquote-text}\ignorespaces}{\unskip\unskip\endlist\end{customblockquote}}

%
% Source Sans as the default font family
% Source Code Pro for monospace text
%
\ifnum 0\ifxetex 1\fi\ifluatex 1\fi=0 % if pdftex
    \usepackage[default]{sourcesans}
  \usepackage{sourcecodepro}
  \else % if not pdftex (e.g. lualatex)
    \fi

%
% heading color
%
\definecolor{heading-color}{RGB}{40,40,40}
% By default, the KOMA-Script classes will typeset sectioning headings in
% sans-serif. Use the normal body font for headings.
\addtokomafont{disposition}{\normalfont\color{heading-color}\bfseries}

%
% variables for title, author and date
%
\usepackage{titling}
\title{Agentic RL on Terminal-Bench with slime and AgentCore Runtime}
\author{Kaige Yang --- Senior Applied Scientist, AWS}
\date{September 30, 2026}

%
% tables
%

\definecolor{table-row-color}{HTML}{F5F5F5}
\definecolor{table-rule-color}{HTML}{999999}

%\arrayrulecolor{black!40}
\arrayrulecolor{table-rule-color}     % color of \toprule, \midrule, \bottomrule
\setlength\heavyrulewidth{0.3ex}      % thickness of \toprule, \bottomrule
\renewcommand{\arraystretch}{1.3}     % spacing (padding)


%
%
% Listings
%
%


%
% header and footer
%
\usepackage[headsepline,footsepline]{scrlayer-scrpage}

\newpairofpagestyles{eisvogel-header-footer}{
  \clearpairofpagestyles
  \ihead*{Terminal-Bench RL on AgentCore}
  \chead*{}
  \ohead*{September 30, 2026}
  \ifoot*{Kaige Yang --- Senior Applied Scientist, AWS}
  \cfoot*{}
  \ofoot*{\thepage}
  \addtokomafont{pageheadfoot}{\upshape}
}
\pagestyle{eisvogel-header-footer}



%
% Define watermark
%

\begin{document}

\begin{titlepage}
\newgeometry{left=6cm}
\newcommand{\colorRule}[3][black]{\textcolor[HTML]{#1}{\rule{#2}{#3}}}
\begin{flushleft}
\noindent
\\[-1em]
\color[HTML]{5F5F5F}
\makebox[0pt][l]{\colorRule[5F5F5F]{1.3\textwidth}{4pt}}
\par
\noindent

{
  \setstretch{1.4}
  \vfill
  \noindent {\huge \textbf{\textsf{Agentic RL on Terminal-Bench with
slime and AgentCore Runtime}}}
    \vskip 1em
  {\Large \textsf{Hands-on guidance, grounded in a working
Terminal-Bench RL demo}}
    \vskip 2em
  \noindent {\Large \textsf{Kaige Yang --- Senior Applied Scientist,
AWS}}
  \vfill
}


\textsf{September 30, 2026}
\end{flushleft}
\end{titlepage}
\restoregeometry
\pagenumbering{arabic}

% don't generate the default title
% \maketitle


{
\setcounter{tocdepth}{3}
\tableofcontents
\newpage
}
\hypertarget{introduction}{%
\section{Introduction}\label{introduction}}

This report is a \textbf{hands-on guide for AWS customers doing agentic
RL}, built around one working demo: \textbf{RL on Terminal-Bench with
slime and AgentCore Runtime} (written for
\texttt{code/demos/demo\_cpu\_terminalbench}, now
\texttt{code/demos/demo\_terminalbench\_kernelbench}; see the update
note above). The policy is Qwen3-4B, trained by slime with fully-async
GRPO. The sandbox is an \textbf{image pool of 218 AgentCore arm64
microVM Runtimes in VPC mode, one per Terminal-Bench task}, with a fresh
session per episode. The agent loop runs \textbf{inside that session},
and the reward is the task's own hidden pytest oracle. The claims are
grounded in the \textbf{official AgentCore docs} {[}D1{]} and in this
project's run logs.

\textbf{After reading it, you will know:}

\begin{itemize}
\tightlist
\item
  \textbf{The reference architecture of agentic RL}: how the trainer,
  inference engine, proxy, harness and sandbox fit together, and which
  seams AgentCore fills (§``Agentic RL and its reference
  architecture'').
\item
  \textbf{What tasks labs train in a sandbox, and how the reward works}:
  the verifiable task families, each scored by a programmatic
  \textbf{RLVR} checker rather than an LLM judge, and why Terminal-Bench
  is the hardest sandbox shape among them (§``Tasks frontier labs run in
  a sandbox'').
\item
  \textbf{What AgentCore Runtime \emph{is}}: the three concepts you must
  not conflate (\textbf{Runtime}, \textbf{session}, \textbf{microVM}),
  their \textbf{1 : N : 1} cardinality, the two compute types, the
  \textbf{\textasciitilde15-minute idle timer}, the \textbf{2 GB arm64
  image cap}, and how to build your own Runtime (§``Foundations of
  AgentCore Runtime'').
\item
  \textbf{How AgentCore plugs into slime without a fork}: one generate
  hook, a sandbox client that turns each operation into one
  \texttt{InvokeAgentRuntime}, and the agent loop in the session,
  dialling the trainer's adapter back over the VPC (§``Using AgentCore
  as the sandbox in an agentic RL workflow'').
\item
  \textbf{How to run agentic RL on hundreds of per-task images} (§``The
  Terminal-Bench demo''):

  \begin{itemize}
  \tightlist
  \item
    one image and one Runtime per task, built from the task's own
    Dockerfile;
  \item
    strict row → image → Runtime routing;
  \item
    a fresh microVM per episode;
  \item
    \textbf{look-ahead pre-open} to hide the cold start;
  \item
    a \textbf{heartbeat} to survive the trainer's weight-sync stall;
  \item
    a \textbf{harness tier} that scales the agent loop with the sandbox
    rather than inside the trainer.
  \end{itemize}
\item
  \textbf{What the runs measured and what is still open}: the 2026-09-23
  and 2026-09-30 smokes, the convergence plan (§``The Terminal-Bench
  demo in depth''), and the production open problems (§``Open problems
  in production agentic RL'').
\end{itemize}

\textbf{Scope.} AgentCore is the execution \textbf{sandbox},
\textbf{slime} is the RL trainer, and the \textbf{harness runs inside
the sandbox session}. The harness is the mini\_claude agent loop, one
per episode, and it calls back to the trainer's per-engine adapter. It
therefore scales with the sandbox, and the trainer keeps only a thin
dispatch tier. Every other demo the project built, including a SWE-bench
Verified pool, was removed on 2026-09-30; they are summarized once in
§``History: removed demos (2026-09-30)''.

\textbf{Roadmap.} The report proceeds in this order:

\begin{enumerate}
\def\labelenumi{\arabic{enumi}.}
\tightlist
\item
  the \textbf{reference architecture} of agentic RL and how slime drives
  it (§2);
\item
  the \textbf{verifiable task families} and their rewards (§3);
\item
  the \textbf{foundations of AgentCore Runtime} (§4);
\item
  \textbf{wiring AgentCore into slime} as the sandbox (§5);
\item
  \textbf{the Terminal-Bench demo} (§6) and its \textbf{in-depth
  results} (§7);
\item
  the removed demos (§8);
\item
  the \textbf{open problems} (§9) and a \textbf{summary} (§10).
\end{enumerate}

\hypertarget{agentic-rl-and-its-reference-architecture}{%
\section{Agentic RL and its reference
architecture}\label{agentic-rl-and-its-reference-architecture}}

\textbf{What is agentic RL?} In \emph{agentic} reinforcement learning
the model under training is not asked for a single answer --- it acts as
an \textbf{agent} over many turns against an \textbf{executable
environment}: it writes code, runs a command, reads the result, and
tries again. Only when the trajectory ends does a \textbf{programmatic
checker} turn the outcome into a scalar reward (tests pass, result-set
matches, flag verified) --- \textbf{no LLM judge} {[}1, 2{]}. This is
reinforcement learning with \emph{verifiable} rewards (RLVR): the
environment, not a model, decides who was right.

Training that loop end-to-end uses a \textbf{disaggregated} design
(following Agent Lightning {[}1{]}): a training side (trainer +
inference engine) and an agent side (harness + sandbox), joined by an
LLM-proxy / rollout-collection gateway. The reference architecture below
traces the full cycle --- a \textbf{generation path} across the top and
a \textbf{training-data path} back along the bottom.

\begin{figure}[htbp]
\centering
\resizebox{\textwidth}{!}{%
\begin{tikzpicture}[
  font=\small,
  box/.style={draw, rounded corners, minimum height=1.15cm, minimum width=2.4cm, align=center, thick, fill=white},
  gpu/.style={box, fill=blue!8},
  cpu/.style={box, fill=black!4},
  focus/.style={box, fill=orange!18},
  lbl/.style={font=\scriptsize, align=center, fill=white, inner sep=1.5pt},
  ptitle/.style={font=\scriptsize\bfseries, text=black!55},
  panel/.style={draw, rounded corners, dashed, thick, inner sep=0.34cm},
  >={Stealth[]},
]
% --- Training server (left): generation top row, training-data bottom row ---
\node[gpu] (trainer) at (0,0)    {\textbf{Training Engine}\\[1pt]\scriptsize Megatron};
\node[gpu] (infer)   at (5.0,0)  {\textbf{Inference Engine}\\[1pt]\scriptsize SGLang};
\node[cpu] (buffer)  at (0,-3.4) {\textbf{Buffer}\\[1pt]\scriptsize training samples};
\node[cpu] (adapter) at (5.0,-3.4){\textbf{Sample Adapter}\\[1pt]\scriptsize assemble + reward};
% --- API gateway (middle) ---
\node[cpu] (proxy)   at (10.0,0)   {\textbf{LLM API Proxy}\\[1pt]\scriptsize policy under training};
\node[cpu] (rollout) at (10.0,-3.4){\textbf{Rollout API}\\[1pt]\scriptsize collect trajectory};
% --- Agent runtime (right) ---
\node[cpu]   (harness) at (15.0,0) {\textbf{Harness}\\[1pt]\scriptsize agent loop};
\node[focus] (sandbox) at (20.0,0){\textbf{Sandbox}\\[1pt]\scriptsize AgentCore:\\tools + checker};

% generation lane (top, left -> right)
\draw[<->, thick] (trainer) -- node[lbl,above]{weight sync} (infer);
\draw[->, thick]  (infer)   -- node[lbl,above]{serve} (proxy);
\draw[->, thick]  (proxy)   -- node[lbl,above]{turn} (harness);
\draw[<->, thick] (harness) -- node[lbl,above]{tool call / reward} (sandbox);
% training-data lane (bottom, right -> left)
\draw[->, thick]  (harness) -- node[lbl,pos=0.32]{rollout} (rollout);
\draw[->, thick]  (rollout) -- node[lbl,above]{trajectory} (adapter);
\draw[->, thick]  (adapter) -- node[lbl,above]{samples} (buffer);
\draw[->, thick]  (buffer)  -- node[lbl,left]{sample batch} (trainer);

% domain panels (drawn behind the boxes)
\begin{scope}[on background layer]
  \node[panel, fill=blue!3,  fit=(trainer)(infer)(buffer)(adapter)] (ptrain) {};
  \node[panel, fill=black!3, fit=(proxy)(rollout)] (pgw) {};
  \node[panel, fill=orange!4, fit=(harness)(sandbox)] (pagent) {};
\end{scope}
\node[ptitle, anchor=south west] at ([yshift=2pt]ptrain.north west) {Training server (GPU) --- slime};
\node[ptitle, anchor=south] at ([yshift=2pt]pgw.north) {API gateway};
\node[ptitle, anchor=south west] at ([yshift=2pt]pagent.north west) {Agent runtime (CPU + sandbox)};
\end{tikzpicture}%
}
\caption{Reference architecture for (harnessed) agentic RL, following the disaggregated design of Agent Lightning [1] with slime [9] as the training server. Three domains: a \textbf{training server} (GPU), an \textbf{API gateway}, and an \textbf{agent runtime}. \textbf{Generation lane (top, $\rightarrow$):} the training engine syncs weights to the inference engine; the LLM API proxy fronts inference so every turn is served by the policy under training; the harness drives the agent loop; the sandbox (AgentCore) runs each tool call and the checker. \textbf{Training-data lane (bottom, $\leftarrow$):} the rollout is collected via the Rollout API, the Sample Adapter assembles it from the recorded token IDs (no re-tokenization --- the TITO property) and attaches the reward to form training samples, which accumulate in the Buffer; the trainer draws a sample batch and updates --- closing the loop.}
\end{figure}

\textbf{The core components.}

\begin{itemize}
\tightlist
\item
  \textbf{Trainer} --- the learner (here \textbf{slime} driving
  \textbf{Megatron}, on GPU). It samples prompts, collects finished
  trajectories with their rewards, computes the policy-gradient update,
  and periodically pushes fresh weights to the inference engine.
\item
  \textbf{Inference engine} --- the generation server (here
  \textbf{SGLang}, on GPU) that produces the model's turns at rollout
  speed. It holds the \emph{current} policy weights; the trainer
  refreshes them via \textbf{weight sync} so rollouts always reflect the
  latest policy.
\item
  \textbf{Proxy} --- a thin adapter (CPU) fronting the inference engine
  so that every model call the harness makes is answered by the
  \textbf{policy under training} (not some external API). It is the seam
  that lets an off-the-shelf agent loop drive an in-training model.
\item
  \textbf{Harness} --- the \textbf{agent loop} (CPU): it holds the
  conversation, calls the proxy for each turn, extracts the model's
  proposed action (a code block, a tool call), submits it to the
  sandbox, and feeds the result back as the next observation until the
  episode ends.
\item
  \textbf{Sandbox} --- the \textbf{isolated execution environment} (CPU,
  or GPU for CUDA work) where the model's untrusted output actually runs
  (write → run → read) and where the \textbf{checker scores the reward}.
  It must be ad-hoc and isolated per episode.
\end{itemize}

\textbf{The training-data path.} The bottom of Figure 1 shows how a
finished episode becomes a gradient step. The \textbf{Rollout API}
collects the trajectory (the sequence of LLM request--response pairs the
harness made); the \textbf{Sample Adapter} assembles those calls from
their recorded token IDs (no re-tokenization ---
§``Token-in-token-out'') and attaches the episode's reward and
advantage, turning one rollout into one or more \textbf{training
samples}; those land in the \textbf{Buffer}; and the trainer draws a
\textbf{sample batch} from the buffer for each update. In
\emph{harnessed} agentic RL the training engine never sees the
environment directly --- only this stream of samples reconstructed from
the proxy's request/response log {[}1{]}.

\hypertarget{where-agentic-rl-is-used}{%
\subsection{Where agentic RL is used}\label{where-agentic-rl-is-used}}

The same loop --- act in an environment, score with a checker --- covers
a wide range of tasks. The common thread is a \textbf{verifiable
reward}: an environment (not a human or an LLM judge) decides success.
Representative use cases:

\begin{itemize}
\tightlist
\item
  \textbf{Software engineering (SWE)} --- edit a real repository to
  resolve an issue; reward = the project's test suite turns green
  (e.g.~SWE-bench {[}6{]}).
\item
  \textbf{Code generation \& competitive programming} --- write a
  function that passes hidden unit tests (e.g.~MBPP {[}8{]}).
\item
  \textbf{GPU-kernel / systems} --- write a CUDA kernel and optimize it;
  reward = numerical \texttt{allclose} plus a speedup target
  (e.g.~KernelBench {[}7{]}).
\item
  \textbf{Math \& logical reasoning} --- produce an answer checked for
  exact/numeric equality against ground truth.
\item
  \textbf{Tool use / function calling} --- orchestrate multi-step API or
  tool calls; reward = correct final result \emph{and} the right tools
  invoked.
\item
  \textbf{SQL \& data analysis} --- generate a query whose result-set
  matches the expected rows.
\item
  \textbf{RAG \& search agents} --- retrieve then answer; reward =
  answer grounded in the retrieved facts.
\item
  \textbf{Web \& computer use} --- drive a browser or GUI to complete a
  task.
\item
  \textbf{Security / CTF} --- exploit a challenge and capture the flag;
  reward = the checker validates the flag.
\item
  \textbf{Multi-agent \& long-horizon workflows} --- several agents or
  many turns cooperate toward a verifiable end state.
\end{itemize}

Every one of these needs an \textbf{isolated place to execute untrusted
output and run the checker} --- the sandbox seam this report is about.

\hypertarget{the-harness-the-agent-loop-that-owns-the-episode}{%
\subsection{The harness --- the agent loop that owns the
episode}\label{the-harness-the-agent-loop-that-owns-the-episode}}

\textbf{What it is.} The harness is the \textbf{agent loop}: the
orchestration code that runs \emph{one} episode end to end. In this
project it is the tool-calling \textbf{mini\_claude} loop
(\texttt{write\_file} / \texttt{bash} / \texttt{read\_file} over the
Anthropic Messages API, \texttt{mini\_agent.py}), and it is
\textbf{task-agnostic} --- the same loop serves every task family; only
the task's prompt, seed files, and checker change. It runs
\textbf{inside the episode's own sandbox session}:
\texttt{in\_session\_harness.py} writes a stdlib runner into the
session, starts it and long-polls it, and every model turn calls back to
the trainer's \texttt{AnthropicAdapter}, which records the exact tokens.
Two reasons put it there rather than in the trainer. The harness has to
scale \textbf{along with} inference and the sandbox, or a few coroutines
silently become the throughput ceiling. And the unit that already scales
per episode is the sandbox session, so one loop per session scales with
it --- DSec's ``move the agent loop out of the trainer'' {[}31{]}.

\textbf{How it works with the other components.} The harness sits
between the proxy and the sandbox and is the only component that sees a
whole episode:

\begin{itemize}
\tightlist
\item
  \textbf{↔ Proxy:} for each turn it sends the running conversation to
  the proxy and receives the model's response (the policy under
  training).
\item
  \textbf{↔ Sandbox:} it extracts the model's proposed action (a code
  block, a tool call) and executes it in the sandbox, then reads back
  stdout / exit-code as the next observation.
\item
  \textbf{↔ Task / checker:} it seeds the workspace from the task and,
  at the end, runs the checker (in the sandbox) to get the reward.
\item
  \textbf{→ Rollout / Sample Adapter:} the trajectory it produces (the
  proxy-logged calls + the final reward) becomes the training sample.
\end{itemize}

\begin{figure}[htbp]
\centering
\resizebox{0.95\textwidth}{!}{%
\begin{tikzpicture}[
  font=\small,
  box/.style={draw, rounded corners, minimum height=1.15cm, minimum width=2.7cm, align=center, thick, fill=white},
  svc/.style={box, fill=blue!8},
  focus/.style={box, fill=orange!18},
  stage/.style={box, fill=black!4},
  lbl/.style={font=\scriptsize, align=center, fill=white, inner sep=1.5pt},
  ptitle/.style={font=\scriptsize\bfseries, text=black!55},
  panel/.style={draw, rounded corners, dashed, thick, inner sep=0.34cm},
  >={Stealth[]},
]
% main lane: seed -> ask -> run -> reward
\node[stage]  (seed)   at (0,0)     {\textbf{Seed}\\[1pt]\scriptsize PROBLEM.md +\\run\_checks.py};
\node[stage]  (ask)    at (4.0,0)   {\textbf{Ask}\\[1pt]\scriptsize model turn\\(parse action)};
\node[stage]  (run)    at (8.6,0)   {\textbf{Run}\\[1pt]\scriptsize exec action +\\checker};
\node[stage]  (reward) at (13.2,0)  {\textbf{Reward}\\[1pt]\scriptsize 1.0 / 0.0 $\rightarrow$\\Sample};
% services the loop calls
\node[svc]   (proxy)  at (4.0,3.5) {\textbf{Proxy} $\rightarrow$ policy\\[1pt]\scriptsize SGLang (in training)};
\node[focus] (sbx)    at (8.6,-2.8){\textbf{AgentCore sandbox}\\[1pt]\scriptsize tools + checker};

% main flow
\draw[->, thick] (seed) -- (ask);
\draw[->, thick] (ask)  -- node[lbl,above]{action} (run);
\draw[->, thick] (run)  -- node[lbl,above]{solved} (reward);
% calls out to services
\draw[<->, thick] (ask) -- node[lbl,left]{turn} (proxy);
\draw[<->, thick] (run) -- node[lbl,right]{write+exec / stdout} (sbx);
% feedback loop: not solved -> next turn (high arc, endpoints offset from the proxy connector)
\draw[->, thick] (run.north) .. controls +(0,2.0) and +(0,2.0) ..
      node[lbl,above,pos=0.5]{observation $\rightarrow$ next turn (if unsolved)} (ask.north east);

% turn-loop panel
\begin{scope}[on background layer]
  \node[panel, fill=black!3, fit=(ask)(run)] (ploop) {};
\end{scope}
\node[ptitle, anchor=north west] at ([yshift=-2pt]ploop.south west) {turn loop ($\leq$ max\_turns)};
\end{tikzpicture}%
}
\caption{The harness's episode loop (\texttt{run\_episode}). It \textbf{seeds} the sandbox once, then repeats a \textbf{turn loop}: ask the proxy for a model turn, parse the proposed action, run it (and the checker) in the AgentCore sandbox, and feed the observation back for the next turn. The loop exits when the checker passes or \texttt{max\_turns} is reached; the final \textbf{reward} and the loss-masked trajectory become one training \texttt{Sample}. The proxy (policy under training) and the sandbox are the only two services the loop calls each turn.}
\end{figure}

\textbf{Step-by-step example --- one MBPP episode from the harness's
view} (a simplified sketch of the loop's logic; the shipped loop is the
tool-calling mini\_claude, run inside the sandbox session).

\begin{Shaded}
\begin{Highlighting}[]
\ControlFlowTok{async} \KeywordTok{def}\NormalTok{ run\_episode(task, policy, sandbox):          }\CommentTok{\# one episode}
    \CommentTok{\# 1. SEED: write the problem + the checker into the fresh sandbox}
    \ControlFlowTok{for}\NormalTok{ name, content }\KeywordTok{in}\NormalTok{ task.setup\_files():           }\CommentTok{\# e.g. PROBLEM.md}
        \ControlFlowTok{await}\NormalTok{ sandbox.write\_file(name, content)}
    \ControlFlowTok{await}\NormalTok{ sandbox.write\_file(}\StringTok{"run\_checks.py"}\NormalTok{, task.runner\_source())}

\NormalTok{    messages }\OperatorTok{=}\NormalTok{ task.initial\_messages()                 }\CommentTok{\# system + PROBLEM.md}
\NormalTok{    solved }\OperatorTok{=} \VariableTok{False}
    \ControlFlowTok{for}\NormalTok{ \_ }\KeywordTok{in} \BuiltInTok{range}\NormalTok{(task.max\_turns):                    }\CommentTok{\# 2. TURN LOOP}
\NormalTok{        text }\OperatorTok{=} \ControlFlowTok{await}\NormalTok{ policy.complete(messages)         }\CommentTok{\#    (a) ask proxy {-}\textgreater{} policy turn}
\NormalTok{        code }\OperatorTok{=}\NormalTok{ task.extract\_submission(text)           }\CommentTok{\#    (b) parse the action (fenced code)}
        \ControlFlowTok{await}\NormalTok{ sandbox.write\_file(}\StringTok{"solution.py"}\NormalTok{, code)  }\CommentTok{\#    (c) place the artifact}
\NormalTok{        exit\_code, out, err }\OperatorTok{=} \ControlFlowTok{await}\NormalTok{ sandbox.}\BuiltInTok{exec}\NormalTok{(       }\CommentTok{\#    (d) run the checker IN the sandbox}
            \StringTok{"python3 run\_checks.py"}\NormalTok{, timeout}\OperatorTok{=}\DecValTok{60}\NormalTok{)}
        \ControlFlowTok{if}\NormalTok{ exit\_code }\OperatorTok{==} \DecValTok{0}\NormalTok{:                              }\CommentTok{\#    (e) checker passed?}
\NormalTok{            solved }\OperatorTok{=} \VariableTok{True}
            \ControlFlowTok{break}
\NormalTok{        messages.append(task.observation(out, err))    }\CommentTok{\#    (f) feed the failure back, loop}

\NormalTok{    reward }\OperatorTok{=} \FloatTok{1.0} \ControlFlowTok{if}\NormalTok{ solved }\ControlFlowTok{else} \FloatTok{0.0}                    \CommentTok{\# 3. REWARD (from the checker)}
    \ControlFlowTok{return}\NormalTok{ build\_token\_stream(messages, reward)        }\CommentTok{\# 4. loss{-}masked trajectory {-}\textgreater{} Sample}
\end{Highlighting}
\end{Shaded}

Walking it: the harness \textbf{seeds} the sandbox (step 1), then
repeats \textbf{ask → parse → run → observe} (step 2a--f), each ``ask''
going through the proxy to the in-training policy and each ``run'' going
to the AgentCore sandbox. When the checker exits \texttt{0} the episode
is \textbf{solved}; otherwise the failure output is appended and the
model tries again, up to \texttt{max\_turns}. The final \textbf{reward}
(step 3) comes from the checker --- not a judge --- and the harness
hands the loss-masked trajectory to the training-data path (step 4). The
harness owns the loop; the proxy and sandbox are the two services it
calls each turn.

\hypertarget{the-proxy-every-model-call-is-the-policy-under-training}{%
\subsection{The proxy --- every model call is the policy under
training}\label{the-proxy-every-model-call-is-the-policy-under-training}}

\textbf{What it is.} The proxy is a thin CPU service that exposes a
\textbf{standard LLM API} (in this project an
Anthropic-Messages--compatible \texttt{AnthropicAdapter}, but an
OpenAI-compatible endpoint works the same) and routes every request to
the \textbf{in-training inference engine} (SGLang) instead of to any
external provider. It is the single seam that lets an
\emph{off-the-shelf agent harness} --- one written to call a normal LLM
API --- drive a model that is still being trained.

\textbf{What it does.} Three jobs, all on the critical path of every
turn:

\begin{enumerate}
\def\labelenumi{\arabic{enumi}.}
\tightlist
\item
  \textbf{Substitute the endpoint.} The harness thinks it is calling a
  hosted LLM; the proxy answers with the \textbf{policy under training},
  holding the latest synced weights. No harness change is needed to
  point it at the trainer.
\item
  \textbf{Record the exact tokens.} For each call the proxy captures the
  \textbf{prompt token IDs and the sampled response token IDs} (and
  log-probs). The training engine only ever observes this stream of
  request/response pairs --- so recording the \emph{exact} sampled
  tokens is what avoids the retokenization/decode-drift problem in
  harnessed agentic RL {[}1{]}.
\item
  \textbf{Bucket by episode.} Each request is tagged with the episode's
  \texttt{session\_id}, so all calls of one rollout are grouped into the
  right trajectory before the Sample Adapter turns them into training
  samples.
\end{enumerate}

\begin{figure}[htbp]
\centering
\resizebox{0.95\textwidth}{!}{%
\begin{tikzpicture}[
  font=\small,
  box/.style={draw, rounded corners, minimum height=1.15cm, minimum width=2.8cm, align=center, thick, fill=white},
  svc/.style={box, fill=blue!8},
  focus/.style={box, fill=orange!18},
  cpu/.style={box, fill=black!4},
  lbl/.style={font=\scriptsize, align=center, fill=white, inner sep=1.5pt},
  >={Stealth[]},
]
\node[cpu]   (harness) at (0,0)     {\textbf{Harness}\\[1pt]\scriptsize agent loop (text)};
\node[focus] (proxy)   at (5.5,0)   {\textbf{Proxy}\\[1pt]\scriptsize LLM API $\rightarrow$ policy};
\node[svc]   (infer)   at (11.0,0)  {\textbf{Inference engine}\\[1pt]\scriptsize SGLang (GPU)};
\node[cpu]   (store)   at (5.5,-2.7) {\textbf{Rollout store}\\[1pt]\scriptsize per \texttt{session\_id}};
\node[cpu]   (adapter) at (11.0,-2.7){\textbf{Sample Adapter}\\[1pt]\scriptsize + reward $\rightarrow$ trainer};

% request / response (text side)
\draw[->, thick] ([yshift=4pt]harness.east) -- node[lbl,above]{POST /v1/messages (text)} ([yshift=4pt]proxy.west);
\draw[->, thick] ([yshift=-4pt]proxy.west)  -- node[lbl,below]{Anthropic-style reply} ([yshift=-4pt]harness.east);
% token side
\draw[->, thick] ([yshift=4pt]proxy.east) -- node[lbl,above]{prompt\_ids} ([yshift=4pt]infer.west);
\draw[->, thick] ([yshift=-4pt]infer.west) -- node[lbl,below]{response\_ids, logprobs} ([yshift=-4pt]proxy.east);
% record + hand off
\draw[->, thick] (proxy) -- node[lbl,right]{record exact tokens} (store);
\draw[->, thick] (store) -- node[lbl,above]{on episode end} (adapter);
\end{tikzpicture}%
}
\caption{The proxy on the seam between the text-based harness and the token-based inference engine. It \textbf{substitutes} the endpoint (the harness's ordinary LLM API call is answered by the policy under training), and it \textbf{records} the exact prompt and response token IDs (with log-probs) into a per-\texttt{session\_id} rollout store. On episode end the Sample Adapter reads that store, attaches the reward, and emits training samples --- so the trainer learns from the exact tokens that were sampled --- the \emph{token-in-token-out} contract discussed next.}
\end{figure}

\textbf{End-to-end example --- one turn through the proxy.} Consider a
single MBPP coding turn:

\begin{Shaded}
\begin{Highlighting}[]
\NormalTok{1. Harness builds the chat and calls the proxy (drop{-}in LLM API):}
\NormalTok{     POST /v1/messages   \{ model: "policy",}
\NormalTok{                           messages: [system, PROBLEM.md, prior observations],}
\NormalTok{                           metadata: \{ session\_id: "ep\_9f3c…" \} \}}

\NormalTok{2. Proxy → Inference engine (SGLang, current policy weights):}
\NormalTok{     generate(prompt\_ids) → response\_ids, text, logprobs}

\NormalTok{3. Proxy records the on{-}policy call for training:}
\NormalTok{     rollout\_store[session\_id].append(}
\NormalTok{        \{ prompt\_ids, response\_ids, logprobs \} )   \# exact sampled tokens}

\NormalTok{4. Proxy returns an Anthropic{-}style response to the harness:}
\NormalTok{     \{ content: "\textasciigrave{}\textasciigrave{}\textasciigrave{}python\textbackslash{}n def solve(): …\textasciigrave{}\textasciigrave{}\textasciigrave{}", stop\_reason: "end\_turn" \}}

\NormalTok{5. Harness extracts solution.py, runs it in the AgentCore sandbox,}
\NormalTok{   appends the checker\textquotesingle{}s stdout/exit{-}code as the next observation, loops.}

\NormalTok{6. Episode ends → reward = 1.0/0.0. Sample Adapter reads}
\NormalTok{   rollout\_store["ep\_9f3c…"], attaches the reward/advantage, and emits}
\NormalTok{   loss{-}masked training samples (mask 0 on prompt/observation tokens,}
\NormalTok{   1 on the policy\textquotesingle{}s own response tokens).}
\end{Highlighting}
\end{Shaded}

The key point: \textbf{without the proxy, the trainer could not observe
--- or take credit for --- the tokens the harness actually sampled.} The
proxy is simultaneously the \emph{substitution} point (external API →
policy) and the \emph{observation} point (turns → training samples),
which is exactly why the disaggregated design in Figure 1 puts it on the
seam between the agent runtime and the training server.

\hypertarget{token-in-token-out-tito-the-proxys-core-responsibility}{%
\subsection{Token-in-token-out (TITO) --- the proxy's core
responsibility}\label{token-in-token-out-tito-the-proxys-core-responsibility}}

\textbf{What it is, and where it comes from.} \emph{Token-in-token-out}
(TITO) is the contract that the rollout path carries \textbf{token IDs},
not just text: the inference engine returns the \textbf{exact token IDs
it sampled} (plus their log-probs), and the trainer consumes
\textbf{those same token IDs} for the policy-gradient update --- never a
sequence re-derived by tokenizing decoded text. The need for it is
intrinsic to \emph{harnessed} agentic RL {[}1{]}. In the traditional
setup the training engine owns the environment loop and keeps one
continuously extended token history, so a rollout is a single linear
token sequence by construction. In the harnessed setup the
\textbf{harness} owns context construction, control flow, and tool
execution; it builds each prompt independently and talks to the model
over an ordinary, \textbf{text} LLM API. The trainer never touches the
environment --- it only observes a sequence of request/response pairs
\((p_1,a_1),(p_2,a_2),\dots\) crossing the LLM endpoint. TITO is the
guarantee that what crosses that seam is captured \emph{token-exact}:
the \(a_i\) the trainer learns from are the very tokens the policy
sampled.

\textbf{The problem --- token-prefix continuity breaks (retokenization
drift).} Harnesses chain calls in \textbf{text}: the previous call is a
clean text-level prefix of the next prompt. RL, however, needs
continuity at the \textbf{token} level, and the two do not coincide ---
a tokenizer is not a clean round-trip, so re-encoding the decoded
history can yield a \emph{different} token sequence than the one
sampled. The paper isolates three mechanisms:

\begin{itemize}
\tightlist
\item
  \textbf{Chat-template non-compositionality.} Rendering a whole message
  history is not the same as concatenating the renderings of its parts
  (\(\mathrm{Template}(A \Vert B) \neq \mathrm{Template}(A)\Vert\mathrm{Template}(B)\)):
  templates insert or drop delimiters at message boundaries. E.g. Qwen's
  chat template can strip an earlier
  \texttt{\textless{}think\textgreater{}} marker.
\item
  \textbf{Decode--retokenize drift.} Decoding is not injective, so
  re-tokenizing the decoded text need not recover the original IDs: the
  word \emph{having}, sampled as \texttt{h}+\texttt{aving}, can
  re-tokenize to \texttt{hav}+\texttt{ing} inside a later prompt ---
  identical text, different token boundaries.
\item
  \textbf{Inference-time output transformation.} Tool-call and
  structured-output handlers parse, normalize, repair, and reserialize a
  response before returning it, changing whitespace, delimiters, or JSON
  --- so even the \emph{text} folded into the next prompt can differ
  from what the sampled IDs represent.
\end{itemize}

If training reconstructs tokens from this text rather than using the
recorded IDs, three things silently misalign:

\begin{itemize}
\tightlist
\item
  \textbf{The action.} The token the loss is computed on is no longer
  the token the policy emitted.
\item
  \textbf{The log-prob / importance ratio.} On-policy RL weights each
  token by the ratio of new to old policy probability; if the ``old''
  tokens are reconstructed rather than recorded, the ratio is computed
  against tokens that were never sampled, biasing the gradient.
\item
  \textbf{The loss mask.} The mask that marks \emph{policy} tokens
  (train) vs.~\emph{prompt/observation} tokens (ignore) is defined per
  token; a shifted boundary leaks environment tokens into the loss, or
  drops the model's own.
\end{itemize}

The same drift also breaks \emph{sample merging}: two consecutive calls
can be stitched into one training sequence only when their token IDs
actually satisfy prefix continuity, so drift turns one rollout into a
\emph{dynamic} number of samples. For RLVR --- where the reward is a
precise 1.0/0.0 from a checker --- this token-level slop quietly
corrupts the very signal the method depends on.

\textbf{The solution space.} The paper frames four strategies, trading
correctness against prefix-computation reuse:

\begin{itemize}
\tightlist
\item
  \textbf{Independent per-call training.} Treat each \((p_i,a_i)\) as
  its own sample; compute loss without assuming any relation to
  neighbors. Token-exact by construction, but recomputes the long shared
  prefixes across calls --- wasteful.
\item
  \textbf{Buffered token replacement} (AReaL, verl Uni-Agent). The proxy
  buffers each call's text and tokens; when a new request's text
  matches, it splices the \emph{recorded} response IDs back in so
  token-prefix continuity always holds. This recovers reuse, but when
  the buffered IDs differ from those in the actual next request it
  evaluates the response under a \emph{stitched} prompt the policy never
  saw --- an off-policy discrepancy.
\item
  \textbf{Prefix-shared / tree-structured training.} Store common
  prefixes once and use a branch-aware attention mask so each token
  attends only to its ancestors. Recovers the reuse of independent
  training, but demands heavy backend support (tree packing, custom
  masks/kernels, distributed-gradient handling).
\item
  \textbf{Best-effort sequence merging} (Agent Lightning v1.0's choice).
  Merge two calls only when their observed token IDs satisfy prefix
  continuity; otherwise close the sequence and start a new one. This
  preserves the prompt actually consumed during rollout, works with
  standard dense causal kernels, and lets drift merely lower the merge
  ratio rather than corrupt the gradient.
\end{itemize}

\textbf{Why the proxy owns it.} Every one of these strategies rests on
the same primitive: \emph{capturing the sampled token IDs at the moment
of generation} instead of guessing them from text later. The proxy is
the only component that straddles both sides of the seam --- it speaks
the \textbf{text} LLM API to the harness and the \textbf{token} API to
the inference engine (SGLang returns \texttt{response\_ids} + log-probs)
--- so it is the natural, and only correct, place to do it. Concretely
the proxy records \texttt{\{prompt\_ids,\ response\_ids,\ logprobs\}}
per turn (Figure above), keyed by \texttt{session\_id} (Agent
Lightning's per-call \texttt{model\_request} event); the \textbf{Sample
Adapter} then assembles the training \texttt{Sample} directly from those
recorded IDs and the loss mask and attaches the reward --- so the
trainer learns from the exact tokens sampled and the whole drift class
is sidestepped before any merging decision. The harness stays blissfully
text-only, yet the trainer still gets token-exact, on-policy data. In
short: \textbf{the proxy is where an off-the-shelf, text-based agent
loop is reconciled with the token-exact requirement of RL --- that
reconciliation is TITO.}

\hypertarget{when-one-episode-becomes-many-samples-subagents-summarization-and-advantage-assignment}{%
\subsection{When one episode becomes many samples --- subagents,
summarization, and advantage
assignment}\label{when-one-episode-becomes-many-samples-subagents-summarization-and-advantage-assignment}}

\textbf{This is a different problem from TITO --- its sibling.} TITO
asks \emph{``are the tokens I train on exactly the ones the policy
sampled?''} --- a question about a single call chain crossing the
text↔token seam. The problem in this subsection asks \emph{``one episode
produced how many training samples, and how do I attach reward,
advantage, and loss across them?''} --- a question about the
\textbf{shape of a whole rollout}. Both descend from the same root ---
in \emph{harnessed} agentic RL the harness (not the trainer) owns
context construction and control flow, so a rollout is a set of
separately-constructed call chains rather than one linear history
{[}1{]} --- but they are handled in different places. Get TITO right and
each individual sample is token-exact; get \emph{this} wrong and you
train on token-exact samples with \textbf{biased rewards and
mis-normalized gradients}.

\textbf{The problem.} In traditional agentic RL one rollout is one
Markov trajectory and maps to exactly one training sample. In harnessed
agentic RL a single episode can fan out into a \textbf{dynamic number of
training samples} \(N_\rho\), known only \emph{after} the harness has
run and samples have been constructed. This is not a corner case: in
Agent Lightning's coding-agent runs only \textbf{36\% of rollouts on
average stay a single training sample, and the mean rollout yields 2.4
samples} {[}1{]}. Three harness behaviours drive the fan-out --- and all
three are exactly the ones you called out:

\begin{itemize}
\tightlist
\item
  \textbf{Context summarization / compaction.} To stay under the context
  window the harness periodically replaces the running message history
  with a summary. That \emph{rewrites the token prefix} the next call is
  built on, so the new call no longer has the previous call as a
  token-level prefix --- the chain must be cut and a new sample started.
\item
  \textbf{Response rewriting / expansion before acting.} The harness may
  parse, repair, reserialize, or expand a model response (tool-call
  normalization, JSON repair, template surgery) before it acts on it or
  folds it into the next prompt --- so what re-enters the history
  differs from the sampled tokens, again breaking the merge condition.
\item
  \textbf{Subagents.} The harness spawns a sub-agent that runs its
  \textbf{own} chain of model calls with its \textbf{own} prompt
  history. These branches do not share one linear history with the
  parent at all; each is its own sample (or its own little tree).
\end{itemize}

\textbf{The cause --- the harness is the only component that sees a
whole episode.} The proxy sees a stream of independent per-call
\texttt{(prompt\_ids,\ response\_ids)} pairs; the sandbox sees a stream
of independent tool ops and one terminal reward. \textbf{Neither sees
the episode as a whole} --- only the harness, sitting between them,
knows that these calls, that summary, and those subagent branches all
belong to \emph{one} rollout \(\rho\) with \emph{one} outcome. So the
harness is the only place that can stamp every fragment with a shared
\texttt{rollout\_id} (and \texttt{prompt\_group\_id}), and the training
side must honour that grouping --- otherwise it silently treats an
incidental fan-out as if it were more independent data. The distinction
from TITO matters here: retokenization drift (§``Token-in-token-out'')
is \emph{incidental} and merely lowers the sample-merge ratio, whereas
summarization and subagent spawning are \emph{structural} --- they are
deliberate harness operations that split a rollout by design.

\begin{figure}[htbp]
\centering
\resizebox{\textwidth}{!}{%
\begin{tikzpicture}[
  font=\small,
  box/.style={draw, rounded corners, minimum height=1.05cm, minimum width=2.2cm, align=center, thick, fill=white},
  cpu/.style={box, fill=black!4},
  focus/.style={box, fill=orange!18},
  call/.style={draw, rounded corners, thick, fill=blue!8, minimum height=0.72cm, minimum width=1.15cm, align=center, font=\scriptsize},
  smp/.style={draw, rounded corners, thick, fill=green!10, minimum height=0.9cm, minimum width=2.1cm, align=center, font=\scriptsize},
  lbl/.style={font=\scriptsize, align=center, fill=white, inner sep=1.5pt},
  op/.style={font=\scriptsize\itshape, text=black!60, align=center},
  panel/.style={draw, rounded corners, dashed, thick, inner sep=0.32cm},
  >={Stealth[]},
]
% left proxy, right sandbox
\node[cpu]   (proxy)   at (-3.6,0.1) {\textbf{LLM Proxy}\\[1pt]\scriptsize records\\per-call tokens};
\node[focus] (sandbox) at (15.6,0.1) {\textbf{Sandbox}\\[1pt]\scriptsize tools + checker\\$\Rightarrow$ reward $R$};

% harness episode = tree of call chains
\node[call] (c1) at (2.2,1.8) {call 1};
\node[call] (c2) at (4.5,1.8) {call 2};
\node[call] (c3) at (8.2,1.8) {call 4};
\node[call] (s1) at (5.6,-1.5) {sub 1};
\node[call] (s2) at (7.8,-1.5) {sub 2};
\draw[->,thick] (c1)--(c2);
\draw[->,thick,dashed] (c2) -- node[op,above]{summarize:\\prefix reset} (c3);
\draw[->,thick,dashed] (c2.south) -- node[op,fill=white,inner sep=2pt,pos=0.5]{spawn\\subagent} (s1.north);
\draw[->,thick] (s1)--(s2);
\begin{scope}[on background layer]
  \node[panel, fill=orange!4, fit=(c1)(c2)(c3)(s1)(s2)] (harness) {};
\end{scope}
\node[op, anchor=south west, font=\scriptsize\bfseries] at ([yshift=2pt]harness.north west) {Harness --- ONE episode (rollout $\rho$): the only whole-episode view};

\draw[<->,thick] (proxy) -- node[lbl,pos=0.6]{turn + tokens} (harness.west);
\draw[<->,thick] (harness.east) -- node[lbl]{tool / reward} (sandbox);

% bottom: the fan-out into samples (ordered left->right = S1,S2,S3, each under its source)
\node[smp] (S1) at (2.6,-4.2) {$S_1$: calls 1--2\\reward $R$,\ id $\rho$};
\node[smp] (S2) at (6.7,-4.2) {$S_2$: sub 1--2\\reward $R$,\ id $\rho$};
\node[smp] (S3) at (10.6,-4.2) {$S_3$: call 4\\reward $R$,\ id $\rho$};
\node[cpu] (adv) at (15.6,-4.2) {\textbf{Advantage + loss}\\[1pt]\scriptsize grouped over\\\emph{rollouts}, not samples};
\draw[->,thick] (c2) -- (S1);
\draw[->,thick] (s2) -- (S2);
\draw[->,thick] (c3) -- (S3);
\draw[->,thick] (S3) -- node[lbl,above]{group by id $\rho$} (adv);
\end{tikzpicture}%
}
\caption{How it works with the other components. The harness sits between the proxy and the sandbox and is the \textbf{only} component that sees a whole episode. One rollout $\rho$ fans out into several training samples: a spawned \emph{subagent} runs its own call chain ($S_2$), and a \emph{summarize} step resets the token prefix so calls 1--2 close as $S_1$ while the post-summary call 4 opens $S_3$. Every sample inherits the episode's terminal reward $R$ and carries the same \texttt{rollout\_id} $\rho$; advantage baselines and loss normalization must then be computed \textbf{per rollout}, not per sample, so an incidental fan-out cannot reweight the update.}
\end{figure}

\textbf{Why the naive handling corrupts training.} Once one rollout
\(\rho\) owns \(N_\rho\) samples that all inherit its outcome reward
\(R_\rho\), two group-level computations become ambiguous:

\begin{itemize}
\tightlist
\item
  \textbf{Advantage baseline (GRPO).} GRPO's advantage subtracts a group
  baseline \(\bar r\) over the samples generated from one prompt. If the
  baseline is taken at the \textbf{sample} level, a rollout that
  happened to split into more samples silently gets more weight in the
  baseline. Concretely: in one group Rollout 1 succeeds (\(R=1\)) and
  Rollout 2 fails (\(R=0\)). If Rollout 1 fanned out into 3 samples and
  Rollout 2 stayed 1, the \emph{rollout-level} baseline is the intended
  \(\bar r=(1+0)/2=\tfrac12\), but the \emph{sample-level} baseline is
  \(\bar r=(1+1+1+0)/4=\tfrac34\) --- a biased baseline that no one
  chose, driven purely by how the harness happened to fragment.
  (Frameworks disagree today: verl Uni-Agent and Polar compute advantage
  at the rollout level; slime and AReaL at the sample level {[}1{]}.)
\item
  \textbf{Loss normalization.} With a dynamic per-rollout sample count,
  sample-level \emph{seq-mean-token-mean} normalization (GRPO's default)
  gives disproportionately more gradient weight to rollouts that
  produced more samples. Since the sample count is an incidental
  artifact of summarization/subagents, that reweighting is unprincipled.
\end{itemize}

\textbf{The solutions --- reconcile at the rollout level.} The fix is a
small, coherent set of choices that make the incidental fan-out
invisible to the gradient:

\begin{enumerate}
\def\labelenumi{\arabic{enumi}.}
\tightlist
\item
  \textbf{Sample construction: best-effort merging, honest splits.}
  Merge consecutive calls into one sample only when their token IDs
  satisfy prefix continuity (the TITO merge condition); when
  summarization, response rewriting, or a subagent breaks that prefix,
  \textbf{close the current sample and open a new one} rather than force
  a stitch. Represent the episode as the resulting \emph{set/tree} of
  samples, all tagged with the shared \texttt{rollout\_id}.
\item
  \textbf{Reward: outcome-inherited.} The reward is terminal and
  per-episode (the sandbox checker's \(1.0/0.0\)); \textbf{every sample
  in the rollout inherits \(R_\rho\).} A subagent's samples carry the
  parent episode's outcome.
\item
  \textbf{Advantage at the rollout level.} Compute the GRPO group
  baseline over \textbf{rollouts}, not samples, so --- as in the worked
  example --- a rollout that fragmented into more samples cannot move
  the baseline. Agent Lightning argues this is the principled default;
  retokenization and harness internals ``should not be allowed to change
  the baseline of the entire group'' {[}1{]}.
\item
  \textbf{Loss normalization at the rollout level.} Use a
  \textbf{rollout-level token-mean} loss: pool all response tokens of a
  rollout, then average uniformly over rollouts (slime's choice). This
  is insensitive to how many samples a rollout produced. Plain
  token-mean is theoretically clean too but empirically unstable when a
  batch has many long negative samples; sample-level seq-mean is the one
  to avoid here.
\item
  \textbf{Preserve provenance to the backend.} When samples are
  flattened into a physical tensor batch, every sequence must retain its
  \texttt{rollout\_id} and \texttt{prompt\_group\_id} so steps 3--4 can
  still be computed after packing onto fixed GPU workers.
\item
  \textbf{Open problem: intra-rollout credit assignment.}
  Outcome-inheritance is a placeholder --- it does not say \emph{which}
  branch, subagent, or turn earned the reward. Finer credit assignment
  across the samples within a rollout (turn-level grouping such as GiGPO
  {[}26{]}, process rewards) is active research; see the
  credit-assignment item in §``Open problems in production agentic RL''.
\end{enumerate}

\textbf{Where this project sits.} The training demos in this report
deliberately use \textbf{short, single-chain episodes}: a binary
end-of-episode reward with \texttt{EPISODE\_MAX\_TURNS=8}, no context
summarization and no subagents, so every rollout collapses to exactly
one sample and the fan-out above never bites --- which is why
\texttt{build\_token\_stream} can safely return a \emph{bare}
\texttt{Sample} (§``Key code --- three pieces''). That is a modelling
convenience, not a limitation of the stack. The moment the harness is a
real, full-featured agent --- one that compacts context and spawns
subagents --- this subsection's machinery becomes mandatory; the removed
\texttt{demo\_cpu\_branching} (§``History: removed demos (2026-09-30)'')
was a standalone demo that switched the mini\_claude harness for the
\textbf{Claude Code CLI}, deliberately provoked summarization and
subagent branching, and was built to show one rollout producing several
samples that are then reconciled with rollout-level advantage and
rollout-level token-mean loss.

\hypertarget{slime-the-trainer-and-its-asynchronous-rollout}{%
\subsection{slime --- the trainer, and its asynchronous
rollout}\label{slime-the-trainer-and-its-asynchronous-rollout}}

The trainer in this report is \textbf{slime} (THUDM) --- an RL
post-training framework built on \textbf{Megatron} (training) +
\textbf{SGLang} (rollout), designed so that agentic workloads plug in as
\emph{data-generation} functions without forking the training kernel
{[}9{]}. Its documented architecture maps cleanly onto Figure 1: a
\textbf{Data Buffer} bridges a \textbf{Megatron} training group and one
or more \textbf{SGLang} rollout servers behind a \textbf{router}; a
\emph{custom rollout generation} hook (where our harness + AgentCore
sandbox live) produces samples into the buffer, and Megatron
periodically pushes updated weights to the SGLang servers.

\begin{figure}[htbp]
\centering
\resizebox{0.95\textwidth}{!}{%
\begin{tikzpicture}[
  font=\small,
  box/.style={draw, rounded corners, minimum height=1.2cm, minimum width=3.0cm, align=center, thick, fill=white},
  gpu/.style={box, fill=blue!8},
  cpu/.style={box, fill=black!4},
  focus/.style={box, fill=orange!18},
  lbl/.style={font=\scriptsize, align=center, fill=white, inner sep=1.5pt},
  ptitle/.style={font=\scriptsize\bfseries, text=black!55},
  panel/.style={draw, rounded corners, dashed, thick, inner sep=0.34cm},
  >={Stealth[]},
]
% central bridge + custom hook
\node[cpu]   (buffer)  at (0,0)      {\textbf{Data Buffer}\\[1pt]\scriptsize prompts + rollout data};
\node[focus] (rollout) at (6.8,0)    {\textbf{Custom rollout gen.}\\[1pt]\scriptsize our harness + AgentCore};
% rollout (SGLang) group
\node[cpu]   (router)  at (6.8,-2.6) {\textbf{SGLang router}};
\node[gpu]   (sg1)     at (5.0,-5.0) {\textbf{SGLang server}\\[1pt]\scriptsize GPU};
\node[gpu]   (sg2)     at (8.6,-5.0) {\textbf{SGLang server}\\[1pt]\scriptsize GPU};
% training group
\node[gpu]   (mega)    at (0,-5.0)   {\textbf{Megatron}\\[1pt]\scriptsize training group (GPU)};

% buffer <-> custom rollout generation
\draw[->, thick] ([yshift=4pt]buffer.east) -- node[lbl,above]{send init prompt} ([yshift=4pt]rollout.west);
\draw[->, thick] ([yshift=-4pt]rollout.west) -- node[lbl,below]{recv rollout data} ([yshift=-4pt]buffer.east);
% custom rollout -> router -> servers  (generation)
\draw[->, thick] (rollout) -- node[lbl,right]{LLM calls} (router);
\draw[->, thick] (router)  -- (sg1);
\draw[->, thick] (router)  -- (sg2);
% buffer -> megatron (train data)
\draw[->, thick] (buffer)  -- node[lbl,left]{send train data} (mega);
% megatron -> servers (weight update)
\draw[->, thick] (mega)    -- node[lbl,above]{update weights} (sg1);

% domain panels (behind boxes)
\begin{scope}[on background layer]
  \node[panel, fill=blue!3,   fit=(mega)] (ptrain) {};
  \node[panel, fill=orange!4, fit=(router)(sg1)(sg2)] (proll) {};
\end{scope}
\node[ptitle, anchor=south west] at ([yshift=2pt]ptrain.north west) {Training (Megatron)};
\node[ptitle, anchor=south west] at ([yshift=2pt]proll.north west) {Rollout (SGLang)};
\end{tikzpicture}%
}
\caption{slime's architecture (redrawn after the slime documentation [9]). A \textbf{Data Buffer} bridges the two sides: it hands prompts to a \emph{custom rollout generation} hook --- where our \textbf{harness + AgentCore sandbox} live --- and receives finished rollout data back. Generation runs through an \textbf{SGLang router} to one or more GPU \textbf{SGLang servers}. On the training side the buffer feeds \textbf{Megatron}, which periodically pushes \textbf{updated weights} to the SGLang servers. Agentic workloads plug into the hook without forking the training kernel.}
\end{figure}

\textbf{Asynchronous mode.} By default training and rollout alternate,
so the GPUs on each side idle while the other runs. slime avoids this
with an \textbf{async training driver} (\texttt{train\_async.py} instead
of \texttt{train.py}): while Megatron trains on the current batch, the
rollout side already generates the \emph{next} rollout's data,
overlapping the two {[}9{]}. For agentic RL --- where rollouts have a
\textbf{long tail} (one slow episode, e.g.~a lengthy compile, can stall
a whole step) --- slime also offers a \textbf{fully-async rollout} path
(\texttt{-\/-rollout-function-path\ slime.rollout.fully\_async\_rollout.generate\_rollout\_fully\_async})
that keeps a fixed pool of in-flight generations across rollout
boundaries, so a training step never waits on the slowest sample. This
is exactly the regime this report targets: \textbf{AgentCore sandboxes
have variable, sometimes long, per-episode latency}, and async rollout
keeps the trainer fed while they run.

\emph{Sources: slime repository \url{https://github.com/THUDM/slime} and
documentation \url{https://thudm.github.io/slime/} {[}9{]}.}

\hypertarget{scaling-the-loop-inference-harness-and-sandbox-in-balance}{%
\subsection{Scaling the loop --- inference, harness, and sandbox in
balance}\label{scaling-the-loop-inference-harness-and-sandbox-in-balance}}

The rollout loop has \textbf{three} components, and they scale as one
system because they are interconnected. The \textbf{inference engine}
(SGLang) produces every model turn; the \textbf{sandbox} (AgentCore)
runs every tool call and the checker; and the \textbf{harness} --- the
agent-loop driver that opens each episode, calls the policy turn by
turn, runs the checker, and packs the result into a \texttt{Sample} ---
sits between them and paces both. Rollout throughput is bounded by the
slowest of the three: \textbf{buffer-fill rate =
\texttt{min(inference,\ sandbox,\ harness)}.} Scale two and leave the
third and the bottleneck merely moves onto the third. So the job has
three parts, in order: \textbf{diagnose} which tier is slow,
\textbf{scale} it, and \textbf{balance load correctly on each tier} ---
the three balance by \emph{different} rules, and applying one tier's
rule to another is a mistake. This subsection treats all three together;
the diagram below is the map.

\textbf{Why scale the sandbox at all?} A policy only \emph{generalizes}
if it is trained across many \textbf{diverse tasks} --- SWE repos, SQL,
math, tool-use, RAG, CTF, GPU kernels --- and across \textbf{many
episodes of each} so the reward signal is not dominated by one family.
That means running a large number of \textbf{concurrent episodes} per
training step, and every concurrent episode needs its own isolated
sandbox. So scaling \emph{up the number of concurrent sandboxes} (the
admission knob \texttt{CONC}: \texttt{S\ =\ CONC\ ×\ engines} episodes
in flight) is the direct lever on task diversity and throughput.

\textbf{But sandboxes do not run in isolation.} Every one of those
concurrent episodes still gets its model turns from the \textbf{same}
inference engine. If we raise sandbox concurrency while holding the
number (and size) of inference replicas \textbf{fixed}, generation
requests start to \textbf{queue} at the inference engine: per-turn
latency climbs, and the extra sandboxes sit idle waiting for their next
turn. Past that point the inference engine --- not the sandbox --- is
the bottleneck, and adding sandboxes buys nothing.

\textbf{So first diagnose which side is slow.} Per episode the
wall-clock is split between (a) \textbf{generation} --- the inference
engine producing the turn --- and (b) \textbf{sandbox execution} ---
running the tool call and the checker. The harness measures both:
\texttt{{[}sglang-gen{]}} → \texttt{perf/sglang\_gen\_ms\_\{p50,p95\}}
for generation, and \texttt{{[}sbx-op{]}} →
\texttt{perf/sandbox\_ms\_\{p50,p95\}} for the sandbox round-trip.
Compare them before scaling:

\begin{itemize}
\tightlist
\item
  \textbf{Generation dominates} (\texttt{sglang\_gen\_ms} ≫
  \texttt{sandbox\_ms}) ⇒ \textbf{inference-bound}. Scale the inference
  engine (more replicas / more GPUs / larger tensor-parallel),
  \emph{not} the sandbox count.
\item
  \textbf{Sandbox dominates} (\texttt{sandbox\_ms} ≫
  \texttt{sglang\_gen\_ms}) ⇒ \textbf{sandbox-bound}. Scale sandbox
  concurrency / instances; one inference replica can feed many episodes.
\end{itemize}

\textbf{Examples.}

\begin{itemize}
\tightlist
\item
  \textbf{Sandbox-bound --- KernelBench (GPU).} The model emits a short
  CUDA kernel in a few hundred ms, but \texttt{nvcc} compilation plus
  \texttt{allclose} on the GPU takes \textbf{seconds}. Sandbox time
  dwarfs generation → add GPU sandbox Instances; a single inference
  replica keeps up.
\item
  \textbf{Inference-bound --- math / long-reasoning.} The checker is a
  one-line stdlib assert (sandbox ≈ tens of ms), but the model generates
  a long chain of thought (generation ≈ seconds). Inference time dwarfs
  the sandbox → adding sandboxes does nothing; add inference replicas.
\item
  \textbf{Balanced --- SWE with a moderate test suite.} Generation and a
  multi-second test run are comparable → scale \textbf{both} roughly in
  proportion to keep either side from queuing.
\end{itemize}

The practical loop: pick a concurrency, read
\texttt{perf/sglang\_gen\_ms\_*} and \texttt{perf/sandbox\_ms\_*}, and
grow whichever side dominates until the two are balanced --- then raise
both together.

\textbf{Balancing load as you scale --- least-loaded engine per episode,
round-robin the sandbox pool.} Once you add replicas on \emph{either}
side, the next question is how episodes get spread across them. The two
sides answer it differently, and it matters that you \emph{not} apply
the same rule to both:

\begin{quote}
\textbf{Do you need a load balancer?} For the \textbf{sandbox: no}
within one Runtime --- the session id \emph{is} the balancer (the
capacity provider places each session for you); across a \emph{pool} of
Runtime ARNs you add only trivial \textbf{sticky round-robin by session
id}, not a real LB. For \textbf{inference: also no} --- the harness tier
(\texttt{harness\_pool.py}) reads the engine list from slime's SGLang
router once and sends each \emph{episode} to the engine with the fewest
in-flight episodes, sticky for all its turns. In short: you never stand
up a load balancer of your own; you mint distinct session ids, count
in-flight episodes per engine, and (only across a Runtime pool) route
sessions stickily.
\end{quote}

\begin{itemize}
\tightlist
\item
  \textbf{Inference --- balance per episode, then stay sticky.} slime
  runs \textbf{one SGLang router} in front of its engines, which
  balances each \emph{request} cache-aware (prefix affinity,
  shortest-queue fallback). A multi-turn episode's prompt grows turn by
  turn on its own prefix, so the harness tier balances one level up: it
  sends each \textbf{episode} to the engine with the fewest in-flight
  episodes (ties rotate, so an idle start does not pile onto engine 0)
  and keeps \textbf{all of that episode's turns} there, where its prefix
  is already in the KV cache. This replaced an earlier crc32 hash of the
  episode onto an engine (\texttt{router\_trace.pin\_engine}), which was
  sticky but blind to load. What to avoid is per-request round-robin: it
  would scatter an episode's prefix across every replica and destroy the
  prefix cache. Scaling inference stays ``add engines''; the per-episode
  \texttt{engine\_inflight} field shows when inference binds.
\item
  \textbf{Sandbox --- the session id \emph{is} the load-balancing key.}
  One \texttt{runtimeSessionId} maps to exactly one isolated microVM /
  instance. Within a \textbf{single} Runtime, AgentCore's capacity
  provider already \textbf{places each new session on an available
  instance}, so simply minting a distinct per-episode session id spreads
  the load for you. You only need \emph{explicit} round-robin when you
  run a \textbf{pool of sandbox Runtimes} (several ARNs). Three things
  trigger a pool --- otherwise a single ARN is enough: \textbf{(1)
  concurrency} --- your in-flight episodes
  (\texttt{S\ =\ CONC\ ×\ num\_engines}) would exceed one Runtime's
  per-region session cap, so you split the load across \texttt{K} ARNs;
  \textbf{(2) AZ/account spread} --- placing Runtimes in different
  availability zones (resilience) or accounts (isolation, separate
  quotas); \textbf{(3) capability routing} --- different families need
  \emph{different sandbox images} (a GPU family, a DB-baked image, a
  per-repo SWE image, §3), so you register one ARN per capability and
  pick by \texttt{family} rather than pure round-robin. Below these
  thresholds --- episode count fits one Runtime's quota and one generic
  image serves every family --- you stay on a \textbf{single ARN} and
  let the capacity provider place sessions for you. When you do pool,
  pin each episode to one ARN and make it \textbf{sticky} so every op of
  that episode returns to the same warm sandbox that holds its
  workspace:
\end{itemize}

\begin{Shaded}
\begin{Highlighting}[]
\ImportTok{import}\NormalTok{ hashlib, os}

\CommentTok{\# 1) Inference: the least{-}loaded engine per EPISODE, sticky for all its turns}
\CommentTok{\#    (harness\_pool.py). The engine URLs come from slime\textquotesingle{}s router; ties rotate.}
\KeywordTok{def}\NormalTok{ acquire\_engine(}\VariableTok{self}\NormalTok{):}
\NormalTok{    n }\OperatorTok{=} \BuiltInTok{len}\NormalTok{(}\VariableTok{self}\NormalTok{.\_inflight)}
    \VariableTok{self}\NormalTok{.\_rr }\OperatorTok{=}\NormalTok{ (}\VariableTok{self}\NormalTok{.\_rr }\OperatorTok{+} \DecValTok{1}\NormalTok{) }\OperatorTok{\%}\NormalTok{ n}
\NormalTok{    i }\OperatorTok{=} \BuiltInTok{min}\NormalTok{(}\BuiltInTok{range}\NormalTok{(n), key}\OperatorTok{=}\KeywordTok{lambda}\NormalTok{ j: (}\VariableTok{self}\NormalTok{.\_inflight[j], (j }\OperatorTok{{-}} \VariableTok{self}\NormalTok{.\_rr) }\OperatorTok{\%}\NormalTok{ n))}
    \VariableTok{self}\NormalTok{.\_inflight[i] }\OperatorTok{+=} \DecValTok{1}                     \CommentTok{\# released when the episode ends}
    \ControlFlowTok{return}\NormalTok{ i, }\VariableTok{self}\NormalTok{.\_engine\_urls[i]}

\CommentTok{\# 2) Sandbox: within ONE Runtime the session id already spreads sessions across}
\CommentTok{\#    instances. Across a POOL of sandbox Runtimes, round{-}robin BY session id so}
\CommentTok{\#    the spread is even AND sticky (all ops of an episode hit one warm microVM).}
\NormalTok{SANDBOX\_ARNS }\OperatorTok{=}\NormalTok{ os.environ[}\StringTok{"AGENTCORE\_SANDBOX\_RUNTIME\_ARNS"}\NormalTok{].split(}\StringTok{","}\NormalTok{)}

\KeywordTok{def}\NormalTok{ pick\_sandbox\_arn(session\_id: }\BuiltInTok{str}\NormalTok{) }\OperatorTok{{-}\textgreater{}} \BuiltInTok{str}\NormalTok{:}
\NormalTok{    h }\OperatorTok{=} \BuiltInTok{int}\NormalTok{(hashlib.sha1(session\_id.encode()).hexdigest(), }\DecValTok{16}\NormalTok{)}
    \ControlFlowTok{return}\NormalTok{ SANDBOX\_ARNS[h }\OperatorTok{\%} \BuiltInTok{len}\NormalTok{(SANDBOX\_ARNS)]   }\CommentTok{\# deterministic, sticky per episode}
\end{Highlighting}
\end{Shaded}

Hashing the session id (rather than a shared counter) gives an even,
\textbf{sticky} spread with no cross-thread lock --- which matters
because episodes run concurrently on \texttt{asyncio.to\_thread}. A
plain counter (\texttt{itertools.cycle} behind a lock) is the simpler
\emph{literal} round-robin, but only use it if you do not need
stickiness --- and here you do: every op of an episode (seed → exec →
read) must land on the microVM that holds that episode's files.

\emph{(Implementation note: the shipped \texttt{make\_sandbox()}
(\texttt{sandbox.py}) actually uses the
\textbf{\texttt{itertools.cycle}-behind-a-lock} counter, not the
session-id hash sketched above. Stickiness still holds --- the ARN is
chosen \textbf{once} when the episode's sandbox object is created and
fixed for all its ops --- so each episode keeps one warm microVM; only
the }selection* is a rotating counter rather than a hash. The hash form
above is the lock-free alternative you would switch to if the shared
counter's lock became a contention point at high \texttt{S}.)*

\begin{figure}[htbp]
\centering
\resizebox{0.95\textwidth}{!}{%
\begin{tikzpicture}[
  font=\small,
  box/.style={draw, rounded corners, minimum height=1.0cm, minimum width=2.5cm, align=center, thick, fill=white},
  gpu/.style={box, fill=blue!8},
  cpu/.style={box, fill=black!4},
  focus/.style={box, fill=orange!18},
  wrk/.style={draw, rounded corners, thick, fill=green!10, minimum height=0.72cm, minimum width=2.3cm, align=center, font=\scriptsize},
  lbl/.style={font=\scriptsize, align=center, fill=white, inner sep=1.5pt},
  ptitle/.style={font=\scriptsize\bfseries, text=black!55, align=center},
  panel/.style={draw, rounded corners, dashed, thick, inner sep=0.34cm},
  >={Stealth[]},
]
% harness tier (left): W dispatch workers (the loop itself runs in the sandbox session)
\node[wrk] (w1) at (0,0.95) {worker 1};
\node[wrk] (w2) at (0,0.05) {worker 2};
\node[font=\scriptsize] (wd) at (0,-0.7) {$\vdots$};
\node[wrk] (wN) at (0,-1.4) {worker $W$};
\begin{scope}[on background layer]
  \node[panel, fill=green!4, fit=(w1)(w2)(wN)] (harness) {};
\end{scope}
\node[ptitle, anchor=south] at ([yshift=2pt]harness.north) {Harness --- $W$ dispatch workers\\(pull queue, no LB)};

% inference branch (top)
\node[cpu] (router) at (5.4,2.7) {\textbf{Engine picker}\\[1pt]\scriptsize least-loaded, per episode};
\node[gpu, minimum width=2.1cm] (sgA) at (10.0,3.8) {\scriptsize replica A};
\node[gpu, minimum width=2.1cm] (sgB) at (10.0,2.7) {\scriptsize replica B};
\node[gpu, minimum width=2.1cm] (sgC) at (10.0,1.6) {\scriptsize replica C};

% sandbox branch (bottom)
\node[cpu] (disp) at (5.4,-2.7) {\textbf{Sandbox dispatcher}\\[1pt]\scriptsize round-robin, sticky};
\node[focus, minimum width=2.1cm] (sbA) at (10.0,-1.6) {\scriptsize sandbox A};
\node[focus, minimum width=2.1cm] (sbB) at (10.0,-2.7) {\scriptsize sandbox B};
\node[focus, minimum width=2.1cm] (sbC) at (10.0,-3.8) {\scriptsize sandbox C};

% harness -> branches
\draw[->, thick] (harness.east) -- node[lbl,above,pos=0.6]{model turns} (router.west);
\draw[->, thick] (harness.east) -- node[lbl,below,pos=0.6]{tool ops} (disp.west);

% router -> replicas (cache-aware)
\draw[->, thick] (router) -- (sgA);
\draw[->, thick] (router) -- node[lbl,pos=0.45]{sticky per episode} (sgB);
\draw[->, thick] (router) -- (sgC);

% dispatcher -> sandboxes (round-robin)
\draw[->, thick] (disp) -- (sbA);
\draw[->, thick] (disp) -- node[lbl,pos=0.45]{by session id} (sbB);
\draw[->, thick] (disp) -- (sbC);

% panels
\begin{scope}[on background layer]
  \node[panel, fill=blue!3,   fit=(router)(sgA)(sgC)] (pinf) {};
  \node[panel, fill=orange!4, fit=(disp)(sbA)(sbC)]   (psbx) {};
\end{scope}
\node[ptitle, anchor=south west] at ([yshift=2pt]pinf.north west) {Inference --- fewest in-flight episodes, all turns on one engine};
\node[ptitle, anchor=north west] at ([yshift=-2pt]psbx.south west) {Sandbox --- round-robin by session id, sticky per episode};
\end{tikzpicture}%
}
\caption{The three interconnected tiers scale together, and \textbf{each balances load by a different rule} --- buffer-fill rate is \texttt{min(inference, harness, sandbox)}. \textbf{Harness (left):} $W$ = \texttt{HARNESS\_WORKERS} dispatch workers in the trainer (default one rollout step, \texttt{RB} $\times$ \texttt{NS}); they need \emph{no} load balancer --- a free worker \emph{pulls} the next episode off one shared queue (work-stealing). The agent loop itself runs \emph{inside} the episode's sandbox session, so the harness scales with the sandbox, not the trainer. \textbf{Inference (top):} each episode takes the SGLang engine with the fewest in-flight episodes (ties rotate) and keeps all its turns there, so its growing prompt prefix stays in one KV cache; scaling is ``add engines.'' \textbf{Sandbox (bottom):} one \texttt{runtimeSessionId} $=$ one isolated microVM/instance, taken \emph{before} the engine. Within a single Runtime the capacity provider spreads sessions for you; across a \emph{pool} of sandbox Runtimes you \textbf{round-robin by session id}, kept \textbf{sticky} so every op of an episode (seed $\rightarrow$ exec $\rightarrow$ read) returns to the same warm microVM that holds its files.}
\end{figure}

\textbf{The third tier --- the harness, the driver behind both dials.}
The two dials above, inference and sandbox, are both \textbf{driven} by
a third component that is easy to forget because it runs neither on a
GPU nor behind a sandbox API: the \textbf{harness} --- the agent loop
that seeds each episode's sandbox, calls the policy turn by turn, runs
the tools and the checker, and packs the result into a \texttt{Sample}.
It has its own throughput ceiling, and if you scale inference and
sandbox to fill a deep async buffer while leaving the harness alone, the
bottleneck simply \textbf{moves onto the driver}: the buffer-fill rate
is \texttt{min(inference,\ sandbox,\ harness)}, and adding replicas to
two of the three buys nothing once the third saturates. So the harness
has to scale \textbf{along with} SGLang and the sandbox --- otherwise a
few coroutines silently become the throughput ceiling.

\textbf{Where the old ceiling was: one process.} slime runs \textbf{one}
\texttt{RolloutManager} Ray actor whose single event loop drives
\emph{every} \texttt{generate} coroutine
(\texttt{slime/ray/rollout.py}); the number of concurrent episodes is
bounded by one \texttt{asyncio.Semaphore} sized
\texttt{S\ =\ CONC\ ×\ num\_engines} --- with the demos' \texttt{CONC=4}
and two SGLang engines, \texttt{S\ =\ 8} in flight
(\texttt{slime/rollout/sglang\_rollout.py}). There is \textbf{no slime
knob for ``number of rollout worker processes''}: \texttt{CONC} raises
concurrency \emph{within} the one process, not the number of processes.
In the original design every episode's agent loop ran on that one event
loop, next to the trainer.

\textbf{Why one process ceilings.} Concurrency there is
\texttt{asyncio}, so it overlaps \emph{I/O waits} --- but not CPU work,
which the Python \textbf{GIL} serializes: every turn JSON-encodes an
invoke payload, decodes the response, and (at episode end) builds the
loss-masked token stream, and those run one-at-a-time. The synchronous
boto3 \texttt{invoke\_agent\_runtime} is offloaded with
\texttt{asyncio.to\_thread} onto the loop's default executor, which is ≈
\texttt{min(32,\ cpu+4)} threads unless configured, and every in-flight
episode also holds open file descriptors / sockets and buffered
trajectory state in RAM. So past some \texttt{S}, raising \texttt{CONC}
just deepens internal queues: per-episode latency climbs and the buffer
stops filling faster.

\textbf{What the harness tier changes --- the default in every demo.}
The fix has two halves (\texttt{harness\_pool.py}, the default rollout
topology of every demo since 2026-09-28, like \texttt{PREOPEN} and
\texttt{HEARTBEAT}). \textbf{(1) The loop moves into the sandbox
session.} The mini\_claude loop runs inside the episode's own session
and calls the engine's adapter back, so N sessions carry N loops and the
harness scales \textbf{with the sandbox, not as part of the trainer} ---
the placement DSec uses {[}31{]}. \textbf{(2) What stays in the trainer
is a dispatch tier sized on its own:} \texttt{W\ =\ HARNESS\_WORKERS}
workers (default \texttt{RB\ ×\ NS}, one rollout step) pulling episodes
off one shared queue, one \texttt{AnthropicAdapter} per engine, the
loop's default executor sized to the pool (\texttt{max(W,\ M)\ +\ 32}
threads instead of the stock 32), and one warm \texttt{aiohttp} session
with no connection limit. The three tiers are sized independently and
joined \textbf{per episode} --- no worker owns an engine or a session: a
free worker pulls the episode, takes its sandbox session first (fresh by
default; \texttt{SANDBOX\_SESSIONS=M} borrows one of M pooled sessions
per image, in a fresh workdir), then the least-loaded engine. Each
episode's \texttt{{[}demo{]}} record names the tier that binds:
\texttt{qdepth} \textgreater{} 0 = harness-bound,
\texttt{sandbox\_wait\_ms} \textgreater{} 0 = sandbox-bound, climbing
\texttt{engine\_inflight} = inference-bound.

\begin{longtable}[]{@{}
  >{\raggedright\arraybackslash}p{(\columnwidth - 8\tabcolsep) * \real{0.1346}}
  >{\raggedright\arraybackslash}p{(\columnwidth - 8\tabcolsep) * \real{0.2308}}
  >{\raggedright\arraybackslash}p{(\columnwidth - 8\tabcolsep) * \real{0.1731}}
  >{\raggedright\arraybackslash}p{(\columnwidth - 8\tabcolsep) * \real{0.2885}}
  >{\raggedright\arraybackslash}p{(\columnwidth - 8\tabcolsep) * \real{0.1731}}@{}}
\toprule\noalign{}
\begin{minipage}[b]{\linewidth}\raggedright
Component
\end{minipage} & \begin{minipage}[b]{\linewidth}\raggedright
Lifetime / isolation
\end{minipage} & \begin{minipage}[b]{\linewidth}\raggedright
Scaling unit
\end{minipage} & \begin{minipage}[b]{\linewidth}\raggedright
Ceiling of one unit
\end{minipage} & \begin{minipage}[b]{\linewidth}\raggedright
Where it runs
\end{minipage} \\
\midrule\noalign{}
\endhead
\bottomrule\noalign{}
\endlastfoot
Inference (SGLang) & persistent / shared & replica behind the router &
GPU KV cache + decode throughput & trainer GPUs \\
\textbf{Harness} & \textbf{the loop: per episode, in its sandbox
session; the dispatch tier: persistent} & \textbf{one loop per session;
\texttt{W} dispatch workers} & \textbf{trainer side: the adapters'
tokenisation and trajectory flattening (GIL)} & \textbf{the loop in the
sandbox session; dispatch + adapters on the trainer box} \\
Sandbox (AgentCore) & ad-hoc, per episode / strong (microVM) & session
id; pooled Runtime ARNs & per-Runtime session cap & AgentCore microVM /
Instance \\
\end{longtable}

\textbf{Where to put the harness --- three options.}

\begin{itemize}
\tightlist
\item
  \textbf{Co-located on the trainer box (the original design).} The one
  \texttt{RolloutManager} runs every loop next to Megatron and SGLang.
  Cheapest and lowest-latency to the buffer, and fine at
  \texttt{S\ ≈\ 8}; but the loops' threads, FDs, and RAM compete with
  the trainer for the box's CPU budget, and the one event loop and the
  GIL become the ceiling as you push \texttt{S} up.
\item
  \textbf{Scaled out to a CPU rollout fleet (disaggregated).} Because
  slime drives generation from a single actor, going past one process is
  an \emph{architecture} change, not a flag flip: run several rollout
  drivers (shard prompts across CPU nodes, each executing the same
  \texttt{generate} hook) that dial the \textbf{shared} SGLang router
  and the \textbf{pooled} sandbox ARNs. This decouples harness
  CPU/FDs/RAM from the GPU box and lets you scale drivers independently
  --- the disaggregated-placement lever of §``Open problems'' applied to
  the harness rather than to the trainer/inference split.
\item
  \textbf{Inside the sandbox session (the default in every demo since
  2026-09-28).} The loop runs where its tools run, one per session, and
  dials the trainer's per-engine adapter back over the VPC (the adapter
  ports are admitted SG→SG). The earlier alternative, a separate pool of
  harness Runtimes (the legacy two-runtime
  \texttt{agentcore\_runtime\_generate.py} /
  \texttt{harness\_runtime\_agent.py}, now deleted), added a second
  fleet to size and a network hop; running the loop in the episode's own
  session gets the loop out of the trainer without a second fleet. There
  is no in-trainer fallback: a session that cannot host the loop fails
  the episode with \texttt{HarnessPlacementError}.
\end{itemize}

\textbf{The joint recipe.} Diagnose the harness the way §``Scaling the
loop'' diagnoses the other two, from the per-episode \texttt{{[}demo{]}}
fields: \texttt{qdepth} \textgreater{} 0 means episodes wait for a free
dispatch worker (raise \texttt{HARNESS\_WORKERS}),
\texttt{sandbox\_wait\_ms} \textgreater{} 0 means they wait for a
session (raise sessions, the Runtime pool or
\texttt{SANDBOX\_SESSIONS}), and a climbing \texttt{engine\_inflight}
means inference binds (add engines). Then grow the three in balance. A
Little's-law sizing keeps them matched: the in-flight episodes needed to
saturate inference ≈
\texttt{inference\ throughput\ ×\ mean\ episode\ latency}; provision
\texttt{S\ =\ CONC\ ×\ engines} and \texttt{W} to hit that number (the
removed \texttt{SCALE=1} preset did exactly this at 4 : 256 : 256; see
§``History: removed demos (2026-09-30)''), then raise all three together
until GPU \textbf{goodput} --- not any single dial --- stops improving.

\hypertarget{tasks-frontier-labs-run-in-a-sandbox}{%
\section{Tasks frontier labs run in a
sandbox}\label{tasks-frontier-labs-run-in-a-sandbox}}

Frontier models are post-trained with \textbf{reinforcement learning
against executable, verifiable environments} (RLVR). The common
denominator is a \textbf{sandbox seam}: the model's untrusted output
runs in an isolated environment, and a \textbf{programmatic verifier}
--- not a human, and increasingly not an LLM judge --- turns the outcome
into a scalar reward.

\hypertarget{what-the-frontier-reports-describe}{%
\subsection{What the frontier reports
describe}\label{what-the-frontier-reports-describe}}

The trend is toward \textbf{verifiable rewards at scale}: Kimi K2, for
example, runs a rollout infrastructure of \textbf{more than 10,000
concurrent execution sandboxes} on Kubernetes for its coding and
software-engineering RL {[}5{]}. What the major open reports describe:

\begin{longtable}[]{@{}
  >{\raggedright\arraybackslash}p{(\columnwidth - 6\tabcolsep) * \real{0.2500}}
  >{\raggedright\arraybackslash}p{(\columnwidth - 6\tabcolsep) * \real{0.2500}}
  >{\raggedright\arraybackslash}p{(\columnwidth - 6\tabcolsep) * \real{0.2500}}
  >{\raggedright\arraybackslash}p{(\columnwidth - 6\tabcolsep) * \real{0.2500}}@{}}
\toprule\noalign{}
\begin{minipage}[b]{\linewidth}\raggedright
Model
\end{minipage} & \begin{minipage}[b]{\linewidth}\raggedright
RL / verifier setup
\end{minipage} & \begin{minipage}[b]{\linewidth}\raggedright
Task domains
\end{minipage} & \begin{minipage}[b]{\linewidth}\raggedright
Reward
\end{minipage} \\
\midrule\noalign{}
\endhead
\bottomrule\noalign{}
\endlastfoot
DeepSeek-R1 {[}2{]} & GRPO; rule-based verifier + code compiler (no
neural reward model, to avoid reward hacking) & math, competitive code &
answer match / tests pass (binary) \\
Qwen3 {[}3{]} & Reasoning RL from \textbf{3,995 query--verifier pairs};
General RL over \textbf{20+ tasks}; multi-turn tool use with
real-environment feedback & math, code, tool-use, RAG, instruction &
rule-based verifiable + reference-based model judge + reward model \\
Llama-Nemotron {[}4{]} & GRPO (science) + RLOO (instruction following);
vLLM + Megatron; prompts difficulty-filtered by pass-rate & science
reasoning, math, instruction & \textbf{LLM-judge} (Llama-3.3-70B)
answer-match + \texttt{\textless{}think\textgreater{}} format reward \\
AceReason-Nemotron {[}10{]} & math-then-code RL; a \textbf{local sandbox
verifier} runs code against \emph{all} test cases (64-worker pool) &
math (\textasciitilde49k problems), code (\textasciitilde8.5k problems)
& binary: pass \textbf{all} tests / \texttt{sympy} symbolic match on
\texttt{\textbackslash{}boxed\{\}} \\
Kimi K2 {[}5{]} & agentic RL; \textbf{verifiable-rewards ``gym'' + a
tool simulator (world model) + real execution sandboxes}
(\textgreater10k concurrent) & agentic tool use, coding, SWE, math,
instruction & executable unit tests + verifiable checks + self-critique
rubric \\
\end{longtable}

NVIDIA also \textbf{publishes the post-training data mix} openly: the
Nemotron post-training datasets {[}12{]} quantify the domain balance
used to train an open reasoning model --- roughly \textbf{Math 239k,
Code 175k, STEM 355k, Chat 628k} samples (plus large multilingual
splits) --- a concrete picture of how much of each family a frontier
post-training run consumes.

\hypertarget{task-families-in-detail}{%
\subsection{Task families in detail}\label{task-families-in-detail}}

Each family below lists its \textbf{dataset, reward function, key
metric, and environment requirement} --- the properties that determine
what a sandbox must provide --- plus an \textbf{episode shape}: the
typical \emph{turns per episode}, the \emph{response length per turn}
(the policy's own tokens), and the resulting \emph{total episode length}
(all turns' responses plus the observations fed back). The episode shape
is what you size the training config against: it sets
\textbf{\texttt{-\/-rollout-max-response-len}} (per-turn generation
budget), the \textbf{max turns} cap (\texttt{EPISODE\_MAX\_TURNS},
default \textbf{3} in this project's demos), the \textbf{sequence-length
/ context window} the trainer and inference engine must hold, and the
\textbf{rollout-buffer} memory per sample. The numbers below are typical
order-of-magnitude ranges (reasoning models sit at the high end of every
response figure), not hard limits --- measure your own before committing
a config.

\hypertarget{single-turn-generation-bound-families}{%
\subsubsection{Single-turn, generation-bound
families}\label{single-turn-generation-bound-families}}

The checker is milliseconds; nearly all the wall-clock is one long
generation. These are the \emph{inference-bound} families --- scale the
SGLang side.

\begin{itemize}
\tightlist
\item
  \textbf{Math \& logical reasoning} --- \emph{Dataset:} GSM8K, MATH-500
  {[}2{]}, AIME / HMMT; \textbf{DeepScaleR + NuminaMath} for RL
  (\textasciitilde49k problems) {[}10{]}. \emph{Reward:}
  \textbf{verifiable answer match} --- e.g.~a \texttt{sympy}-based check
  of symbolic equivalence on the \texttt{\textbackslash{}boxed\{\}}
  answer, binary, with no format or length bonus {[}10{]}.
  \emph{Metric:} accuracy, pass@k (Qwen3-235B's AIME'24 rose
  \textbf{70.1 → 85.1} over 170 RL steps {[}3{]}). \emph{Env:} a
  lightweight \textbf{CPU checker} (milliseconds). \emph{Episode shape:}
  \textbf{almost always single-turn} --- there is no environment
  feedback loop, just one answer. But that one response is a
  \textbf{long reasoning trace: \textasciitilde1k--8k tokens, and up to
  16k--32k} for hard competition problems on a reasoning model. Episode
  length ≈ that one response; strongly \textbf{inference-bound} (the
  checker is a one-line assert).
\item
  \textbf{Code generation \& competitive programming} ---
  \emph{Dataset:} MBPP {[}8{]}, HumanEval, \textbf{LiveCodeBench}
  {[}11{]}; competitive sets from \textbf{Codeforces, AtCoder, LeetCode}
  {[}10{]}. \emph{Reward:} \textbf{hidden unit tests pass} --- AceReason
  grants reward \emph{only if the code passes every test case} {[}10{]}.
  \emph{Metric:} pass@1 / pass@k, Codeforces ELO, LiveCodeBench (Kimi K2
  v6: \textbf{53.7\%} {[}5{]}). \emph{Env:} a language interpreter plus
  a \textbf{test-runner sandbox with per-case time limits}, run at scale
  (AceReason: a 64-worker process pool {[}10{]}; Kimi K2:
  \textgreater10k concurrent {[}5{]}). \emph{Episode shape:} \textbf{few
  turns --- often 1, up to \textasciitilde5} with retry on failed tests.
  The submitted program is short (\textbf{\textasciitilde100--800
  tokens}), but a \emph{reasoning} model prepends a long
  chain-of-thought, pushing the per-turn response to \textbf{several
  thousand tokens}; total episode is typically
  \textbf{\textasciitilde2k--16k tokens}, generation-heavy rather than
  context-heavy.
\item
  \textbf{Instruction following} --- \emph{Dataset:} synthetic prompts
  carrying \textbf{one to ten verifiable instructions} (Llama-Nemotron's
  RLOO stage {[}4{]}; part of Qwen3's 20+-task General RL {[}3{]}).
  \emph{Reward:} \textbf{deterministic verification via a code
  interpreter} for the checkable constraints, plus a \textbf{format
  reward}, with an LLM judge for the open-ended remainder {[}3, 4{]}.
  \emph{Metric:} IFEval-style constraint satisfaction. \emph{Env:} a
  \textbf{code interpreter} to evaluate the constraints. \emph{Episode
  shape:} \textbf{typically single-turn} --- one response scored against
  the constraints. Response length tracks the asked-for output,
  \textbf{\textasciitilde200 tokens--2k}; episode ≈ that one response.
\end{itemize}

\hypertarget{iterative-few-turn-families}{%
\subsubsection{Iterative, few-turn
families}\label{iterative-few-turn-families}}

A short retry loop against a checker --- write, run, read the failure,
fix. One family (GPU kernels) is where the \emph{checker} dominates
cost.

\begin{itemize}
\tightlist
\item
  \textbf{GPU kernels / systems} --- \emph{Dataset:} KernelBench
  {[}7{]}. \emph{Reward:} numerical \textbf{\texttt{allclose}} vs.~a
  reference \textbf{plus a speedup target}. \emph{Metric:} correctness
  rate, speedup. \emph{Env:} an \textbf{NVIDIA GPU + CUDA toolchain}
  (\texttt{nvcc}) --- the one family that requires a GPU, not just CPU.
  \emph{Episode shape:} \textbf{a few turns
  (\textsubscript{1--5)\textbf{, iterating on compile errors and
  speedup. Each response is one CUDA kernel, }}200--1,000 tokens}; total
  episode a \textbf{few thousand tokens}. Cost lives in the
  \emph{sandbox} (\texttt{nvcc} + \texttt{allclose} take seconds), not
  in generation --- the canonical sandbox-bound family.
\item
  \textbf{SQL / data analysis} --- \emph{Dataset:} text-to-SQL suites
  (Spider-style). \emph{Reward:} \textbf{result-set equality} against
  gold rows (the query is executed against a real engine).
  \emph{Metric:} execution accuracy. \emph{Env:} a local SQL engine
  (e.g.~SQLite), self-contained. \emph{Episode shape:} \textbf{1--3
  turns} (write query; retry on a SQL error or empty result-set).
  Response is a single query, \textbf{\textasciitilde30--300 tokens};
  episode is \textbf{small (a few hundred to \textasciitilde2k tokens)}
  --- cheap on both sides.
\end{itemize}

\hypertarget{multi-turn-observation-heavy-families}{%
\subsubsection{Multi-turn, observation-heavy
families}\label{multi-turn-observation-heavy-families}}

The response each turn is short, but the sandbox feeds a large
observation (tool result, retrieved passages, command output) back every
turn --- cost moves into \emph{context growth and the sandbox}.

\begin{itemize}
\tightlist
\item
  \textbf{Tool use / function calling} --- \emph{Dataset:} live and
  \textbf{simulated} tool environments; Kimi K2 builds a repository of
  \textbf{3,000+ real MCP tools plus 20,000+ synthetic tools} and
  generates multi-turn trajectories against a \textbf{tool simulator
  that acts as a world model} {[}5{]}; Qwen3 trains with
  real-environment execution feedback {[}3{]}. \emph{Reward:}
  \textbf{correct final result \emph{and} the right tools invoked}
  (verifiable), sometimes with a model judge. \emph{Metric:} task
  success on agentic tool-use benchmarks (Kimi K2: \textbf{τ²-Bench
  66.1, ACEBench 76.5} {[}5{]}). \emph{Env:} tool stubs / an API sandbox
  / a \textbf{stateful simulator} that returns realistic feedback across
  turns. \emph{Episode shape:} \textbf{multi-turn, one turn per tool
  call --- typically \textasciitilde3--10+ turns}. The policy's response
  each turn is short (a tool call plus brief reasoning,
  \textbf{\textasciitilde50--300 tokens}), but the \emph{observations}
  (tool results) can be large; total episode is
  \textbf{\textasciitilde2k--10k tokens}, dominated by fed-back tool
  output rather than by generation.
\item
  \textbf{RAG \& search agents} --- \emph{Dataset:} open-domain QA over
  a corpus or live web search (part of Qwen3's General RL, which
  explicitly targets RAG {[}3{]}). \emph{Reward:} \textbf{answer
  grounded in retrieved facts} --- Kimi K2 adds a \textbf{sentence-level
  faithfulness judge} {[}5{]}. \emph{Metric:} exact-match / F1.
  \emph{Env:} a retrieval index or search API plus \textbf{network
  egress}. \emph{Episode shape:} \textbf{multi-turn ---
  \textasciitilde2--8 turns} (issue query → read passages → refine or
  answer). The policy's responses (queries, final answer) are short, but
  each retrieved-document observation is large, so the \textbf{episode
  context is retrieval-heavy --- commonly 8k--32k+ tokens} even though
  the model generates comparatively little.
\item
  \textbf{Security / CTF} --- \emph{Dataset:} CTF challenge sets.
  \emph{Reward:} a \textbf{checker validates the captured flag}.
  \emph{Metric:} solve rate. \emph{Env:} an isolated shell, sometimes
  with network to a remote challenge service. \emph{Episode shape:}
  \textbf{multi-turn exploration --- \textasciitilde5--30+ turns} of
  shell commands. Per-turn responses are short (a command plus
  reasoning, \textbf{\textasciitilde50--400 tokens}); observations
  (command output) vary from a line to a hex dump, so the
  \textbf{episode ranges from a few thousand to tens of thousands of
  tokens}.
\end{itemize}

\hypertarget{long-horizon-session-bound-families}{%
\subsubsection{Long-horizon, session-bound
families}\label{long-horizon-session-bound-families}}

Tens to hundreds of turns; the accumulated transcript is the largest of
all, and the session must stay alive for the whole workflow --- the
families that most need a warm, long-lived sandbox.

\begin{itemize}
\tightlist
\item
  \textbf{Software engineering (SWE)} --- \emph{Dataset:} SWE-bench /
  SWE-bench Verified {[}6{]}; Kimi K2 additionally mines
  \textbf{executable unit tests from real GitHub pull requests and
  issues} {[}5{]}. \emph{Reward:} the repo's \textbf{test suite turns
  green} (\texttt{FAIL\_TO\_PASS} + \texttt{PASS\_TO\_PASS}), binary.
  \emph{Metric:} \% issues resolved (Kimi K2 reports \textbf{65.8\%}
  single-attempt on SWE-bench Verified {[}5{]}). \emph{Env:} a
  per-repository toolchain --- usually a \textbf{prebuilt Docker image}
  carrying the exact dependencies, run in isolation per task.
  \emph{Episode shape:} \textbf{long-horizon and multi-turn} --- real
  SWE agents take \textbf{\textasciitilde15--50+ turns} (locate → edit →
  run tests → read failure → re-edit); each turn's response is a command
  or a \textbf{diff/patch of \textasciitilde100--1,000 tokens}, but the
  observations (file contents, test logs) are large, so the
  \textbf{episode context routinely reaches tens of thousands of tokens
  (16k--128k+)} --- the most demanding family for context length.
\item
  \textbf{Web / computer use} --- \emph{Dataset:} WebArena-style browser
  tasks. \emph{Reward:} \textbf{goal state reached} (a DOM / state
  checker). \emph{Metric:} task success. \emph{Env:} a real
  \textbf{browser / rendering engine} and network. \emph{Episode shape:}
  \textbf{long-horizon --- \textasciitilde10--50+ turns}, one UI action
  (click/type/scroll) per turn. Each response is tiny
  (\textbf{\textasciitilde20--150 tokens}), but every observation (a DOM
  dump or accessibility tree, sometimes a screenshot) is very large,
  making this --- with SWE --- the \textbf{heaviest family for episode
  context (tens of thousands of tokens, growing per turn)}.
\item
  \textbf{Multi-agent / long-horizon} --- \emph{Dataset:} cooperative or
  long-workflow tasks. \emph{Reward:} a \textbf{verifiable end state}.
  \emph{Metric:} task success. \emph{Env:} many turns / long wall-clock,
  with state persisted across a long session. \emph{Episode shape:}
  \textbf{the longest horizon --- tens to hundreds of turns}, per-turn
  responses short-to-moderate, but the accumulated transcript makes the
  \textbf{episode the largest of all (often 32k--128k+ tokens)} and
  demands a session that stays alive for the whole workflow.
\end{itemize}

\hypertarget{language-and-environment-per-family-and-how-one-microvm-serves-them-all}{%
\subsection{Language and environment per family --- and how one microVM
serves them
all}\label{language-and-environment-per-family-and-how-one-microvm-serves-them-all}}

The families above differ sharply in \emph{what must be installed to run
and check them}. The table names, per family, the \textbf{language the
model's artifact is written in}, the \textbf{environment dependencies
the checker needs}, and whether it fits the \textbf{CPU microVM} or
requires the \textbf{GPU Instances} substrate. This is the concrete
``what runs inside'' that the sandbox image (or a per-session install)
must provide.

\begin{longtable}[]{@{}
  >{\raggedright\arraybackslash}p{(\columnwidth - 8\tabcolsep) * \real{0.2000}}
  >{\raggedright\arraybackslash}p{(\columnwidth - 8\tabcolsep) * \real{0.2000}}
  >{\raggedright\arraybackslash}p{(\columnwidth - 8\tabcolsep) * \real{0.2000}}
  >{\raggedright\arraybackslash}p{(\columnwidth - 8\tabcolsep) * \real{0.2000}}
  >{\raggedright\arraybackslash}p{(\columnwidth - 8\tabcolsep) * \real{0.2000}}@{}}
\toprule\noalign{}
\begin{minipage}[b]{\linewidth}\raggedright
Family (§3)
\end{minipage} & \begin{minipage}[b]{\linewidth}\raggedright
Dataset(s)
\end{minipage} & \begin{minipage}[b]{\linewidth}\raggedright
Artifact language
\end{minipage} & \begin{minipage}[b]{\linewidth}\raggedright
Environment / dependencies the checker needs
\end{minipage} & \begin{minipage}[b]{\linewidth}\raggedright
Substrate
\end{minipage} \\
\midrule\noalign{}
\endhead
\bottomrule\noalign{}
\endlastfoot
Math \& logical reasoning & GSM8K, MATH-500, DeepScaleR & Python (answer
string) & Python stdlib + \texttt{sympy} (symbolic/numeric match) &
CPU \\
Code gen \& competitive & MBPP, HumanEval, LiveCodeBench & Python
(competitive: often C++) & Python interpreter + test runner
(\texttt{unittest}/asserts) + per-case timeout & CPU \\
Instruction following & synthetic (IFEval-style) & text / Python & code
interpreter (stdlib) to verify each constraint & CPU \\
GPU kernels / systems & KernelBench & \textbf{CUDA C++} + Python driver
& \textbf{NVIDIA GPU + CUDA toolkit (\texttt{nvcc}) + PyTorch}
(\texttt{allclose}, speedup) & \textbf{GPU} \\
SQL / data analysis & Spider-style, Gretel text-to-SQL & SQL (+ Python
harness) & a SQL engine --- \texttt{sqlite3} (in the Python stdlib) &
CPU \\
Tool use / function calling & xLAM, ToolACE & Python / JSON tool calls &
Python + tool stubs / API sandbox / a \textbf{stateful simulator}; MCP
via Gateway & CPU \\
RAG \& search agents & HotpotQA, open-domain QA & Python (queries +
answer) & a retrieval index / search API + \textbf{network egress} &
CPU \\
Security / CTF & CTF challenge sets & shell + C / Python / binaries & an
isolated shell + common tooling, sometimes \textbf{network} to a
challenge service & CPU \\
Software engineering (SWE) & SWE-bench / Verified & Python (repo
language varies) & \textbf{per-repo toolchain} --- pinned deps + system
packages, \texttt{pytest}; a large image & CPU (amd64) \\
Web / computer use & WebArena-style & Python + browser actions & a
\textbf{headless browser / rendering engine} + network & CPU \\
Multi-agent / long-horizon & cooperative / long-workflow & varies & the
sub-tasks' deps + a \textbf{long-lived session} holding state & CPU \\
\end{longtable}

\textbf{How can one microVM support such different environments?} The
key is that \textbf{the microVM does not need one universal image} ---
it reconciles the variation three ways, in increasing order of how
structural the difference is (this is the §``One image per Runtime, but
a different env per row'' logic applied to \emph{languages/deps} rather
than to per-repo images):

\begin{enumerate}
\def\labelenumi{\arabic{enumi}.}
\tightlist
\item
  \textbf{Task-agnostic sandbox → most differences are just
  \emph{files}, not images.} The sandbox exposes only
  \texttt{write\_file} / \texttt{exec} / \texttt{read\_file}, and every
  session is an \textbf{isolated microVM} booting \emph{your} OCI image.
  All task specificity --- the prompt, seed files, and the hidden
  checker --- is \emph{written into that episode's session at seed
  time}, not baked into the image. So a \textbf{single generic Python
  image} already serves the CPU families that need nothing beyond the
  interpreter and stdlib: math (\texttt{sympy}), code
  (\texttt{unittest}), SQL (\texttt{sqlite3} is stdlib), instruction
  following, tool-use stubs, and light CTF. A math episode, a SQL
  episode, and a tool-use episode are just \emph{different files} in
  three identical, isolated microVMs.
\item
  \textbf{Per-session runtime install → heavier or unusual deps go into
  the ephemeral workspace.} This has two variants with \emph{different}
  implementation status. \textbf{(a) Deferred to \texttt{CMD}} --- the
  sandbox image installs a large dep when the container boots, before
  the server binds, keeping the \emph{built} image small; this is
  \textbf{implemented} and is how the removed GPU demos' sandbox
  installed the \textasciitilde2.7 GB CUDA torch wheel to stay under the
  2 GB image cap. \textbf{(b) Installed at \emph{seed time}} --- the
  harness would \texttt{git\ clone} / \texttt{pip\ install} /
  \texttt{apt} a family's deps into that session's workspace on turn 0.
  This is a \textbf{described design, not yet implemented}: the shipped
  demos seed only small files (\texttt{PROBLEM.md},
  \texttt{run\_checks.py}, \texttt{solution.py}), and the
  seed-time-install path is the SWE workaround the report explicitly
  flags as \emph{``a design answer, not a measured one''} (§``One image
  per Runtime, but a different env per row''). Either way, because each
  \texttt{runtimeSessionId} is a \textbf{separate microVM}, one
  episode's installed environment \textbf{never leaks into another's}
  --- the isolation is what would make per-session customization safe.
\item
  \textbf{Per-capability image → a \emph{pool} of Runtimes when the
  difference is structural.} Some differences cannot be installed at
  runtime: a \textbf{GPU toolchain} (KernelBench needs \texttt{nvcc} + a
  GPU --- the \textbf{Instances/amd64} substrate, not the CPU microVM),
  a \emph{pinned OS image} per task (Terminal-Bench), root \texttt{apt}
  the base image can't grant, or an image over the 2 GB microVM cap. For
  those you register \textbf{one Runtime ARN per image}, and routing is
  now \textbf{implemented}: every task carries an \texttt{image} tag,
  and \texttt{make\_sandbox(image=…)} →
  \texttt{\_runtime\_arn\_for\_image} looks it up in
  \texttt{AGENTCORE\_SANDBOX\_RUNTIME\_ARN\_BY\_IMAGE} (falling back to
  round-robin over \texttt{AGENTCORE\_SANDBOX\_RUNTIME\_ARNS}).
  AgentCore itself never sees the task --- the routing is entirely
  client-side (the Terminal-Bench demo's ``Strict routing'' design point
  walks the chain, §``The Terminal-Bench demo'').
\end{enumerate}

In short: \textbf{most task diversity needs no task-specific image.} The
managed AgentCore \textbf{Code Interpreter} (the training default until
2026-09-28, a probe since) already ships Python 3.12 with numpy, pandas,
scipy, scikit-learn, sympy, torch (CPU), sqlite3, gcc and node, and a
live probe (2026-09-25) passed \textbf{all 13 task families} at reward
1.0 on it (\#1); since the agent loop moved into the sandbox session,
the demos run the same tasks on one generic own-Runtime image instead,
because the loop must reach the trainer's adapter. Only genuine
image/substrate differences --- a pinned per-task OS (Terminal-Bench),
an image over 2 GB, git/internet-dependent installs, or CUDA --- need
your own Runtime (\#3); seed-time installs (\#2b) remain a design for
SWE workloads, not a path the current demos exercise.

\hypertarget{the-one-family-that-stresses-the-image-swe-and-the-microvms-arm64-2-gb-limits}{%
\subsection{The one family that stresses the image: SWE and the
microVM's arm64 + 2 GB
limits}\label{the-one-family-that-stresses-the-image-swe-and-the-microvms-arm64-2-gb-limits}}

Most families above run inside a small, generic image --- a Python
interpreter and a stdlib checker fit comfortably. \textbf{Software
engineering (SWE)} is the exception, and it is worth calling out before
the platform section because it collides with two properties of
AgentCore's \emph{managed serverless microVM} --- and the two collide
with each other:

\begin{itemize}
\tightlist
\item
  \textbf{arm64-only (documented).} The managed microVM runs
  \textbf{only arm64} images: the official Runtime service contract
  states the \emph{``Platform: ARM64 container --- Required for
  compatibility with the AgentCore Runtime environment''} {[}D17{]} (the
  GPU capacity-provider substrate, by contrast, is amd64 --- §``Building
  a customized Runtime''). A per-repository SWE image must therefore be
  \textbf{cross-built to arm64}.
\item
  \textbf{A tight built-image ceiling.} The managed microVM image must
  stay small --- this project measured a \textbf{\textasciitilde2 GB
  (2048 MB)} ceiling on the built image, which is exactly why its GPU
  sandbox installs the CUDA torch wheel at \texttt{CMD} rather than
  baking it in (§``Building a customized Runtime''). A realistic
  SWE-bench repo image --- the exact compiler, system packages, and
  pinned dependencies a project needs --- routinely exceeds a couple of
  GB.
\end{itemize}

The two interact badly for SWE specifically, because \textbf{SWE-bench's
own tooling treats arm64 as second-class}: its README states plainly
that \emph{``Support for \texttt{arm64} machines is experimental''}, and
on ARM hosts it does not ship prebuilt images --- you must pass
\texttt{-\/-task-repo} so \emph{``the images are built locally with
Docker Buildx''} {[}6b{]}. So the very images you would reuse are (a)
only \emph{experimentally} available on the architecture the microVM
forces, and (b) sized for x86 developer boxes, not a 2 GB serverless
cap. The heaviest-image family lands on the substrate least suited to
it.

\textbf{Workarounds --- in rough order of how well they fit this
project's architecture:}

\begin{enumerate}
\def\labelenumi{\arabic{enumi}.}
\tightlist
\item
  \textbf{Fetch the repo and deps at \emph{seed} time, not in the
  image.} The sandbox is task-agnostic (\texttt{write\_file} /
  \texttt{exec} / \texttt{read\_file}), so the harness can
  \texttt{git\ clone} the target repo and \texttt{pip\ install}/build
  its dependencies into the \emph{session workspace} on the first turn,
  instead of shipping them in the OCI image. This moves the repo bulk
  off the image entirely and into the per-session ephemeral filesystem
  (good for 8 h) --- sidestepping the image cap. Cost: per-episode setup
  latency and \textbf{network egress}, so the Runtime needs public
  egress or VPC interface endpoints (§``Scaling, sharing, and network
  isolation''; {[}D11{]}).
\item
  \textbf{Defer heavy installs to \texttt{CMD}.} Keep the \emph{built}
  image lean and \texttt{pip}/\texttt{apt} the large layers when the
  session boots --- the trick the removed GPU sandbox used for torch
  (§``Building a customized Runtime''). Cost: added cold-start (the
  removed GPU demo paid \textasciitilde90 s for its torch install).
\item
  \textbf{Slim the repo image.} Multi-stage builds,
  \texttt{-\/-no-install-recommends}, drop build toolchains and package
  caches after compiling, and start from a slim arm64 base --- standard
  Docker hygiene, often enough to bring a repo image under the cap.
\item
  \textbf{Cross-build to arm64 with Buildx, and \emph{validate the
  checker there}.}
  \texttt{docker\ buildx\ build\ -\/-platform\ linux/arm64}; for
  SWE-bench use its \texttt{-\/-task-repo} local-build path. Because
  arm64 is experimental upstream, re-run the
  \texttt{FAIL\_TO\_PASS}/\texttt{PASS\_TO\_PASS} checker on arm64
  before trusting the reward --- an arch-specific test flake would
  silently corrupt the RLVR signal.
\item
  \textbf{Move SWE to the amd64 GPU/Instances substrate.} Instances run
  \textbf{amd64}, allow far larger images, and add a \textbf{persistent
  EBS volume} (so a one-time repo/deps install survives stop/resume) and
  \textbf{14-day} sessions (§``Two compute types''). This sidesteps both
  the arch and the image cap at once; the cost is spending an EC2
  instance (and its GPU) that a CPU-only SWE checker does not otherwise
  need.
\end{enumerate}

This is the concrete reason the §``Can AgentCore host every frontier
sandbox?'' table marks the per-repo Docker path with a ◑ rather than a
clean ✓: for SWE, the sandbox \emph{image} --- not the checker or the
GPU --- is the binding constraint.

\hypertarget{one-image-per-runtime-but-a-different-env-per-row-how-agentcore-reconciles-it}{%
\subsection{One image per Runtime, but a different env per row: how
AgentCore reconciles
it}\label{one-image-per-runtime-but-a-different-env-per-row-how-agentcore-reconciles-it}}

SWE-bench raises a sharper version of the same problem: \textbf{each
task instance ships its own Docker environment} (a specific interpreter
version, system packages, and pinned dependencies checked out at one
commit) --- hundreds of distinct images across the benchmark.
AgentCore's binding constraint is that a \textbf{Runtime ARN is bound to
exactly one container image}, chosen once at \texttt{CreateAgentRuntime}
(\texttt{containerConfiguration.containerUri}, an ECR image ---
§``Building a customized Runtime''). An \texttt{invoke\_agent\_runtime}
call passes only a \textbf{\texttt{runtimeSessionId} + payload --- never
an image {[}D8{]}}, so you \textbf{cannot swap the Docker environment
per invocation}. Three ways reconcile ``one image per ARN'' with ``a
different env per row'' --- and the practical answer is usually a
hybrid:

\begin{itemize}
\tightlist
\item
  \textbf{A --- Move the per-row difference \emph{out of the image, into
  per-session setup} (this project's native model).} The sandbox is
  task-agnostic (\texttt{write\_file} / \texttt{exec} /
  \texttt{read\_file}) and all task specificity already lives in files
  the harness writes into \emph{that episode's session} at seed time.
  Extend that: on turn 0 the harness runs \texttt{git\ clone} +
  \texttt{git\ checkout\ \textless{}base\_commit\textgreater{}},
  installs the pinned deps, and drops in the hidden
  \texttt{run\_checks.py} --- so ``a different image per row'' becomes
  ``a different \textbf{setup script} per row, run inside one generic
  image.'' Each \texttt{runtimeSessionId} is an isolated microVM, so row
  \emph{i}'s installed env never leaks into row \emph{j}. \emph{Cost:}
  per-episode setup latency + \textbf{network egress} (public or VPC
  endpoints, {[}D11{]}). \emph{Limit:} works when the delta is
  \emph{installable at runtime}; it struggles when a row needs a
  \textbf{different interpreter version}, \textbf{root \texttt{apt}} the
  base image can't grant, or a heavy compiled toolchain --- those belong
  in the image.
\item
  \textbf{B --- A \emph{pool} of Runtimes, one image per ARN, routed by
  row.} Register \textbf{K Runtime ARNs, each with a different baked
  image}, and have the harness pick the ARN by the row's repo/family
  instead of round-robin (the \emph{capability-routed} dispatch the
  Terminal-Bench demo uses, §``The Terminal-Bench demo''). But hundreds
  of per-instance images cannot each get a Runtime (operationally heavy,
  per-region limits), so group at \textbf{repo-family} granularity (one
  \texttt{django} image, one \texttt{sympy} image, \ldots), \textbf{not}
  per-instance.
\item
  \textbf{C --- The practical hybrid (family base image + per-row
  delta).} Bake the \emph{expensive, shared} layer per repo-family ARN
  (interpreter version, system libs, repo at a common base) → B; install
  the \emph{cheap per-commit delta} (that instance's pip deps) at seed →
  A. This keeps each image small --- under the microVM's
  \textasciitilde2 GB arm64 cap above --- while one ARN still covers
  many instances. If images are simply too large or too x86-bound, fall
  back to the amd64 \textbf{Instances} substrate (workaround \#5).
\end{itemize}

\emph{Status:} A--C above are a \textbf{design answer}. The project also
built a fourth, simpler route and then removed it on 2026-09-30: one
Runtime per SWE-bench instance on its Epoch AI arm64 image (§``History:
removed demos (2026-09-30)'').

\hypertarget{foundations-of-agentcore-runtime}{%
\section{Foundations of AgentCore
Runtime}\label{foundations-of-agentcore-runtime}}

Amazon Bedrock AgentCore {[}D1{]} is a family of managed services for
running and securing agents and tools --- any framework, any model,
fully serverless. This report uses one of them, \textbf{AgentCore
Runtime}, as the execution \textbf{sandbox} for RL. This section defines
its three core concepts, its two compute types, how AgentCore manages
their lifecycle, and how you build and deploy a customized Runtime ---
cross-checking the official documentation {[}D1--D13, D17{]} against
AWS's own SDK and sample code {[}D14--D16{]}.

\hypertarget{the-microvm-at-a-glance-what-it-is-what-runs-in-it-what-it-supports-how-you-customize-it}{%
\subsection{The microVM at a glance --- what it is, what runs in it,
what it supports, how you customize
it}\label{the-microvm-at-a-glance-what-it-is-what-runs-in-it-what-it-supports-how-you-customize-it}}

Before the precise definitions, a compact answer to the four questions a
practitioner asks first. Each is expanded in the subsection named after
it.

\begin{itemize}
\tightlist
\item
  \textbf{What is the microVM?} The default AgentCore Runtime substrate
  is a \textbf{Firecracker-class, hardware-virtualized microVM} --- a
  lightweight VM with its \textbf{own guest kernel, CPU, memory, and
  filesystem}, one per session, that still boots \textbf{your OCI
  container image} inside it and is \textbf{sanitized on termination}.
  It gives VM-strength isolation for untrusted, model-generated code
  while keeping container packaging. It is \textbf{CPU-only} (no GPU
  {[}D5{]}), runs an \textbf{ARM64} image serving
  \texttt{POST\ /invocations} + \texttt{GET\ /ping} on
  \texttt{0.0.0.0:8080}, and lasts up to \textbf{8 hours} with a
  15-minute idle timeout. When a task needs a GPU, Runtime switches to
  the \textbf{Instances} (EC2) substrate instead (§``Two compute
  types''). \emph{Details:} §``Runtime, session, and microVM'' and
  §``microVM vs a Docker container vs E2B''.
\item
  \textbf{What environment / dependencies does it have?} Only what
  \textbf{your image} provides --- there is no fixed ``AgentCore OS''
  you inherit. A minimal sandbox is a slim Python base plus the
  \texttt{bedrock-agentcore} SDK; the built image must stay
  \textbf{ARM64 and under \textasciitilde2 GB}. You add dependencies
  three ways, in increasing weight: bake them into the image;
  \textbf{install at \texttt{CMD}} when the container boots (how the GPU
  image adds the large CUDA torch wheel without blowing the 2 GB cap);
  or \textbf{install at seed time} into the per-session ephemeral
  workspace (\texttt{git\ clone}/\texttt{pip}/\texttt{apt} on turn 0,
  needing egress). Each session's installs are isolated and wiped on
  termination. \emph{Details:} §``Language and environment per family''
  and §``Building a customized Runtime''.
\item
  \textbf{What tasks can it support?} Effectively every \textbf{CPU}
  sandbox family frontier labs train against: code + unit tests,
  math/SQL checkers, instruction following, tool-use simulators, RAG,
  and light CTF all fit one \textbf{generic} image (task specificity is
  just \emph{files} seeded per session). \textbf{SWE-bench}-style
  per-repo images fit with workarounds (arm64 cross-build, seed-time
  clone, or moving to the amd64 Instances substrate). The one family the
  microVM \textbf{cannot} serve is \textbf{GPU kernels}
  (\texttt{nvcc}/\texttt{allclose}) --- that needs the GPU Instances
  substrate. \emph{Details:} §``Language and environment per family''
  and §``Can AgentCore host every frontier sandbox?''.
\item
  \textbf{How do you customize it as a Runtime?} Author a container that
  satisfies the contract (ARM64; \texttt{0.0.0.0:8080};
  \texttt{/invocations} + \texttt{/ping}), cross-build it with
  \texttt{docker\ buildx\ -\/-platform\ linux/arm64}, push to ECR, and
  register it with \texttt{CreateAgentRuntime}'s
  \texttt{containerConfiguration.containerUri}. AWS also offers a
  lower-effort \textbf{\texttt{agentcore} CLI} (CodeZip or Container
  build) and a \textbf{\texttt{bedrock-agentcore} Python SDK} whose
  \texttt{@app.entrypoint} implements the HTTP contract for you.
  \emph{Details:} §``Building a customized Runtime''.
\end{itemize}

\hypertarget{runtime-session-and-microvm-three-concepts-precisely}{%
\subsection{Runtime, session, and microVM --- three concepts,
precisely}\label{runtime-session-and-microvm-three-concepts-precisely}}

These three words are used loosely elsewhere; the integration depends on
keeping them distinct.

\begin{itemize}
\tightlist
\item
  \textbf{Runtime (agent runtime).} A \emph{versioned, containerized
  host} you create with \texttt{CreateAgentRuntime} {[}D9{]}. It
  packages your code --- a container image in ECR, or a zipped Python
  bundle --- behind an ARN. Creating a Runtime auto-creates version
  \textbf{V1} and a \textbf{DEFAULT endpoint}; you can publish further
  versions and custom endpoints and route invocations to them. The
  Runtime is the \emph{deployable unit}; it holds no per-user state of
  its own.
\item
  \textbf{Session.} One logical task or conversation, addressed by a
  \textbf{\texttt{runtimeSessionId}} you pass on
  \texttt{InvokeAgentRuntime} {[}D8{]} (client-supplied or
  Runtime-generated; \textbf{minimum 33 characters}, max 256, carried in
  the \texttt{X-Amzn-Bedrock-AgentCore-Runtime-Session-Id} header). The
  \textbf{first} invoke for a new id \textbf{cold-starts} dedicated
  compute; every \textbf{subsequent} invoke with the same id
  \textbf{re-routes to that same warm compute}, so files, memory, and
  processes persist \emph{across turns within an episode}. Different ids
  are fully isolated.
\item
  \textbf{microVM.} The \emph{unit of isolation} backing a session on
  the default compute: a lightweight virtual machine with its
  \textbf{own CPU, memory, and filesystem}, sanitized on termination.
  One session maps to one microVM; a microVM is never shared across
  sessions. Its container contract {[}D17{]}: an \textbf{ARM64} image
  listening on \textbf{\texttt{0.0.0.0:8080}} that serves
  \textbf{\texttt{POST\ /invocations}} (the work) and
  \textbf{\texttt{GET\ /ping}} (health, returning
  \texttt{\{"status":\ "Healthy"\}} or \texttt{"HealthyBusy"}) --- the
  same contract AWS's SDK scaffolds for you {[}D15{]}.
\end{itemize}

\textbf{The cardinality.} These three nest \textbf{one-to-many-to-one}:
\textbf{one Runtime serves many concurrent sessions, and each session
maps to exactly one microVM (1 : N : 1)}. The Runtime is a stateless
\emph{template} (one image, one ARN); every distinct
\texttt{runtimeSessionId} you invoke against it mints its own isolated
microVM, and AgentCore's capacity provider places each one for you ---
so the session id, not a second Runtime, is how you get a second
sandbox. You create \emph{additional} Runtimes only when the
\textbf{image}, the per-region \textbf{session cap}, or the
\textbf{AZ/account} must differ (§``Scaling the loop''); otherwise a
single Runtime backs all your concurrent sessions. Put another way:
\textbf{the number of Runtimes tracks the number of distinct sandbox
\emph{images} (capabilities), not the number of tasks} --- because the
sandbox is task-agnostic (\texttt{write\_file} / \texttt{exec} /
\texttt{read\_file}), one generic image serves every task family that
runs on the same substrate, so even a mixed-family training run stays on
\textbf{one Runtime, many sessions}.

In RL terms: \textbf{the Runtime is the deployable sandbox image, the
session is one episode, and the microVM is the isolated box that episode
executes in.} So \emph{N} concurrent episodes = \emph{N} session ids on
one Runtime --- not \emph{N} Runtimes.

\hypertarget{a-mixed-family-batch-how-many-runtimes-sessions-and-microvms-and-who-picks-the-image}{%
\subsubsection{A mixed-family batch --- how many Runtimes, sessions, and
microVMs, and who picks the
image?}\label{a-mixed-family-batch-how-many-runtimes-sessions-and-microvms-and-who-picks-the-image}}

The sharpest way to feel the distinction is a \textbf{heterogeneous
rollout step}: a mixed-family dataset (the removed
\texttt{demo\_cpu\_tasks} with a comma-list \texttt{DATASET=},
e.g.~\texttt{math,sql,ctf,tool\_use}) draws a batch whose rows span
several families --- say a math problem, a text-to-SQL query, a CTF
challenge, and a tool-use task --- with \texttt{n-samples-per-prompt}
episodes each. A natural worry is that four task types need four sandbox
images, so four Runtimes. They do not. Count each concept:

\begin{itemize}
\tightlist
\item
  \textbf{Runtimes: one} (a small pool of \emph{K} only for throughput,
  not for variety). Every one of these families is a \textbf{pure-CPU /
  stdlib checker} --- \texttt{sympy}, \texttt{unittest},
  \texttt{sqlite3}, an isolated shell, tool stubs --- so a
  \textbf{single generic CPU image} runs and checks all of them. The
  Runtime count tracks distinct \emph{images} (capabilities), and here
  there is exactly one capability. You would add a second Runtime only
  for a genuinely different image --- the \textbf{GPU-kernel} family
  (\texttt{nvcc}, a different substrate) --- not for a different CPU
  family.
\item
  \textbf{Sessions: one per in-flight episode, recycled.} If the batch
  is \texttt{RB} prompts × \texttt{NS} samples, that is \texttt{RB×NS}
  \emph{episodes}, but they do \textbf{not} all run at once: only
  \texttt{S\ =\ CONC\ ×\ num\_engines} are concurrent, and those few
  sessions \textbf{stream through all \texttt{RB×NS} episodes} over the
  step. A math episode and a SQL episode are just two different session
  ids against the same Runtime.
\item
  \textbf{microVMs: one per session, 1:1.} Each
  \texttt{runtimeSessionId} mints its own isolated microVM (booting that
  same generic image), sanitized when the session ends. So at any
  instant there are \texttt{S} live microVMs, whatever mix of families
  they happen to be running.
\end{itemize}

\textbf{So how do slime and AgentCore know which image a given data row
needs? They don't --- and by design, they don't have to.} No per-row
image selection happens anywhere, because the image is
\emph{task-agnostic}. The routing that \emph{does} happen is a division
of labor:

\begin{itemize}
\tightlist
\item
  \textbf{slime / the harness picks the \emph{task}, not the image.}
  Each dataset row carries \texttt{metadata.task.family}; the generate
  hook (\texttt{\_task\_from\_sample}) rebuilds the matching
  \textbf{Task subclass} --- its prompt, seed files, and hidden checker
  --- and the harness \textbf{writes those task-specific files into that
  episode's session} at seed time (\texttt{write\_file}). The family
  discriminator chooses \emph{what gets written into the box}, never
  \emph{which box}.
\item
  \textbf{AgentCore never sees a task at all.} An
  \texttt{invoke\_agent\_runtime} call carries only a
  \textbf{\texttt{runtimeSessionId} + payload --- never an image}
  {[}D8{]}; the image was fixed once, at \texttt{CreateAgentRuntime}.
  AgentCore's only job is to route the session id to its microVM. It is
  oblivious to whether the payload is a math check or a SQL query.
\end{itemize}

In other words, ``a different environment per row'' is realized as
\textbf{a different set of \emph{files} per session inside one shared
image}, not a different image per row. The one time a row truly needs a
different \emph{image} --- a CPU family vs.~the GPU-kernel family ---
that is a \textbf{substrate} choice made per run by an env var, not per
row by a router; genuine per-row \emph{image} routing is what the
Terminal-Bench demo does, one Runtime per task (§``The Terminal-Bench
demo''). This is exactly why a mixed-family run stays on \textbf{one
Runtime, \texttt{S} sessions, \texttt{S} microVMs}.

\hypertarget{microvm-vs-a-docker-container-vs-e2b}{%
\subsection{microVM vs a Docker container vs
E2B}\label{microvm-vs-a-docker-container-vs-e2b}}

Because ``sandbox'' and ``container'' are often used interchangeably, it
is worth being precise about what the microVM actually is --- especially
since we are running \textbf{untrusted, model-generated code}, where the
strength of the isolation boundary is the whole point.

\begin{itemize}
\tightlist
\item
  \textbf{A Docker container} isolates via Linux namespaces and cgroups
  but \textbf{shares the host kernel}. That is fine for packaging and
  for trusted workloads, but a shared kernel is a comparatively soft
  boundary for hostile code: a kernel exploit escapes to the host and to
  co-tenants. Containers alone are not a safe multi-tenant boundary for
  arbitrary generated code.
\item
  \textbf{An AgentCore microVM} is a lightweight,
  \textbf{hardware-virtualized} virtual machine (a Firecracker-class
  microVM) with its \textbf{own guest kernel}, its own
  CPU/memory/filesystem, and a VM-strength boundary --- while still
  booting \textbf{your OCI container image} inside it. So it is not
  ``container \emph{or} VM'': you keep container packaging \emph{and}
  gain VM isolation, provisioned per session and sanitized on
  termination. On top of that it is \textbf{AWS-native} --- reached only
  by in-account SigV4/IAM, placeable in your VPC, and integrated with
  the other AgentCore services (Gateway, Memory, Identity,
  Observability) --- and it can carry a \textbf{GPU} (via Instances) and
  \textbf{long sessions} (8 h microVM / 14 d Instances).
\item
  \textbf{E2B} is an independent sandbox platform that, like AgentCore,
  runs code in \textbf{Firecracker microVMs} (its templates are also
  defined from Dockerfiles), so the \emph{isolation model is
  comparable}. The differences are operational, not architectural: E2B
  is its own cloud (or self-hosted) service with its own SDK and
  identity, oriented to short-lived code-interpreter sandboxes, whereas
  AgentCore is a managed AWS service that runs \textbf{in your own
  account} under IAM/VPC, adds \textbf{GPU on Instances} and
  \textbf{multi-hour/multi-day sessions}, and slots into the rest of the
  Bedrock AgentCore stack.
\end{itemize}

\begin{longtable}[]{@{}
  >{\raggedright\arraybackslash}p{(\columnwidth - 6\tabcolsep) * \real{0.2500}}
  >{\raggedright\arraybackslash}p{(\columnwidth - 6\tabcolsep) * \real{0.2500}}
  >{\raggedright\arraybackslash}p{(\columnwidth - 6\tabcolsep) * \real{0.2500}}
  >{\raggedright\arraybackslash}p{(\columnwidth - 6\tabcolsep) * \real{0.2500}}@{}}
\toprule\noalign{}
\begin{minipage}[b]{\linewidth}\raggedright
\end{minipage} & \begin{minipage}[b]{\linewidth}\raggedright
Docker container
\end{minipage} & \begin{minipage}[b]{\linewidth}\raggedright
AgentCore microVM
\end{minipage} & \begin{minipage}[b]{\linewidth}\raggedright
E2B sandbox
\end{minipage} \\
\midrule\noalign{}
\endhead
\bottomrule\noalign{}
\endlastfoot
Isolation boundary & shared host kernel (namespaces/cgroups) &
\textbf{own kernel} --- Firecracker-class microVM & own kernel ---
Firecracker microVM \\
Runs your image? & yes (OCI) & yes --- your OCI image \emph{inside} a
microVM & Dockerfile-defined templates \\
Managed by & you / your orchestrator & \textbf{AWS, in your account}
(IAM/SigV4/VPC) & E2B service (cloud or self-host) \\
GPU & host-dependent & \textbf{via Instances} (EC2) & not the core
offering \\
Session length & you manage & 8 h microVM / 14 d Instances & short-lived
(code-interpreter) \\
\end{longtable}

For an RL sandbox running untrusted policy output, the microVM's
per-session hardware isolation is the property that matters: one
episode's generated code cannot see another episode's files, the
trainer's weights, or the host --- and the box is destroyed and its
memory wiped when the session ends.

\hypertarget{two-compute-types-managed-serverless-microvm-vs-instances-gpu-ec2}{%
\subsection{Two compute types --- managed serverless (microVM) vs
Instances (GPU
EC2)}\label{two-compute-types-managed-serverless-microvm-vs-instances-gpu-ec2}}

AgentCore Runtime offers two substrates, chosen at create time; the loop
code is identical across them.

\begin{itemize}
\tightlist
\item
  \textbf{Serverless microVM (default).} No servers to manage: AgentCore
  provisions a per-session microVM on demand. \textbf{CPU only} --- a
  microVM has \emph{no GPU} {[}D5{]}. Sessions last up to \textbf{8
  hours} with a \textbf{15-minute idle timeout} (reset on each invoke).
  This fits the large majority of RL sandboxes: code execution + unit
  tests, math/SQL checkers, tool simulators.
\item
  \textbf{Instances (GPU).} When a task needs a GPU, Runtime runs on the
  \textbf{Instances} compute type --- real \textbf{EC2 instances in your
  own account}, provisioned through a \textbf{capacity provider}
  (\texttt{capacityProviderConfiguration} → \texttt{ec2Configuration})
  {[}D5, D9{]}. Supported GPU families include \textbf{g4dn, g5, g6,
  g6e, gr6, g6f, gr6f, g7e} (plus \texttt{inf2}); AgentCore provisions
  the NVIDIA drivers, so a standard CUDA image works unmodified.
  Sessions can run up to \textbf{14 days}, a \textbf{persistent EBS
  volume} survives a session stop and re-attaches on resume, and ---
  unlike the microVM --- one instance/session can host \textbf{many
  agents} (1:N). Crucially, one \texttt{runtimeSessionId} still maps to
  exactly \textbf{one instance = one GPU}, so there is \textbf{no
  cross-session GPU time-slicing}: you scale GPUs by adding
  sessions/instances, not by fractioning one.
\end{itemize}

\begin{quote}
\textbf{What about a multi-GPU instance --- two sessions, one GPU each?}
No, not natively: \textbf{one AgentCore session = one whole instance}
(session→instance is 1:1; the capacity provider hands you the
\emph{instance}, not a GPU). So you \textbf{cannot} point two distinct
\texttt{runtimeSessionId}s at the \emph{same} 2-GPU box and expect one
GPU each --- two session ids try to acquire two \emph{separate}
instances, and you never get to pin a session to a particular GPU on a
shared box. To put several episodes on the GPUs of \emph{one} multi-GPU
instance you must instead use \textbf{one shared session} plus a
\textbf{self-managed in-instance farm} (a Ray farm keyed by per-episode
\texttt{sub\_session}), which is what the farm demo
\texttt{demo\_gpu\_farm} did until its removal on 2026-09-30 ---
\texttt{B4\_GPU\_FRACTION=1.0} with \texttt{B4\_FARM\_GPUS=N} leased
\textbf{one whole GPU per sub-session} (\texttt{FARM=whole}), and
\texttt{B4\_GPU\_FRACTION=0.5} on a single GPU packed two fractional
episodes onto it (\texttt{FARM=fractional}) --- no MIG and no
cross-session time-slicing required, but also no VM boundary between
episodes. The shipped demos use no Instance at all (the removed
\texttt{demo\_gpu\_single} ran one session per g6).
\end{quote}

\begin{quote}
\textbf{Can you pick a CPU-only EC2 instance for an Instance?} Not as a
documented path. The \texttt{ec2Configuration} families AWS lists {[}D5,
D9{]} are all \textbf{accelerated} (g4dn/g5/g6/g6e/gr6/g6f/gr6f/g7e +
\texttt{inf2}) --- there is \textbf{no general-purpose (m/c/t) CPU
family}. The split is deliberate: \textbf{CPU-only work belongs on the
serverless microVM}, and the Instances substrate exists to add a
\textbf{GPU}. So if you want a larger/x86/longer-lived \emph{CPU} box
(e.g.~to escape the microVM's arm64 + 2 GB image limits, §3), the
documented route is to launch a GPU-family Instance and leave the GPU
idle --- you pay for an accelerator you don't use. (The supported list
is per the docs as of this work; re-check \texttt{ec2Configuration}. In
practice this project's account accepts an \texttt{m5.xlarge} capacity
provider: the removed \texttt{demo\_cpu\_instance} ran one m5 Instance
per session on \texttt{acs\_cpu\_cp\_instance}.)
\end{quote}

\begin{longtable}[]{@{}
  >{\raggedright\arraybackslash}p{(\columnwidth - 4\tabcolsep) * \real{0.3333}}
  >{\raggedright\arraybackslash}p{(\columnwidth - 4\tabcolsep) * \real{0.3333}}
  >{\raggedright\arraybackslash}p{(\columnwidth - 4\tabcolsep) * \real{0.3333}}@{}}
\toprule\noalign{}
\begin{minipage}[b]{\linewidth}\raggedright
Property
\end{minipage} & \begin{minipage}[b]{\linewidth}\raggedright
Serverless microVM
\end{minipage} & \begin{minipage}[b]{\linewidth}\raggedright
Instances (GPU)
\end{minipage} \\
\midrule\noalign{}
\endhead
\bottomrule\noalign{}
\endlastfoot
Compute & managed, per session & EC2 in your account (capacity
provider) \\
GPU & none {[}D5{]} & g4dn/g5/g6/g6e/g7e\ldots{} + inf2 \\
vCPU / RAM & \textbf{allocated per session; not a customer-set knob} ---
metered by vCPU-hour + GB-hour {[}D2, D13{]} & the chosen EC2 type's
spec (e.g.~\texttt{g6.xlarge} = 4 vCPU / 16 GB RAM / 1× L4) \\
Disk & \textbf{ephemeral only --- no persistent volume}; in-VM state
lost at session end {[}D3{]} & ephemeral \textbf{plus a sized,
persistent EBS volume} (survives stop/resume) \\
Container image & \textbf{arm64}, built image \textbf{≤ \textasciitilde2
GB} (project-measured, §``Building a customized Runtime'') &
\textbf{amd64}, far larger images allowed \\
Max session & 8 h & 14 d \\
Isolation unit & 1 microVM / session & 1 instance / session (1 GPU) \\
Filesystem & ephemeral & persistent EBS across stop/resume \\
Agents / session & 1 & many (1:N) \\
\end{longtable}

\textbf{On the microVM's resource envelope.} Two limits bite in practice
and are worth stating plainly. First, \textbf{the filesystem is
ephemeral}: ``any data stored in memory or written to disk persists only
for the compute lifecycle'' {[}D3{]}, so there is \emph{no} persistent
disk on the serverless substrate --- a repo you \texttt{git\ clone} or a
wheel you install lives only for that session (up to 8 h) and is wiped
on termination; durable state needs \textbf{session storage / Memory},
or the Instances substrate's EBS volume. Second, the \textbf{built
container image must stay small} --- this project measured a
\textbf{\textasciitilde2 GB} ceiling (§``Building a customized
Runtime''), which is why the GPU image installs its multi-GB CUDA wheel
at \texttt{CMD} rather than baking it in. AWS does \textbf{not} publish
a fixed per-session vCPU/RAM \emph{number} for the microVM --- capacity
is provisioned per session and billed by \textbf{vCPU-hour and GB-hour}
(the \texttt{agent.runtime.vcpu.hours.used} /
\texttt{agent.runtime.memory.gb\_hours.used} usage metrics {[}D13{]}),
so treat CPU/RAM as \emph{metered, not sized}, and confirm the current
hard caps on the \textbf{Service Quotas} page for your account before
sizing a heavy checker. When a task genuinely needs a large, fixed
RAM/disk footprint (a big SWE repo image, a data-heavy checker), that is
a reason to move it to \textbf{Instances}, where the resources are
simply the EC2 instance type's and the EBS volume you attach (§3, §``Two
compute types'').

\hypertarget{how-agentcore-manages-the-lifecycle-runtime-session-microvm-and-gpu-instance-end-to-end}{%
\subsection{How AgentCore manages the lifecycle --- Runtime, session,
microVM, and GPU instance, end to
end}\label{how-agentcore-manages-the-lifecycle-runtime-session-microvm-and-gpu-instance-end-to-end}}

An RL workload touches \textbf{four} AgentCore objects, with
\textbf{nested lifetimes}. Keeping them straight is what makes cost and
capacity predictable; conflating them is what makes a run's instance
count surprising. The whole section below is one answer to a single
question --- \emph{from the moment you deploy to the moment the run
ends, what is created, how long does it live, and what ends it?}

\begin{longtable}[]{@{}
  >{\raggedright\arraybackslash}p{(\columnwidth - 8\tabcolsep) * \real{0.2000}}
  >{\raggedright\arraybackslash}p{(\columnwidth - 8\tabcolsep) * \real{0.2000}}
  >{\raggedright\arraybackslash}p{(\columnwidth - 8\tabcolsep) * \real{0.2000}}
  >{\raggedright\arraybackslash}p{(\columnwidth - 8\tabcolsep) * \real{0.2000}}
  >{\raggedright\arraybackslash}p{(\columnwidth - 8\tabcolsep) * \real{0.2000}}@{}}
\toprule\noalign{}
\begin{minipage}[b]{\linewidth}\raggedright
Object
\end{minipage} & \begin{minipage}[b]{\linewidth}\raggedright
Created by
\end{minipage} & \begin{minipage}[b]{\linewidth}\raggedright
Lives for
\end{minipage} & \begin{minipage}[b]{\linewidth}\raggedright
Ended by
\end{minipage} & \begin{minipage}[b]{\linewidth}\raggedright
How many in a run
\end{minipage} \\
\midrule\noalign{}
\endhead
\bottomrule\noalign{}
\endlastfoot
\textbf{Runtime} (the deployable + endpoint) &
\texttt{CreateAgentRuntime}, at \textbf{deploy time} & the \textbf{whole
run} (and beyond) & \texttt{DeleteAgentRuntime} & \textbf{K} --- 1 by
default; a small pool only to raise invoke throughput \\
\textbf{Session} (\texttt{runtimeSessionId}) & first invoke with a
\textbf{fresh id} & \textbf{one episode} & \texttt{StopRuntimeSession} /
idle (15 min) / max-life & \textbf{≈ CONC} live at once,
\textbf{recycled} through every row \\
\textbf{microVM} (serverless / CPU substrate) & provisioned on session
launch & its \textbf{session} & session end → \textbf{memory sanitized}
& 1 per live CPU session \\
\textbf{GPU instance} (EC2, on Instances) & \textbf{cold-launched} on
session launch & its session (the instance may stay \textbf{warm} in the
capacity provider) & session end / capacity-provider scale-in & 1
instance = \textbf{one GPU}, ≈ CONC live \\
\end{longtable}

\textbf{The Runtime is a deploy-time artifact --- it spans the whole
run.} \texttt{CreateAgentRuntime} builds version V1 and the DEFAULT
endpoint; updates publish new versions; you route traffic by endpoint.
Deleting the Runtime removes the deployable --- \textbf{sessions are not
part of the Runtime's own state.} slime is simply handed the ARN and
never creates or deletes it (§``Who controls the AgentCore lifecycle'').
A run normally has \textbf{one} Runtime; pooling \textbf{K} of them
raises invoke throughput but changes neither how many tasks nor how many
images you can serve (§``A mixed-family batch'').

\textbf{The session is the RL unit --- one per episode, Active → Idle →
Terminated} {[}D3, D4{]}. The generate hook mints a fresh
\texttt{runtimeSessionId} per episode --- \emph{``one session id == one
isolated sandbox for the whole episode''} --- so the session id is the
isolation unit.

\begin{itemize}
\tightlist
\item
  \textbf{Launch / cold start.} There is no ``create session'' API. The
  first \texttt{invoke\_agent\_runtime} under a fresh id (the hook's
  \texttt{\_\_aenter\_\_} sends a \texttt{ping}) provisions the compute:
  a microVM on the serverless substrate, or a \textbf{cold-launched EC2
  instance} on Instances (\emph{minutes} for a GPU box). Same-id invokes
  thereafter hit warm compute.
\item
  \textbf{Active / Idle.} Each op (\texttt{write\_file} / \texttt{exec}
  / \texttt{read\_file}) is one more invoke under the \emph{same} id and
  resets the \textbf{15-minute idle timer}; between invokes the session
  sits Idle but warm.
\item
  \textbf{Termination.} Explicit via
  \textbf{\texttt{StopRuntimeSession}} {[}D12{]} (this project now calls
  it at episode end --- see below), or automatic on \textbf{idle expiry
  (15 min)}, \textbf{max lifetime} (microVM \textbf{8 h} / Instances
  \textbf{14 d}), or a failed health check → the microVM is
  \textbf{terminated and its memory sanitized}.
\end{itemize}

\textbf{The microVM and the GPU instance are the compute a session rides
on.} One live session == one microVM (CPU) or \textbf{one EC2 instance =
one GPU} (Instances), a strict 1:1 with \textbf{no cross-session GPU
time-slicing} (§``Two compute types''). The microVM is
\textbf{ephemeral} --- its filesystem and memory vanish at session end.
A GPU instance's \textbf{EBS volume persists across a stop}, and the
capacity provider may keep the \emph{instance} warm for the next session
even after the \emph{session} ends. So the session's lifetime and the
instance's lifetime are related but not identical: the session is what
your episode owns; the instance is what the fleet manages.

\textbf{Terminate at episode end --- this project now stops sessions
explicitly (per backend).} A session can be ended explicitly or
abandoned to the idle timer. The three sandbox backends differ at
episode end (\texttt{\_\_aexit\_\_}):

\begin{longtable}[]{@{}
  >{\raggedright\arraybackslash}p{(\columnwidth - 4\tabcolsep) * \real{0.3333}}
  >{\raggedright\arraybackslash}p{(\columnwidth - 4\tabcolsep) * \real{0.3333}}
  >{\raggedright\arraybackslash}p{(\columnwidth - 4\tabcolsep) * \real{0.3333}}@{}}
\toprule\noalign{}
\begin{minipage}[b]{\linewidth}\raggedright
Backend (\texttt{SANDBOX=})
\end{minipage} & \begin{minipage}[b]{\linewidth}\raggedright
At episode end
\end{minipage} & \begin{minipage}[b]{\linewidth}\raggedright
Session ends by
\end{minipage} \\
\midrule\noalign{}
\endhead
\bottomrule\noalign{}
\endlastfoot
\texttt{agentcore} (Code Interpreter microVM) & calls the client
\texttt{stop()} & \textbf{explicit}, immediately \\
\texttt{agentcore-runtime} (sandbox-as-Runtime) & calls
\textbf{\texttt{StopRuntimeSession}} (best-effort;
\texttt{AGENTCORE\_STOP\_SESSION=0} to opt out) & \textbf{explicit},
immediately --- idle reclaim is only the fallback if the stop call
fails \\
\texttt{agentcore-ray-farm} (B4; no shipped demo since 2026-09-30) &
sends \texttt{op=close} for the sub\_session & frees the
\textbf{fractional-GPU actor} immediately; the shared g6 session
persists \\
\end{longtable}

Because each episode mints a \textbf{fresh} \texttt{runtimeSessionId},
there is no warm session to preserve across episodes: stopping at
\texttt{\_\_aexit\_\_} returns the microVM (and its concurrent-session
quota slot) \textbf{immediately} instead of holding it idle for up to 15
minutes --- which matters when a batch runs many episodes concurrently
against a bounded session quota. Set \texttt{AGENTCORE\_STOP\_SESSION=0}
for a GPU fleet that deliberately keeps instances warm for the next
session, letting idle reclaim apply instead.

\textbf{A batch of prompts --- row count ≠ instance count.} A 100-row
dataset (KernelBench level\_1) does \textbf{not} launch 100 instances.
Rows stream through episodes across training steps; at any instant only
\texttt{-\/-sglang-server-concurrency} (default \textbf{4}) episodes are
in flight, so the live-instance count = concurrent sessions ≈ CONC, and
those few instances \textbf{recycle through all rows} over time. What
multiplies instances is \emph{concurrency}, not data:

\begin{itemize}
\tightlist
\item
  \textbf{One session per episode (\texttt{agentcore-runtime} on an
  Instance, e.g.~a g6 or an m5):} one session per episode × 1:1 ⇒
  \textbf{N concurrent episodes = N concurrent instances} (bounded by
  \texttt{HARNESS\_WORKERS} + \texttt{PREOPEN} and the capacity
  provider's max-instance cap).
\item
  \textbf{Historical --- B4 (\texttt{agentcore-ray-farm},
  \texttt{demo\_gpu\_farm}, removed 2026-09-30):} all concurrent
  episodes shared \textbf{one} \texttt{runtimeSessionId} → \textbf{one}
  g6 instance; each episode's sub\_session leased a fractional-GPU Ray
  actor, so N concurrent episodes still = one instance.
\end{itemize}

This is the concrete cost lever: packing episodes onto one instance (B4)
raises sandbox concurrency on a GPU task \textbf{without} paying for N
instances, but gives up the VM boundary between episodes; the project
now pays for the instances (one session per episode) and hides the cold
boot with pre-open.

\textbf{A mixed-task batch --- same session budget, different task per
row.} When a batch spans task families (a comma-list \texttt{DATASET=}),
the lifecycle above is \textbf{unchanged}: one Runtime image serves
every family, and the harness writes each row's task files into its
session and picks the \emph{task}, not the image (§``A mixed-family
batch''). A mixed-family batch therefore consumes exactly the same
Runtime / session / microVM / instance budget as a single-family batch
of the same size --- the family mix changes what runs \emph{inside} each
session, not how many sessions exist. Capability routing (a different
image per family) changes that by one Runtime per extra image --- the
Terminal-Bench demo does exactly this at task scale (one Runtime per
task, strict routing).

\textbf{Latency --- measure, don't assume.} The docs promise ``fast cold
starts'' but publish no number {[}D2{]}; a GPU first-invoke can take
\textbf{minutes} to launch the EC2. This project therefore instruments
every sandbox round-trip (\texttt{{[}sbx-op{]}} →
\texttt{perf/sandbox\_ms\_\{p50,p95\}}) and every generation
(\texttt{{[}sglang-gen{]}} → \texttt{perf/sglang\_gen\_ms\_*}) --- the
same signals the scaling analysis in §``Scaling the loop'' depends on.

\hypertarget{building-a-customized-runtime-bring-your-own-container-end-to-end}{%
\subsection{Building a customized Runtime --- bring your own container,
end to
end}\label{building-a-customized-runtime-bring-your-own-container-end-to-end}}

AWS offers three ways to put code on a Runtime --- the
\texttt{agentcore} CLI's zero-Docker \textbf{CodeZip} build {[}D16{]},
the \textbf{\texttt{bedrock-agentcore}} Python SDK's
\texttt{@app.entrypoint} {[}D15{]}, and \textbf{bring your own
container}. This project uses the \textbf{bring-your-own-container}
path, because it is the only one that gives you the whole image: base
OS, system packages, CUDA toolchain, a cloned repo --- exactly what a
task-specific sandbox needs. \texttt{CreateAgentRuntime} takes a
\textbf{union} --- \texttt{containerConfiguration} (an ECR image)
\textbf{or} \texttt{codeConfiguration} (a zip); BYO container is the
former {[}D9{]}. This section walks the full path from a Dockerfile to a
training run that uses it.

\textbf{Step 1 --- write a container that satisfies the contract.}
AgentCore only requires that your image serve HTTP on
\texttt{0.0.0.0:8080} with two routes --- \textbf{\texttt{/invocations}}
(one request = one op) and \textbf{\texttt{/ping}} (health) {[}D17{]}
--- and, for the managed microVM fleet, that it be built for
\textbf{arm64}. The \texttt{bedrock-agentcore} SDK's
\texttt{BedrockAgentCoreApp} + \texttt{@app.entrypoint} implement both
routes, so the image stays tiny. A minimal generic sandbox image (build
context: the demo folder, which vendors
\texttt{agentcore\_slime\_demo/}):

\begin{Shaded}
\begin{Highlighting}[]
\KeywordTok{FROM}\NormalTok{ public.ecr.aws/docker/library/python:3.12{-}slim}
\KeywordTok{WORKDIR}\NormalTok{ /app}
\KeywordTok{RUN} \ExtensionTok{pip}\NormalTok{ install }\AttributeTok{{-}{-}no{-}cache{-}dir}\NormalTok{ bedrock{-}agentcore}
\CommentTok{\# sandbox\_server\_agent.py is self{-}contained (stdlib + bedrock\_agentcore, no relative}
\CommentTok{\# imports), so copy it flat — no package, nothing heavy.}
\KeywordTok{COPY}\NormalTok{ agentcore\_slime\_demo/sandbox\_server\_agent.py /app/sandbox\_server\_agent.py}
\KeywordTok{ENV}\NormalTok{ SANDBOX\_WORKDIR=/workspace}
\KeywordTok{EXPOSE}\NormalTok{ 8080}
\KeywordTok{CMD}\NormalTok{ [}\StringTok{"python"}\NormalTok{, }\StringTok{"{-}m"}\NormalTok{, }\StringTok{"sandbox\_server\_agent"}\NormalTok{]      }\CommentTok{\# BedrockAgentCoreApp serves :8080}
\end{Highlighting}
\end{Shaded}

The Terminal-Bench demo does not write its own image: its deploy script
(\texttt{deploy\_terminalbench.py}) appends one wrapper layer to each
task's own Dockerfile, so the task's interpreter and files are
untouched:

\begin{Shaded}
\begin{Highlighting}[]
\KeywordTok{USER}\NormalTok{ root}
\KeywordTok{COPY}\NormalTok{ py/ /opt/acs/py/                       }\CommentTok{\# standalone arm64 CPython}
\KeywordTok{COPY}\NormalTok{ acs\_http\_server.py sandbox\_server\_agent.py /opt/acs/}
\KeywordTok{ENTRYPOINT}\NormalTok{ []}
\KeywordTok{CMD}\NormalTok{ [}\StringTok{"/opt/acs/py/bin/python3"}\NormalTok{, }\StringTok{"/opt/acs/acs\_http\_server.py"}\NormalTok{]   }\CommentTok{\# stdlib server on :8080}
\end{Highlighting}
\end{Shaded}

A \textbf{GPU} sandbox (the removed GPU demos' image) changes two
things: AgentCore injects \texttt{libcuda.so.1} (point
\texttt{LD\_LIBRARY\_PATH} at it), and the CUDA torch wheel is large ---
so install it \textbf{at \texttt{CMD}, not baked into the image}, to
stay under AgentCore's \textbf{\textasciitilde2 GB} image cap (§``Two
compute types''):

\begin{Shaded}
\begin{Highlighting}[]
\KeywordTok{ENV}\NormalTok{ LD\_LIBRARY\_PATH=/usr/lib64          }\CommentTok{\# AgentCore injects libcuda.so.1 here}
\KeywordTok{CMD}\NormalTok{ [}\StringTok{"sh"}\NormalTok{,}\StringTok{"{-}c"}\NormalTok{,}\StringTok{"pip install torch==2.5.1 {-}{-}index{-}url https://download.pytorch.org/whl/cu124 }\OperatorTok{\textbackslash{}}
\StringTok{                \&\& exec python {-}m sandbox\_server\_agent"}\NormalTok{]}
\end{Highlighting}
\end{Shaded}

\textbf{Step 2 --- create the ECR repo and log Docker in.}

\begin{Shaded}
\begin{Highlighting}[]
\BuiltInTok{export} \VariableTok{AWS\_REGION}\OperatorTok{=}\NormalTok{us{-}west{-}2}
\VariableTok{ACCT}\OperatorTok{=}\VariableTok{$(}\ExtensionTok{aws}\NormalTok{ sts get{-}caller{-}identity }\AttributeTok{{-}{-}query}\NormalTok{ Account }\AttributeTok{{-}{-}output}\NormalTok{ text}\VariableTok{)}
\VariableTok{REG}\OperatorTok{=}\VariableTok{$ACCT}\NormalTok{.dkr.ecr.}\VariableTok{$AWS\_REGION}\NormalTok{.amazonaws.com}
\ExtensionTok{aws}\NormalTok{ ecr create{-}repository }\AttributeTok{{-}{-}repository{-}name}\NormalTok{ acs{-}sandbox{-}cpu }\AttributeTok{{-}{-}region} \VariableTok{$AWS\_REGION}  \CommentTok{\# idempotent}
\ExtensionTok{aws}\NormalTok{ ecr get{-}login{-}password }\AttributeTok{{-}{-}region} \VariableTok{$AWS\_REGION} \KeywordTok{|} \ExtensionTok{docker}\NormalTok{ login }\AttributeTok{{-}{-}username}\NormalTok{ AWS }\AttributeTok{{-}{-}password{-}stdin} \VariableTok{$REG}
\end{Highlighting}
\end{Shaded}

\textbf{Step 3 --- cross-build for arm64 and push straight to ECR.} The
managed microVM fleet is \textbf{arm64} (\texttt{CreateAgentRuntime}
rejects amd64 with \emph{``Supported platforms: {[}arm64{]}''}), while a
GPU capacity provider is \textbf{amd64} --- so a build host of the other
arch must cross-build. Use a \texttt{docker-container} buildx builder
(with QEMU) and emit a \textbf{plain} single-arch image ---
\texttt{-\/-provenance=false\ -\/-sbom=false} keeps it out of the
attestation-manifest-list shape that \texttt{containerUri} rejects:

\begin{Shaded}
\begin{Highlighting}[]
\ExtensionTok{docker}\NormalTok{ buildx create }\AttributeTok{{-}{-}name}\NormalTok{ xbuild }\AttributeTok{{-}{-}driver}\NormalTok{ docker{-}container }\AttributeTok{{-}{-}bootstrap}   \CommentTok{\# once}
\ExtensionTok{docker}\NormalTok{ buildx build }\AttributeTok{{-}{-}builder}\NormalTok{ xbuild }\AttributeTok{{-}{-}platform}\NormalTok{ linux/arm64 }\DataTypeTok{\textbackslash{}}
  \AttributeTok{{-}{-}provenance}\OperatorTok{=}\NormalTok{false }\AttributeTok{{-}{-}sbom}\OperatorTok{=}\NormalTok{false }\DataTypeTok{\textbackslash{}}
  \AttributeTok{{-}f}\NormalTok{ Dockerfile.sandbox }\AttributeTok{{-}t} \VariableTok{$REG}\NormalTok{/acs{-}sandbox{-}cpu:latest }\AttributeTok{{-}{-}push}\NormalTok{ .}
\end{Highlighting}
\end{Shaded}

\textbf{Step 4 --- register the image as a Runtime.} An
\textbf{execution role} (trusts
\texttt{bedrock-agentcore.amazonaws.com}; grants ECR pull + CloudWatch
logs; for a harness Runtime that calls the sandbox, also
\texttt{InvokeAgentRuntime}) is passed to
\texttt{create\_agent\_runtime}, whose \texttt{agentRuntimeArtifact}
carries the \texttt{containerUri}:

\begin{Shaded}
\begin{Highlighting}[]
\NormalTok{control }\OperatorTok{=}\NormalTok{ boto3.client(}\StringTok{"bedrock{-}agentcore{-}control"}\NormalTok{, region\_name}\OperatorTok{=}\NormalTok{REGION)}
\NormalTok{resp }\OperatorTok{=}\NormalTok{ control.create\_agent\_runtime(}
\NormalTok{    agentRuntimeName}\OperatorTok{=}\StringTok{"acs{-}sandbox{-}cpu"}\NormalTok{,}
\NormalTok{    roleArn}\OperatorTok{=}\NormalTok{exec\_role\_arn,}
\NormalTok{    agentRuntimeArtifact}\OperatorTok{=}\NormalTok{\{}\StringTok{"containerConfiguration"}\NormalTok{: \{}\StringTok{"containerUri"}\NormalTok{: }\SpecialStringTok{f"}\SpecialCharTok{\{}\NormalTok{REG}\SpecialCharTok{\}}\SpecialStringTok{/acs{-}sandbox{-}cpu:latest"}\NormalTok{\}\},}
\NormalTok{    protocolConfiguration}\OperatorTok{=}\NormalTok{\{}\StringTok{"serverProtocol"}\NormalTok{: }\StringTok{"HTTP"}\NormalTok{\},}
\NormalTok{    networkConfiguration}\OperatorTok{=}\NormalTok{\{}\StringTok{"networkMode"}\NormalTok{: }\StringTok{"PUBLIC"}\NormalTok{\},   }\CommentTok{\# or VPC + networkModeConfig}
\NormalTok{    environmentVariables}\OperatorTok{=}\NormalTok{\{}\StringTok{"SANDBOX\_WORKDIR"}\NormalTok{: }\StringTok{"/workspace"}\NormalTok{\},}
\NormalTok{)}
\CommentTok{\# poll DescribeAgentRuntime until status == READY, then keep the ARN:}
\NormalTok{sandbox\_arn }\OperatorTok{=}\NormalTok{ resp[}\StringTok{"agentRuntimeArn"}\NormalTok{]}
\end{Highlighting}
\end{Shaded}

For a \textbf{GPU} sandbox this step is preceded by
\textbf{\texttt{CreateCapacityProvider}} (\texttt{ec2Configuration}:
instance type e.g.~\texttt{g6.xlarge}, VPC subnets/SGs, EBS,
\texttt{idleInstanceTimeout=600}, \texttt{maxLifetime=86400}), attached
via \texttt{capacityProviderConfiguration} {[}D5, D9{]}; the execution
setup then also needs the operator policy
\texttt{BedrockAgentCoreRuntimeInstancesOperatorRolePolicy}.

The Terminal-Bench demo bundles these steps into its numbered scripts.
Both are idempotent and resumable: \texttt{bash\ 2\_build\_image.sh}
does the ECR create, buildx and push per task, and
\texttt{bash\ 4\_deploy\_runtime.sh} calls
\texttt{terminalbench/deploy\_terminalbench.py}, which uses
\texttt{deploy/agentcore\_deploy.py} to mint the roles, register one VPC
Runtime per image and write the image → ARN map:

\begin{Shaded}
\begin{Highlighting}[]
\BuiltInTok{cd}\NormalTok{ code/demos/demo\_terminalbench\_kernelbench}
\FunctionTok{bash}\NormalTok{ 2\_build\_image.sh      }\CommentTok{\# ECR create → buildx {-}{-}platform linux/arm64 → push, per task}
\FunctionTok{bash}\NormalTok{ 4\_deploy\_runtime.sh   }\CommentTok{\# exec role → CreateAgentRuntime (VPC) per image → arn\_by\_image.json}
\end{Highlighting}
\end{Shaded}

\textbf{Step 5 --- use it in a training run.} Nothing in slime
references the image or the ARN at build time; the sandbox client reads
the ARN from the environment at rollout time (and raises if unset). So
``using'' a custom Runtime in training is one export --- the ARN that
Step 4 printed --- before launching \texttt{run.sh} (a pool like
Terminal-Bench's exports the map
\texttt{AGENTCORE\_SANDBOX\_RUNTIME\_ARN\_BY\_IMAGE} instead):

\begin{Shaded}
\begin{Highlighting}[]
\BuiltInTok{export} \VariableTok{AGENTCORE\_SANDBOX\_RUNTIME\_ARN}\OperatorTok{=}\NormalTok{arn:aws:bedrock{-}agentcore:us{-}west{-}2:}\OperatorTok{\textless{}}\NormalTok{acct}\OperatorTok{\textgreater{}}\NormalTok{:runtime/acs{-}sandbox{-}cpu{-}…}
\BuiltInTok{export} \VariableTok{AGENTCORE\_REGION}\OperatorTok{=}\NormalTok{us{-}west{-}2}
\BuiltInTok{export} \VariableTok{SANDBOX}\OperatorTok{=}\NormalTok{agentcore{-}runtime          }\CommentTok{\# select the sandbox{-}as{-}Runtime backend}
\FunctionTok{bash}\NormalTok{ run.sh                               }\CommentTok{\# slime\textquotesingle{}s generate hook now invokes YOUR Runtime per op}
\end{Highlighting}
\end{Shaded}

From here the lifecycle in §``How AgentCore manages the lifecycle''
takes over: the first \texttt{invoke\_agent\_runtime} under each
episode's fresh session id cold-starts a microVM (or GPU instance) from
\textbf{your} image, every tool call is one more invoke, and
\texttt{\_\_aexit\_\_} stops the session at episode end.

\hypertarget{using-agentcore-as-the-sandbox-in-an-agentic-rl-workflow}{%
\section{Using AgentCore as the sandbox in an agentic RL
workflow}\label{using-agentcore-as-the-sandbox-in-an-agentic-rl-workflow}}

With the foundations in place, this section shows how AgentCore Runtime
plugs into a real RL trainer --- \textbf{slime} --- as the execution
\textbf{sandbox}, with the harness running \textbf{inside the sandbox
session}. The integration is deliberately small: \textbf{no fork of
slime}, two hook flags, and a sandbox client that turns each tool call
into one \texttt{InvokeAgentRuntime}.

\hypertarget{where-agentcore-attaches-to-slime-hooks-not-a-fork}{%
\subsection{Where AgentCore attaches to slime --- hooks, not a
fork}\label{where-agentcore-attaches-to-slime-hooks-not-a-fork}}

slime {[}9{]} runs Megatron (training) and SGLang (rollout)
asynchronously (§``slime --- the trainer''). It exposes the rollout path
as a \textbf{custom generation function}, so an agentic workload plugs
in without touching the training kernel:

\begin{itemize}
\tightlist
\item
  \texttt{-\/-custom-generate-function-path\ agentcore\_slime\_demo.generate.generate}
  --- the per-prompt rollout: it hands the episode to the harness tier
  (\texttt{harness\_pool.py}), which takes an AgentCore sandbox session,
  then an SGLang engine, and runs the mini\_claude loop inside that
  session.
\item
  \texttt{-\/-custom-rm-path} --- an optional custom reward hook (here
  the reward comes straight from the in-sandbox checker, so the generate
  hook attaches it directly).
\end{itemize}

The key design choice: \textbf{AgentCore is the sandbox, and the harness
rides in it.} The harness (agent loop) is independent per episode, so it
runs inside the episode's own sandbox session and calls the trainer's
per-engine \texttt{AnthropicAdapter} back over the VPC; the trainer
keeps only the dispatch tier and the adapters, which record the exact
tokens. There is no in-trainer fallback: a session that cannot host the
loop (no \texttt{python3} in the image, or no route back to the adapter)
fails the episode with \texttt{HarnessPlacementError}. The legacy
two-runtime path (a separate harness Runtime,
\texttt{agentcore\_runtime\_generate.py} /
\texttt{harness\_runtime\_agent.py}) was deleted.

\hypertarget{the-call-path-one-diagram}{%
\subsection{The call path --- one
diagram}\label{the-call-path-one-diagram}}

\begin{figure}[htbp]
\centering
\resizebox{0.95\textwidth}{!}{%
\begin{tikzpicture}[
  font=\small,
  box/.style={draw, rounded corners, minimum height=1.2cm, minimum width=3.0cm, align=center, thick, fill=white},
  gpu/.style={box, fill=blue!8},
  cpu/.style={box, fill=black!4},
  focus/.style={box, fill=orange!18},
  lbl/.style={font=\scriptsize, align=center, fill=white, inner sep=1.5pt},
  ptitle/.style={font=\scriptsize\bfseries, text=black!55},
  panel/.style={draw, rounded corners, dashed, thick, inner sep=0.34cm},
  >={Stealth[]},
]
% nodes
\node[cpu]   (slime)  at (0,0)     {\textbf{slime}\\[1pt]\scriptsize Data Buffer $\leftrightarrow$ Megatron\\\scriptsize + SGLang router};
\node[focus] (hook)   at (5.4,0)   {\textbf{generate() hook}\\[1pt]\scriptsize harness tier\\\scriptsize ($W$ dispatch workers)};
\node[cpu]   (proxy)  at (5.4,3.0) {\textbf{AnthropicAdapter}\\[1pt]\scriptsize proxy $\rightarrow$ SGLang};
\node[focus] (sbx)    at (12.6,0)  {\textbf{AgentCore Runtime}\\[1pt]\scriptsize sandbox microVM\\\scriptsize mini\_claude loop + tools\\\scriptsize + hidden checker};

% slime <-> hook
\draw[->, thick] ([yshift=5pt]slime.east) -- node[lbl,above]{prompt} ([yshift=5pt]hook.west);
\draw[->, thick] ([yshift=-5pt]hook.west) -- node[lbl,below]{Sample + reward} ([yshift=-5pt]slime.east);
% sandbox (in-session loop) <-> proxy: the loop dials the adapter back
\draw[->, thick] ([xshift=-10pt]sbx.north) -- node[lbl,sloped,above,pos=0.5]{model turn (dial-back)} ([yshift=-6pt]proxy.east);
\draw[->, thick] ([yshift=6pt]proxy.east) -- node[lbl,sloped,above,pos=0.5]{response (tokens recorded)} ([xshift=10pt]sbx.north);
% hook <-> sandbox
\draw[->, thick] ([yshift=5pt]hook.east) -- node[lbl,above,pos=0.6]{invoke: seed, start, long-poll} ([yshift=5pt]sbx.west);
\draw[->, thick] ([yshift=-5pt]sbx.west) -- node[lbl,below,pos=0.6]{status + checker result} ([yshift=-5pt]hook.east);
% proxy -> hook: the recorded tokens become the Sample
\draw[->, thick, dashed] (proxy.south) -- node[lbl,left,pos=0.5]{tokens + logprobs} (hook.north);

% panels
\begin{scope}[on background layer]
  \node[panel, fill=blue!3,   fit=(slime)(proxy)(hook)] (phost) {};
  \node[panel, fill=orange!4, fit=(sbx)]                (psbx)  {};
\end{scope}
\node[ptitle, anchor=south west] at ([yshift=2pt]phost.north west) {Trainer host --- dispatch tier + adapters};
\node[ptitle, anchor=north] at ([yshift=-2pt]psbx.south) {AgentCore session = 1 episode (isolated; runs the agent loop)};
\end{tikzpicture}%
}
\caption{How AgentCore drops into slime. slime hands a \textbf{prompt} to the \texttt{generate()} hook, whose \textbf{harness tier} queues the episode: a free dispatch worker takes a fresh \textbf{AgentCore sandbox} session, then the least-loaded SGLang engine, seeds the task, and starts the \textbf{mini\_claude agent loop inside the session}. Each \textbf{model turn} dials back from the session to that engine's \textbf{proxy} (\texttt{AnthropicAdapter}), which records the exact sampled \textbf{tokens + logprobs} --- the TITO property of §\"Token-in-token-out\"; each \textbf{tool call} runs locally in the session. The trainer only long-polls the loop over \texttt{invoke\_agent\_runtime}, then runs the hidden \textbf{checker} (the reward source). The adapter's record becomes the loss-masked \textbf{Sample + reward} in slime's buffer. The trainer host keeps only the dispatch tier and the adapters; the loop and its tools live in the isolated per-episode microVM.}
\end{figure}

\hypertarget{key-code-three-pieces}{%
\subsection{Key code --- three pieces}\label{key-code-three-pieces}}

\textbf{1. The generate hook (the harness tier).} slime calls this once
per prompt. It queues the episode on the harness tier: a free dispatch
worker takes \textbf{one fresh sandbox session}, then the SGLang engine
with the fewest in-flight episodes (sticky for all its turns, §``Scaling
the loop''), runs the mini\_claude loop inside that session, and returns
a single loss-masked \texttt{Sample}:

\begin{Shaded}
\begin{Highlighting}[]
\CommentTok{\# agentcore\_slime\_demo/generate.py  (wired via {-}{-}custom{-}generate{-}function{-}path)}
\ControlFlowTok{async} \KeywordTok{def}\NormalTok{ generate(args, sample, sampling\_params, evaluation}\OperatorTok{=}\VariableTok{False}\NormalTok{):}
\NormalTok{    task }\OperatorTok{=}\NormalTok{ \_task\_from\_sample(sample)                        }\CommentTok{\# rebuild Task from row metadata}
\NormalTok{    kind, image }\OperatorTok{=}\NormalTok{ os.environ.get(}\StringTok{"SANDBOX"}\NormalTok{, }\StringTok{"agentcore{-}runtime"}\NormalTok{), task.image}
\NormalTok{    pool }\OperatorTok{=}\NormalTok{ harness\_pool.get\_pool(args)                      }\CommentTok{\# the harness tier (W workers)}

    \ControlFlowTok{async} \KeywordTok{def}\NormalTok{ episode(worker):                              }\CommentTok{\# run by a FREE worker (pull queue)}
        \ControlFlowTok{async} \ControlFlowTok{with}\NormalTok{ pool.sandbox((kind, image), }\KeywordTok{lambda}\NormalTok{: maybe\_keep\_warm(}
\NormalTok{                make\_sandbox(kind, image}\OperatorTok{=}\NormalTok{image))) }\ImportTok{as}\NormalTok{ (sb, wait\_ms, slot):  }\CommentTok{\# sandbox FIRST}
            \ControlFlowTok{with}\NormalTok{ pool.engine() }\ImportTok{as}\NormalTok{ (e, engine):              }\CommentTok{\# then the least{-}loaded engine}
                \CommentTok{\# mini\_claude runs INSIDE sb and dials this engine\textquotesingle{}s adapter back; the}
                \CommentTok{\# adapter\textquotesingle{}s token record becomes ONE loss{-}masked Sample + the reward}
                \ControlFlowTok{return} \ControlFlowTok{await}\NormalTok{ mini\_claude\_episode.run(args, sample, task, sampling\_params,}
\NormalTok{                                                     sb, engine, fields}\OperatorTok{=}\NormalTok{pool.tier\_fields(}
\NormalTok{                                                         worker, e, wait\_ms, slot))}

    \ControlFlowTok{return} \ControlFlowTok{await}\NormalTok{ pool.run(episode)                          }\CommentTok{\# queue it; await the bare Sample}
\end{Highlighting}
\end{Shaded}

Two RL-critical details: the loss mask is \textbf{1 on the policy's own
tokens, 0 on prompt/observation/tool tokens}, and the adapter's
per-session record (\texttt{finish\_session}) reuses the
\textbf{recorded \texttt{token\_ids} per turn} --- no re-tokenization,
the TITO property of §``Token-in-token-out''. The pool returns a
\emph{bare} \texttt{Sample} (not a list): slime's fully-async path (the
only rollout driver) chokes on nested sample lists.

\textbf{2. The sandbox client (one op = one invoke).} Each
\texttt{write\_file} / \texttt{exec} / \texttt{read\_file} becomes a
single SigV4 \texttt{invoke\_agent\_runtime}, keyed by the episode's
session id. boto3 is synchronous, so it is offloaded with
\texttt{asyncio.to\_thread}, letting N episodes run concurrently (the
harness tier sizes the loop's default executor to the pool, so the
thread pool is not a hidden cap):

\begin{Shaded}
\begin{Highlighting}[]
\CommentTok{\# agentcore\_slime\_demo/agentcore\_runtime\_sandbox.py}
\KeywordTok{def}\NormalTok{ \_invoke\_sync(}\VariableTok{self}\NormalTok{, payload: }\BuiltInTok{dict}\NormalTok{) }\OperatorTok{{-}\textgreater{}} \BuiltInTok{dict}\NormalTok{:}
\NormalTok{    resp }\OperatorTok{=} \VariableTok{self}\NormalTok{.\_client.invoke\_agent\_runtime(               }\CommentTok{\# boto3 "bedrock{-}agentcore"}
\NormalTok{        agentRuntimeArn}\OperatorTok{=}\VariableTok{self}\NormalTok{.\_arn,}
\NormalTok{        runtimeSessionId}\OperatorTok{=}\VariableTok{self}\NormalTok{.\_session\_id,                  }\CommentTok{\# \textgreater{}=33 chars; one per episode}
\NormalTok{        payload}\OperatorTok{=}\NormalTok{json.dumps(payload).encode())}
    \ControlFlowTok{return}\NormalTok{ json.loads(resp[}\StringTok{"response"}\NormalTok{].read().decode())}

\ControlFlowTok{async} \KeywordTok{def} \BuiltInTok{exec}\NormalTok{(}\VariableTok{self}\NormalTok{, cmd, }\OperatorTok{*}\NormalTok{, timeout}\OperatorTok{=}\DecValTok{120}\NormalTok{, }\OperatorTok{**}\NormalTok{kw):}
\NormalTok{    res }\OperatorTok{=} \ControlFlowTok{await}\NormalTok{ asyncio.to\_thread(}\VariableTok{self}\NormalTok{.\_invoke\_sync,}
\NormalTok{              \{}\StringTok{"op"}\NormalTok{: }\StringTok{"exec"}\NormalTok{, }\StringTok{"command"}\NormalTok{: cmd, }\StringTok{"timeout"}\NormalTok{: }\BuiltInTok{int}\NormalTok{(timeout), }\StringTok{"workdir"}\NormalTok{: }\VariableTok{self}\NormalTok{.\_workdir\})}
    \ControlFlowTok{return} \BuiltInTok{int}\NormalTok{(res.get(}\StringTok{"exit\_code"}\NormalTok{, }\DecValTok{1}\NormalTok{)), res.get(}\StringTok{"stdout"}\NormalTok{, }\StringTok{""}\NormalTok{), res.get(}\StringTok{"stderr"}\NormalTok{, }\StringTok{""}\NormalTok{)}
\end{Highlighting}
\end{Shaded}

The real client also \textbf{retries} transient 5xx/throttle blips with
backoff (one bad round-trip must not crash a multi-hour run) and carries
a much larger budget for GPU ``insufficient capacity'' while a g6
launches.

\textbf{3. The Runtime server (what runs inside the sandbox).} On the
Runtime, the \texttt{bedrock-agentcore} SDK's \texttt{@app.entrypoint}
receives each payload and dispatches on \texttt{op}. This is the whole
sandbox server --- and \emph{the checker runs here, so the microVM is
the reward source}:

\begin{Shaded}
\begin{Highlighting}[]
\CommentTok{\# agentcore\_slime\_demo/sandbox\_server\_agent.py  (deployed as the Runtime image)}
\ImportTok{from}\NormalTok{ bedrock\_agentcore.runtime }\ImportTok{import}\NormalTok{ BedrockAgentCoreApp}

\KeywordTok{def}\NormalTok{ handle(request: }\BuiltInTok{dict}\NormalTok{) }\OperatorTok{{-}\textgreater{}} \BuiltInTok{dict}\NormalTok{:              }\CommentTok{\# pure dispatch (unit{-}testable without AgentCore)}
\NormalTok{    op }\OperatorTok{=}\NormalTok{ (request }\KeywordTok{or}\NormalTok{ \{\}).get(}\StringTok{"op"}\NormalTok{, }\StringTok{"ping"}\NormalTok{)}
    \ControlFlowTok{if}\NormalTok{ op }\OperatorTok{==} \StringTok{"ping"}\NormalTok{:       \_ensure\_workdir(request.get(}\StringTok{"workdir"}\NormalTok{))}\OperatorTok{;} \ControlFlowTok{return}\NormalTok{ \{}\StringTok{"ok"}\NormalTok{: }\VariableTok{True}\NormalTok{\}}
    \ControlFlowTok{if}\NormalTok{ op }\OperatorTok{==} \StringTok{"write\_file"}\NormalTok{: }\ControlFlowTok{return}\NormalTok{ \_do\_write(request)     }\CommentTok{\# {-}\textgreater{} \{"ok": True\}}
    \ControlFlowTok{if}\NormalTok{ op }\OperatorTok{==} \StringTok{"read\_file"}\NormalTok{:  }\ControlFlowTok{return}\NormalTok{ \_do\_read(request)      }\CommentTok{\# {-}\textgreater{} \{"text": ...\}}
    \ControlFlowTok{if}\NormalTok{ op }\OperatorTok{==} \StringTok{"exec"}\NormalTok{:       }\ControlFlowTok{return}\NormalTok{ \_do\_exec(request)      }\CommentTok{\# {-}\textgreater{} \{exit\_code, stdout, stderr\}  (runs the checker)}

\NormalTok{app }\OperatorTok{=}\NormalTok{ BedrockAgentCoreApp()}

\AttributeTok{@app.entrypoint}
\KeywordTok{def}\NormalTok{ entrypoint(request):}
    \ControlFlowTok{return}\NormalTok{ handle(request)}

\NormalTok{app.run()      }\CommentTok{\# serves 0.0.0.0:8080 — POST /invocations + GET /ping}
\end{Highlighting}
\end{Shaded}

\hypertarget{who-controls-the-agentcore-lifecycle-deploy-driver-vs.-the-hook}{%
\subsection{Who controls the AgentCore lifecycle --- deploy driver
vs.~the
hook}\label{who-controls-the-agentcore-lifecycle-deploy-driver-vs.-the-hook}}

A precise point that the code above implies but is worth stating
outright: \textbf{slime's training kernel controls neither the Runtime
nor the session.} Lifecycle is split across two planes, and slime's core
(\texttt{RolloutManager}, Data Buffer, Megatron, the SGLang router) is
entirely \textbf{AgentCore-unaware} --- it only calls the function
registered at \texttt{-\/-custom-generate-function-path} and gets a
\texttt{Sample} back.

\begin{itemize}
\tightlist
\item
  \textbf{Runtime lifecycle (create/delete the ARN) --- outside slime,
  at deploy time.} No slime component calls
  \texttt{CreateAgentRuntime}/\texttt{DeleteAgentRuntime}. The Runtime
  (and, for GPU, the capacity provider) is provisioned \textbf{once
  before training} by the deploy driver
  (\texttt{terminalbench/deploy\_terminalbench.py}), and slime is simply
  handed the ARN via \texttt{AGENTCORE\_SANDBOX\_RUNTIME\_ARN}, or a
  per-image map via
  \texttt{AGENTCORE\_SANDBOX\_RUNTIME\_ARN\_BY\_IMAGE}. The sandbox
  client only \emph{reads} that env var and raises if it is unset --- it
  never creates the Runtime. The Runtime is a \textbf{deploy-time
  artifact}, not a training-loop object.
\item
  \textbf{Session lifecycle (launch/terminate per episode) --- the
  custom \texttt{generate} hook and its sandbox client.} All session
  control lives \emph{below} slime's seam, inside
  \texttt{agentcore\_slime\_demo.generate.generate}, which opens one
  sandbox per prompt
  (\texttt{async\ with\ make\_sandbox(...)\ as\ sb:}). The client
  (\texttt{AgentCoreRuntimeSandbox}) mints \textbf{one
  \texttt{runtimeSessionId} per episode}, and:

  \begin{itemize}
  \tightlist
  \item
    \textbf{Launch is implicit} --- there is no ``create session'' API.
    \texttt{\_\_aenter\_\_} sends the \textbf{first
    \texttt{invoke\_agent\_runtime}} (a \texttt{ping}) under the fresh
    session id, and that first invoke is what \textbf{cold-starts the
    microVM} (and, on GPU, warms the g6).
  \item
    \textbf{Each op} (\texttt{write\_file} / \texttt{exec} /
    \texttt{read\_file}) is one more \texttt{invoke\_agent\_runtime}
    under the \emph{same} session id.
  \item
    \textbf{Termination is explicit} --- \texttt{\_\_aexit\_\_} calls
    \textbf{\texttt{StopRuntimeSession}} (best-effort;
    \texttt{AGENTCORE\_STOP\_SESSION=0} opts out) so the microVM and its
    concurrent-session quota slot are released \textbf{immediately} at
    episode end. AgentCore's \textbf{\textasciitilde15-minute idle
    timeout} remains the fallback if the stop call fails. slime's kernel
    still never terminates a session itself --- the control lives in the
    hook's sandbox client.
  \end{itemize}
\end{itemize}

\begin{longtable}[]{@{}
  >{\raggedright\arraybackslash}p{(\columnwidth - 6\tabcolsep) * \real{0.2500}}
  >{\raggedright\arraybackslash}p{(\columnwidth - 6\tabcolsep) * \real{0.2500}}
  >{\raggedright\arraybackslash}p{(\columnwidth - 6\tabcolsep) * \real{0.2500}}
  >{\raggedright\arraybackslash}p{(\columnwidth - 6\tabcolsep) * \real{0.2500}}@{}}
\toprule\noalign{}
\begin{minipage}[b]{\linewidth}\raggedright
Action
\end{minipage} & \begin{minipage}[b]{\linewidth}\raggedright
API
\end{minipage} & \begin{minipage}[b]{\linewidth}\raggedright
Controlled by
\end{minipage} & \begin{minipage}[b]{\linewidth}\raggedright
When
\end{minipage} \\
\midrule\noalign{}
\endhead
\bottomrule\noalign{}
\endlastfoot
Create / delete Runtime (+ capacity provider) &
\texttt{CreateAgentRuntime} / \texttt{Delete…} /
\texttt{CreateCapacityProvider} & \textbf{deploy driver} --- \emph{not
slime} & once, before/after the run \\
Drive rollout & --- & slime \texttt{RolloutManager} / Data Buffer ---
\textbf{AgentCore-unaware} & per prompt \\
\textbf{Launch} session (cold-start microVM) & first
\texttt{invoke\_agent\_runtime} (\texttt{ping}) & the hook's sandbox
client \texttt{\_\_aenter\_\_} & per episode \\
Run ops & \texttt{invoke\_agent\_runtime} (write/exec/read) & sandbox
client, same session id & each turn \\
\textbf{Terminate} session & \textbf{\texttt{StopRuntimeSession}}
(best-effort; idle timeout is the fallback) & the hook's sandbox client
\texttt{\_\_aexit\_\_} & per episode, at end \\
Defer termination across weight-sync (optional) & repeated \texttt{ping}
& \texttt{SessionKeeper}, \emph{outside} the harness & during the sync
barrier \\
\end{longtable}

Explicit \texttt{StopRuntimeSession} at episode end is the happy-path
teardown, but a session can still be lost \textbf{mid-episode}: a long
weight-sync pause between two turns can trip the 15-minute idle clock
while the episode is still running and reclaim the session out from
under it --- exactly the failure mode the \texttt{HEARTBEAT=1}
\texttt{SessionKeeper} heartbeat addresses (§``Hiding the container
launch: pre-open and heartbeat''). \textbf{One-sentence version:} the
deploy driver owns the Runtime, the \texttt{generate} hook + its sandbox
client own the session (launch on first invoke, terminate with
\texttt{StopRuntimeSession} at episode end), and slime's training kernel
owns neither.

\hypertarget{ending-the-session-at-episode-end-explicit-teardown-and-what-end-means-per-substrate}{%
\subsection{Ending the session at episode end --- explicit teardown, and
what ``end'' means per
substrate}\label{ending-the-session-at-episode-end-explicit-teardown-and-what-end-means-per-substrate}}

\textbf{We always end the session explicitly when the episode ends ---
and \emph{what} gets ended differs by substrate, deliberately.} The
teardown is not inside the loop: the episode
(\texttt{mini\_claude\_episode.run}, and the in-session loop it starts)
runs \emph{inside an already-open sandbox} and closes nothing itself ---
it grades and returns the sample. The close fires because \textbf{every
caller wraps the episode in a context manager} --- the harness tier's
\texttt{async\ with\ pool.sandbox(...)} around
\texttt{make\_sandbox(...)} (\texttt{harness\_pool.py},
\texttt{generate.py}) --- so the sandbox's \texttt{\_\_aexit\_\_} runs
the instant the episode returns. (With the opt-in
\texttt{SANDBOX\_SESSIONS=M}, a pooled session instead has its
\texttt{ep-\textless{}hex\textgreater{}} workdir removed and goes back
to the free list; the pool stops every pooled session at exit.) One
teardown path, exercised identically by all demos.

What \texttt{\_\_aexit\_\_} does depends on the backend the
\texttt{SANDBOX} env selected:

\begin{itemize}
\tightlist
\item
  \textbf{Per-episode session (managed microVM or capacity-provider
  Instance, \texttt{agentcore-runtime}; every shipped demo): stop the
  whole session.} \texttt{AgentCoreRuntimeSandbox.\_\_aexit\_\_} calls
  \textbf{\texttt{StopRuntimeSession(runtimeSessionId,\ agentRuntimeArn)}}
  (\texttt{agentcore\_runtime\_sandbox.py:287-324}; opt-out
  \texttt{AGENTCORE\_STOP\_SESSION=0}). Each episode minted its own
  \texttt{sbx-\textless{}uuid\textgreater{}} session, so ending the
  whole thing is exactly right --- there is nothing to preserve.
\item
  \textbf{Warm-shared instance (CPU/GPU Instances,
  \texttt{agentcore-ray-farm}; historical, its last demo
  \texttt{demo\_gpu\_farm} was removed 2026-09-30): close only the
  episode's \texttt{sub\_session}, never the shared session.}
  \texttt{RayFarmRuntimeSandbox.\_\_aexit\_\_} sends a single
  \textbf{\texttt{op=close} for the per-episode \texttt{sub\_session}}
  and releases its client-side concurrency slot
  (\texttt{agentcore\_runtime\_sandbox.py:484-493}); it \textbf{never}
  calls \texttt{StopRuntimeSession}. Server-side, that wipes just that
  episode's working directory and returns the fractional-GPU Ray actor
  to the pool, while the warm CUDA context, the shared
  \texttt{runtimeSessionId}, and the instance stay alive for the next
  episode (\texttt{sandbox\_ray\_farm.py:378-385}, \texttt{:219-226}).
\item
  \textbf{Default / local:} \texttt{AgentCoreSandbox.\_\_aexit\_\_} →
  \texttt{client.stop()}; \texttt{LocalSandbox} → temp-workspace cleanup
  (\texttt{sandbox.py:194}, \texttt{:110}).
\end{itemize}

\textbf{Why we do it, and why it matters.} A session is not free to
leave dangling: it holds a microVM (or an Instance's actor slot) and
counts against AgentCore's \textbf{concurrent-session quota}, and the
only automatic reclaim is the \textbf{\textasciitilde15-minute idle
timeout} --- far too slow at rollout throughput, where leaked sessions
throttle concurrency and inflate cost. Explicit
\texttt{StopRuntimeSession} releases the slot \textbf{immediately}
instead of waiting out the idle clock, and closing the per-episode
workspace guarantees the next episode starts from a \textbf{clean,
isolated env} (no cross-episode file leakage). For the warm-shared path
the \emph{asymmetry is the whole point}: closing \textbf{only} the
\texttt{sub\_session} frees the episode's env and hands the next episode
a fresh workspace \textbf{without re-paying the \textasciitilde60--80 s
instance cold-boot} --- whereas stopping the shared session would defeat
warm reuse (re-pay the boot every episode) and could tear the instance
out from under sibling episodes still running on it. So the single
``always end at episode end'' discipline yields two correct behaviors:
\textbf{kill the per-episode session} where the session \emph{is} the
instance, and \textbf{kill only the sub-session} where one warm instance
is shared. (The complementary hazard --- a session lost
\emph{mid-episode} to the idle clock during a long weight-sync pause ---
is what the \texttt{HEARTBEAT=1} \texttt{SessionKeeper} addresses.)

\hypertarget{one-harness-many-task-families}{%
\subsection{One harness, many task
families}\label{one-harness-many-task-families}}

The episode (\texttt{mini\_claude\_episode.py}) is
\textbf{task-agnostic} --- the same seed → agent loop → hidden check →
reward cycle serves every family; only the \emph{prompt, seed files, and
checker} change (\texttt{task\_families.py}), dispatched by
\texttt{metadata.task.family}:

\begin{Shaded}
\begin{Highlighting}[]
\ControlFlowTok{for}\NormalTok{ name, content }\KeywordTok{in}\NormalTok{ task.setup\_files():                   }\CommentTok{\# seed the workspace}
    \ControlFlowTok{await}\NormalTok{ sb.write\_file(name, content)}
\NormalTok{res }\OperatorTok{=} \ControlFlowTok{await}\NormalTok{ harness.run(sb, session\_id}\OperatorTok{=}\NormalTok{sid,                }\CommentTok{\# mini\_claude loop INSIDE the session:}
\NormalTok{                        adapter\_url}\OperatorTok{=}\NormalTok{session\_url,           }\CommentTok{\#   write\_file / bash / read\_file tools,}
\NormalTok{                        prompt}\OperatorTok{=}\NormalTok{task.initial\_user\_prompt()) }\CommentTok{\#   model turns dial the adapter back}
\NormalTok{exit\_code, stdout }\OperatorTok{=} \ControlFlowTok{await}\NormalTok{ \_hidden\_grade(sb, task)          }\CommentTok{\# hidden checker, after the loop stops}
\NormalTok{reward }\OperatorTok{=}\NormalTok{ task.reward(}\DecValTok{0} \ControlFlowTok{if}\NormalTok{ exit\_code }\OperatorTok{==} \DecValTok{0} \ControlFlowTok{else} \DecValTok{1}\NormalTok{)           }\CommentTok{\# binary 1.0 / 0.0}
\NormalTok{samples }\OperatorTok{=} \ControlFlowTok{await}\NormalTok{ adapter.finish\_session(sid, base\_sample}\OperatorTok{=}\NormalTok{sample, reward}\OperatorTok{=}\NormalTok{reward)  }\CommentTok{\# exact tokens}
\end{Highlighting}
\end{Shaded}

Datasets are built once and seeded per episode: each row carries
\texttt{metadata.task} (family + spec), reference answers are dropped so
they cannot leak into prompts, and \texttt{generate.py} rebuilds the
\texttt{Task} from that metadata. The same code runs the checker for
MBPP, SWE, SQL, math, tool-use, RAG, CTF, and the rest --- the sandbox
is the only thing that changes.

\hypertarget{reading-the-episode-shape-into-a-config}{%
\subsection{Reading the episode shape into a
config}\label{reading-the-episode-shape-into-a-config}}

The families of §``Tasks frontier labs run in a sandbox'' differ mostly
in their \textbf{episode shape} --- turns per episode, response length
per turn, total episode length --- and that shape is what you size the
training config against. Two axes drive every knob. \textbf{Turns per
episode} sets \texttt{EPISODE\_MAX\_TURNS} and, multiplied by the
per-turn response, the \emph{sequence length} the trainer and inference
engine must hold: single-turn families (math, instruction following)
need only a long \texttt{-\/-rollout-max-response-len}; long-horizon
ones (SWE, web/computer use, multi-agent) need a large \textbf{context
window} and a \textbf{session that stays warm for the whole episode}.
\textbf{Where the tokens are} sets which side you scale (§``Scaling the
loop''): generation-heavy families (math, reasoning code) are
\emph{inference-bound} --- long responses, trivial checker;
observation-heavy families (SWE, RAG, web, tool use) push cost into
\emph{context growth and the sandbox}. Rule of thumb: budget context ≈
\texttt{max\_turns\ ×\ (response\_per\_turn\ +\ observation\_per\_turn)},
size \texttt{-\/-rollout-max-response-len} to the per-turn response at
the high end (reasoning traces), and set rollout-buffer memory from the
total episode length so long-horizon samples do not overflow it.

\hypertarget{the-episode-end-to-end}{%
\subsection{The episode, end to end}\label{the-episode-end-to-end}}

Taking one MBPP episode:

\begin{enumerate}
\def\labelenumi{\arabic{enumi}.}
\tightlist
\item
  \texttt{run.sh} sets \texttt{SANDBOX=agentcore-runtime},
  \texttt{DATASET}, \texttt{CONC}, \texttt{HARNESS\_WORKERS}, then
  \texttt{ray\ job\ submit\ …\ train\_async.py\ -\/-custom-generate-function-path\ …generate.generate}.
\item
  slime hands a prompt to \texttt{generate()}, which rebuilds the Task
  and queues the episode on the harness tier; a free dispatch worker
  opens \texttt{make\_sandbox("agentcore-runtime")} →
  \textbf{cold-start} the session's microVM (or takes a pre-opened one),
  then takes the least-loaded SGLang engine.
\item
  \textbf{Seed:} \texttt{write\_file} the task's files
  (\texttt{PROBLEM.md}, \ldots) and the stdlib loop runner.
\item
  \textbf{Agent loop, inside the session:} each model turn dials back to
  the engine's adapter (\texttt{{[}sglang-gen{]}}), each tool call
  (\texttt{write\_file} / \texttt{bash} / \texttt{read\_file}) runs
  locally in the session; the trainer long-polls the runner
  (\texttt{{[}sbx-op{]}}).
\item
  After the loop stops, the hidden checker runs in the session:
  \texttt{exit\_code\ ==\ 0} → solved; \textbf{reward = 1.0 / 0.0}.
\item
  The adapter's token record (\texttt{finish\_session}) → one
  loss-masked \texttt{Sample} → buffer → Megatron step → weight sync to
  SGLang.
\end{enumerate}

\textbf{CPU → GPU is one env var:} a CPU Runtime ARN (microVM, MBPP) → a
GPU-Runtime ARN (Instances, KernelBench, \texttt{allclose} on cuda) ---
\textbf{zero change to the loop}.

\hypertarget{scaling-sharing-and-network-isolation}{%
\subsection{Scaling, sharing, and network
isolation}\label{scaling-sharing-and-network-isolation}}

\begin{itemize}
\tightlist
\item
  \textbf{Scaling} follows §``Scaling the loop'': grow the side that
  dominates the per-op latencies, and balance with the harness tier: a
  pull queue of dispatch workers, the least-loaded SGLang engine per
  episode (sticky for its turns), and a pool of sandbox Runtimes
  round-robined sticky per episode.
\item
  \textbf{Sharing a GPU.} One session = one GPU, with no cross-session
  time-slicing {[}D5{]}. To pack several episodes onto one physical GPU
  you either run \textbf{many agents in one session} (they share the
  instance --- no isolation) or run \textbf{your own fractional
  scheduler} (this project's \texttt{RayGpuFarm}, \texttt{num\_gpus=0.5}
  actors, with an \texttt{nvcc} \texttt{flock} because a non-MIG L4
  wedges under compile storms). Isolated scaling means \emph{more
  instances}; there is \textbf{no MIG on L4}.
\item
  \textbf{Warm-instance reuse --- only for the Instances, and not used
  today.} A managed \textbf{microVM} cold-starts in a second or two, so
  the Terminal-Bench demo \textbf{does not reuse it}: every episode
  mints its own session, its own fresh microVM, which is also perfectly
  isolated. The \textbf{Instances} compute type is the opposite: a
  capacity-provider \texttt{g6} or \texttt{m5} takes
  \textbf{\textasciitilde60--80 s to cold-boot}. The removed Instance
  demos therefore tried one shared \texttt{runtimeSessionId} (one warm
  instance, cold boot paid once) plus a per-episode
  \texttt{sub\_session} with a fresh environment. Measured on hardware,
  this cut \texttt{m5} cold boots \textbf{201 → 1} over a 10-step run,
  and a fractional-GPU Ray farm on a g6 worked in the same shape
  (§``History: removed demos (2026-09-30)''). Detail:
  \texttt{docs/sandbox-fresh-env-on-persistent-instance.md}.
\item
  \textbf{Network isolation --- pull, don't expose.} Every Runtime is
  \emph{reached} by an outbound SigV4 \texttt{invoke\_agent\_runtime};
  \textbf{no Runtime has an inbound port}. The one connection the other
  way stays inside our VPC: the in-session agent loop dials the
  trainer's adapter ports, admitted SG→SG. IAM/SigV4 is the auth; the
  \texttt{session\_id} is routing, not auth. For
  \texttt{networkMode=VPC} the Runtime has no internet route, so it
  needs interface endpoints (\texttt{ecr.api}, \texttt{ecr.dkr},
  \texttt{s3}, \texttt{logs}, \texttt{sts}, \texttt{bedrock-agentcore}),
  and the ENI must land in an AZ that serves \texttt{bedrock-agentcore}
  --- miss one and create fails with a null reason {[}D10, D11{]}.
\end{itemize}

\hypertarget{can-agentcore-host-every-frontier-sandbox}{%
\subsection{Can AgentCore host every frontier
sandbox?}\label{can-agentcore-host-every-frontier-sandbox}}

Mapping the task families of §``Tasks frontier labs run in a sandbox''
onto AgentCore, the only hard limit is GPU \textbf{fractioning} --- not
any missing capability (✓ native · ◑ needs a BYO container / egress):

\begin{longtable}[]{@{}
  >{\raggedright\arraybackslash}p{(\columnwidth - 4\tabcolsep) * \real{0.3333}}
  >{\raggedright\arraybackslash}p{(\columnwidth - 4\tabcolsep) * \real{0.3333}}
  >{\raggedright\arraybackslash}p{(\columnwidth - 4\tabcolsep) * \real{0.3333}}@{}}
\toprule\noalign{}
\begin{minipage}[b]{\linewidth}\raggedright
Task / sandbox
\end{minipage} & \begin{minipage}[b]{\linewidth}\raggedright
✓?
\end{minipage} & \begin{minipage}[b]{\linewidth}\raggedright
How --- and the limit
\end{minipage} \\
\midrule\noalign{}
\endhead
\bottomrule\noalign{}
\endlastfoot
Code exec + unit tests (DeepSeek, Qwen) {[}2,3{]} & ✓ & Code
Interpreter, or your image on a Runtime \\
Math / logic verifier & ✓ & stdlib checker in the CPU microVM \\
Tool-use env (Kimi K2) {[}5{]} & ✓ & Python tools; MCP via Gateway \\
Per-repo Docker image (SWE) {[}6{]} & ◑ & Runtime runs your OCI image,
isolated per session --- but the managed microVM is \textbf{arm64-only}
and caps the image at \textbf{2 GB} (§3); slim it, defer installs to
\texttt{CMD}, or use the amd64 Instances substrate \\
Browser / computer use & ✓ & managed Browser tool \\
Networked RAG / web read & ◑ & egress --- public, or VPC + endpoints
{[}D11{]} \\
GPU kernel / CUDA (KernelBench) {[}7{]} & ◑ & Instances only; 1 session
= 1 GPU; \textbf{no MIG on L4} \\
\end{longtable}

\hypertarget{related-work-deepseek-dsec-a-purpose-built-sandbox-platform}{%
\subsection{Related work --- DeepSeek DSec, a purpose-built sandbox
platform}\label{related-work-deepseek-dsec-a-purpose-built-sandbox-platform}}

DSec (DeepSeek Elastic Compute) {[}31{]} is DeepSeek's in-house sandbox
platform for agentic RL, evaluation and environment building. It has
served every sandbox workload from DeepSeek V3.2 to V4.1. It is the best
published reference point for what a sandbox layer looks like when it is
built for RL rather than rented. One unit is about 160 CPU nodes with
30K cores. It serves about \textbf{3M sandboxes a day}, with peaks of
\textbf{380K running at once} and \textbf{more than 5,000 created per
second}.

Its design choices line up closely with the questions this report asks
of AgentCore:

\begin{itemize}
\tightlist
\item
  \textbf{Four backends behind one SDK:}

  \begin{itemize}
  \tightlist
  \item
    warm pre-created containers for short stateless calls, including GPU
    kernels on MIG slices;
  \item
    containers for SWE and tool use;
  \item
    Firecracker microVMs for security-sensitive and computer-use tasks;
  \item
    QEMU VMs for Android and GUI tasks.
  \end{itemize}

  The caller picks the backend, which is the same idea as our three
  sandbox tiers.
\item
  \textbf{Layered images, loaded on demand.} Base OS, task workspace and
  toolkit are versioned separately (11K bases × 100K workspaces in one
  week) and stacked at start time. They are read on demand from
  DeepSeek's shared filesystem: a sandbox reads only 4--13 \% of its
  image, and on an 8,192-container burst a full pull took 1.71× longer.
\item
  \textbf{Heavy overcommit.} About 90 \% of sandboxes use at most 5 \%
  of the CPU they request, so one node runs about 3,200 containers or
  800 microVMs. Latency-sensitive sandboxes get protected CPU
  scheduling.
\item
  \textbf{Coupled to the trainer.} The agent loop runs on the sandbox
  platform, outside the preemptible GPU job. When a job is preempted,
  its sandboxes are paused (containers frozen, microVMs snapshotted) and
  resume on the next request.
\item
  \textbf{Reward-hacking controls.} Their agents forged requests to the
  in-sandbox shell service, overwrote \texttt{/bin/bash}, port-scanned
  for code mirrors and crashed the kernel through \texttt{/proc}. The
  fixes are AppArmor rules plus per-sandbox network allowlists that can
  change between task stages.
\end{itemize}

\begin{longtable}[]{@{}
  >{\raggedright\arraybackslash}p{(\columnwidth - 4\tabcolsep) * \real{0.3333}}
  >{\raggedright\arraybackslash}p{(\columnwidth - 4\tabcolsep) * \real{0.3333}}
  >{\raggedright\arraybackslash}p{(\columnwidth - 4\tabcolsep) * \real{0.3333}}@{}}
\toprule\noalign{}
\begin{minipage}[b]{\linewidth}\raggedright
\end{minipage} & \begin{minipage}[b]{\linewidth}\raggedright
DSec {[}31{]}
\end{minipage} & \begin{minipage}[b]{\linewidth}\raggedright
AgentCore (this account, us-west-2, 2026-09-26)
\end{minipage} \\
\midrule\noalign{}
\endhead
\bottomrule\noalign{}
\endlastfoot
Sandbox creation rate & \textgreater{} 5,000 /s & \textbf{25 /s} new
Runtime sessions (adjustable) \\
Running at once & 380K & 5,000 Runtime sessions · 1,000 Code Interpreter
sessions \\
Distinct environments & 11K bases × 100K workspaces, composed at start &
one image per Runtime; max \textbf{1,000} Runtimes per account \\
Image delivery & on demand; start time barely depends on image size &
full image per cold boot (≈ 0.8 s stdlib, ≈ 15 s at 447 MB, grows with
size); 2 GB microVM cap. \textbf{Ours:} no image at all on Code
Interpreter (pandas/scipy/sklearn/torch-cpu preinstalled); install heavy
deps at boot to fit 2 GB; \textgreater{} 2 GB on a Runtime Instance
(60--80 s m5, ≈ 207 s g6 cold boot, paid once through a warm-shared
session) \\
Hiding the cold boot & not needed: lazy loading keeps the pull time low;
pre-warming rejected, since DSec's bottleneck is I/O volume, not latency
& \textbf{pre-open (ours)}: while the current batch runs, the sandboxes
of the next 1--2 batches are opened (image pulled, microVM booted),
which hides the image boot time; median wait for a ready sandbox 774 ms
→ \textbf{5 ms} (warm-start probe, small stdlib image), makespan −20
\% \\
Sandbox lifetime & p50 15--17 min, p99 \textgreater{} 3 h & 15-min idle
timeout by default; max lifetime 8 h on a microVM, 14 days on a Runtime
Instance \\
Survives a trainer stall & pause/resume, state kept & \textbf{no
pause/resume, no state kept}; a heartbeat keeps the sandbox alive \\
Agent loop & outside the GPU job & \textbf{outside the trainer too
(ours, every demo since 2026-09-28)}: the mini\_claude loop + its tools
in the episode's own Runtime session, one loop per sandbox (the removed
\texttt{demo\_cpu\_branching} ran the Claude Code CLI there); the
trainer keeps a dispatch tier + one adapter per engine. Results
committed before 2026-09-28 used the in-trainer loop \\
Reward-hack controls & AppArmor + network allowlists & microVM boundary;
VPC mode (a Landlock + no-TCP jail on a shared \texttt{m5} pool was
measured, then removed) \\
\end{longtable}

The headline gaps are the three things AgentCore does not expose: image
layers, session pause, and a creation rate in the thousands. One number
bears directly on this project. DSec's median sandbox lives 15--17
minutes, almost exactly AgentCore's default idle timeout, so bridging
trainer stalls --- now the default: heartbeat on for every remote own
Runtime, plus a 3600 s idle timeout --- (§``Hiding the container launch:
pre-open and heartbeat'') is the normal case rather than an edge case.
Two DSec ideas can be copied cheaply on AgentCore: moving the agent loop
out of the trainer process (now done in every demo: the loop runs in the
sandbox session and scales with it, §``Scaling the loop''), and
allowlisting the network per task through VPC mode. We also adopted
three of its metrics (batch time-to-ready, peak creation rate and
concurrency compared with the quotas, and session lifetime compared with
the idle limit); see §``The Terminal-Bench demo in depth'' for how they
are computed and what the first log shows.

In this project the policy is always \textbf{Qwen3-4B}, trained by
slime, with AgentCore as the sandbox and the harness running inside the
sandbox session. The next section applies all of the above to one demo,
RL on Terminal-Bench (§``The Terminal-Bench demo''), and §``The
Terminal-Bench demo in depth'' gives its pipeline, metrics and measured
results.

\hypertarget{the-terminal-bench-demo}{%
\section{The Terminal-Bench demo}\label{the-terminal-bench-demo}}

The demo lives in \texttt{code/demos/demo\_terminalbench\_kernelbench/}
(it was \texttt{code/demos/demo\_cpu\_terminalbench/}, the microVM-only
version described here, until 2026-10-06). It is a self-contained
recipe: a vendored library (\texttt{agentcore\_slime\_demo/}), the
Terminal-Bench tooling (\texttt{terminalbench/}), the deploy helpers
(\texttt{deploy/}) and six numbered step scripts. The policy is
\textbf{Qwen3-4B}, trained by \textbf{slime} with fully-async GRPO.
\textbf{AgentCore is the sandbox}, and the \textbf{harness is the
mini\_claude loop, running inside each episode's sandbox session}. It is
\textbf{REMOTE-ONLY}: \texttt{SANDBOX=agentcore-runtime} is the only
backend.

\hypertarget{which-sandbox-does-a-task-need-three-tiers}{%
\subsection{Which sandbox does a task need? Three
tiers}\label{which-sandbox-does-a-task-need-three-tiers}}

Pick the \textbf{lowest tier that works}: move to tier 2 for a run-level
need, and to tier 3 only when the task's environment itself must change.

\begin{longtable}[]{@{}
  >{\raggedright\arraybackslash}p{(\columnwidth - 6\tabcolsep) * \real{0.2500}}
  >{\raggedright\arraybackslash}p{(\columnwidth - 6\tabcolsep) * \real{0.2500}}
  >{\raggedright\arraybackslash}p{(\columnwidth - 6\tabcolsep) * \real{0.2500}}
  >{\raggedright\arraybackslash}p{(\columnwidth - 6\tabcolsep) * \real{0.2500}}@{}}
\toprule\noalign{}
\begin{minipage}[b]{\linewidth}\raggedright
Tier
\end{minipage} & \begin{minipage}[b]{\linewidth}\raggedright
Sandbox
\end{minipage} & \begin{minipage}[b]{\linewidth}\raggedright
Use it when
\end{minipage} & \begin{minipage}[b]{\linewidth}\raggedright
In this project
\end{minipage} \\
\midrule\noalign{}
\endhead
\bottomrule\noalign{}
\endlastfoot
\textbf{1 · Managed Code Interpreter} & no image, no Runtime &
\textbf{probes only}: a Code Interpreter session has no route back to
the trainer's adapter, so the in-session agent loop cannot run there &
warmstart probe \\
\textbf{2 · Own Runtime, generic image} & own Runtime + a stdlib image &
the task needs no task-specific environment (code, math, SQL, RAG, tool
calls) & warmstart probe \\
\textbf{3 · Own Runtime per row, custom image} & one image + one Runtime
per task & the environment differs per row & \textbf{Terminal-Bench} \\
\end{longtable}

Terminal-Bench is tier 3 on the managed arm64 microVM fleet. The fleet's
2048 MB image cap cannot be raised, so tasks whose image does not fit
are reported and left out.

\hypertarget{use-case-covered}{%
\subsection{Use case covered}\label{use-case-covered}}

Agentic RL where \textbf{every task needs a different OS image}
(packages, services, files, compilers) and the agent changes the whole
machine through the shell. Terminal-Bench {[}6c{]} ships each task as
its own container: a task \texttt{Dockerfile}, a hidden pytest oracle
(\texttt{tests/} + \texttt{run-tests.sh}) and a reference solution. This
is the \textbf{``large image pool''} shape: hundreds of bespoke images,
each served by its own AgentCore Runtime, with strict row → image →
Runtime routing.

\hypertarget{why-this-demo}{%
\subsection{Why this demo}\label{why-this-demo}}

\begin{itemize}
\tightlist
\item
  \textbf{It is the hardest sandbox shape in the literature datasets.}
  No generic image can serve it, and a shared session would leak one
  episode's changes into the next. The reasoning, code, math, SQL and
  tool-use datasets all run on one generic image (§``Language and
  environment per family''), so they do not exercise the sandbox layer
  the way Terminal-Bench does.
\item
  \textbf{It shows AgentCore scaling to a pool of \textasciitilde218
  Runtimes} built from the tasks' own Dockerfiles, with no per-task
  code: the row names its image, and \texttt{image\_routing.py} maps
  image → ARN.
\item
  \textbf{It shows the two latency fixes a per-episode microVM needs}:
  \textbf{pre-open}, which boots the next batches' sessions ahead of
  time, and a \textbf{heartbeat}, which keeps them alive across the
  weight-sync stall.
\item
  \textbf{It exercises every tier of the loop at once}: the sandbox (a
  fresh microVM per episode), the harness (one loop per session,
  dispatched by the harness tier) and inference (the least-loaded SGLang
  engine per episode).
\end{itemize}

\hypertarget{design-considerations}{%
\subsection{Design considerations}\label{design-considerations}}

\begin{itemize}
\tightlist
\item
  \textbf{One fresh microVM session per episode}
  (\texttt{terminalbench\_pool.py}): the tasks mutate the OS, so no
  state may cross episodes or GRPO siblings. The session is stopped when
  the episode is graded.
\item
  \textbf{Wrapping, not rewriting, the task image.}
  \texttt{deploy\_terminalbench.py} appends one layer to each task's own
  Dockerfile: a standalone Python (\texttt{/opt/acs/py}) and the stdlib
  AgentCore HTTP server. The task's own interpreter is never touched;
  Terminal-Bench images range from python3.8 to no Python at all.
\item
  \textbf{Strict routing.} The episode's \texttt{task.image} (the task
  name) is looked up in the image → ARN map
  (\texttt{arn\_by\_image.json}, exported as
  \texttt{AGENTCORE\_SANDBOX\_RUNTIME\_ARN\_BY\_IMAGE}). An unmapped
  task is an error. It is never sent to a fallback Runtime, which would
  run it on another task's image.
\item
  \textbf{Look-ahead pre-open} (\texttt{tb\_lookahead\_source.py},
  \texttt{TB\_PREOPEN\_BATCHES=2}): the data source peeks at the next
  two batches' rows, and the pool boots their sessions on the right task
  Runtime, so the cold start leaves the critical path.
\item
  \textbf{Heartbeat} (\texttt{HEARTBEAT=1}, default): a
  \texttt{SessionKeeper} pings every idle pre-opened session below the
  microVM idle timeout, so a trainer stall cannot reclaim it. A
  reclaimed session id comes back as an empty sandbox, and nothing
  reports the loss.
\item
  \textbf{Harness in the session, harness tier in the trainer.} The
  mini\_claude loop runs inside the session and dials the trainer's
  per-engine \texttt{AnthropicAdapter} back over the VPC.
  \texttt{HARNESS\_WORKERS} dispatch workers (default \texttt{RB\ ×\ NS}
  = 32) pull episodes off one queue. Each episode takes its session
  first, then the least-loaded SGLang engine, which stays sticky for all
  its turns.
\item
  \textbf{Oracle validation before training.}
  \texttt{validate\_oracle.py} checks each task on its own Runtime: an
  untouched environment must fail and the reference solution must pass.
  Tasks that fail the check are dropped from the data.
\item
  \textbf{What doesn't fit.} Of 229 single-container tasks, \textbf{7
  fail to build for arm64} and \textbf{4 exceed the 2 GB microVM cap}.
  They are listed in \texttt{terminalbench/tasks/images.csv} and left
  out of the data. Multi-container tasks are skipped, since one image
  maps to one sandbox.
\item
  \textbf{Network}: VPC mode with NAT egress (\texttt{TB\_NETWORK=VPC}).
  The loop dials the adapter back through an SG→SG rule on ports
  18100-18163, and \texttt{run-tests.sh} installs uv + pytest from the
  internet.
\item
  \textbf{Names shared with nothing else}: Runtimes are
  \texttt{acs\_tb\_\textless{}task\textgreater{}} and ECR repos are
  \texttt{acs-tb-\textless{}task\textgreater{}}.
\end{itemize}

\textbf{Isolation.} One episode = one microVM = one session. Every
episode gets a separate kernel, and nothing is shared between episodes
except the task's read-only image. The adapter port hole (SG→SG) is the
only inbound path.

\hypertarget{the-pool}{%
\subsection{The pool}\label{the-pool}}

\texttt{bash\ 2\_build\_image.sh} fetches the corpus
(\texttt{fetch\_terminalbench.py}) and builds one arm64 image per task
with \texttt{docker\ buildx}, then pushes it to ECR
\texttt{acs-tb-\textless{}task\textgreater{}}. The step is resumable and
takes about 2 h cold. \texttt{bash\ 4\_deploy\_runtime.sh} registers one
VPC Runtime per image and writes the map \texttt{arn\_by\_image.json}.
The files named below live in \texttt{terminalbench/tasks/}. Measured on
2026-09-30:

\begin{longtable}[]{@{}
  >{\raggedright\arraybackslash}p{(\columnwidth - 4\tabcolsep) * \real{0.3333}}
  >{\raggedleft\arraybackslash}p{(\columnwidth - 4\tabcolsep) * \real{0.3333}}
  >{\raggedright\arraybackslash}p{(\columnwidth - 4\tabcolsep) * \real{0.3333}}@{}}
\toprule\noalign{}
\begin{minipage}[b]{\linewidth}\raggedright
Stage
\end{minipage} & \begin{minipage}[b]{\linewidth}\raggedleft
Count
\end{minipage} & \begin{minipage}[b]{\linewidth}\raggedright
Source
\end{minipage} \\
\midrule\noalign{}
\endhead
\bottomrule\noalign{}
\endlastfoot
single-container tasks in the corpus & 229 & \texttt{images.csv} \\
failed to build for arm64 & 7 & same \\
over the 2 GB microVM cap & 4 & same \\
\textbf{Runtimes registered} & \textbf{218} &
\texttt{arn\_by\_image.json} \\
oracle-usable (untouched fails, reference passes) & 166 &
\texttt{oracle.csv} \\
dataset rows (train / eval, split by task) & 134 / 31 & \texttt{data/}
(train, eval) \\
\end{longtable}

\textbf{Cold start of all 218 Runtimes} (\texttt{probe\_coldstart.py} →
\texttt{coldstart.csv}, 1,090 opens). A first open of each Runtime, one
at a time, took p50 \textbf{0.60 s}, p90 0.96 s and max 3.3 s. A burst
of four concurrent opens per Runtime took p50 \textbf{4.7 s}, p90 12.3 s
and max 49 s. 1,085 of 1,090 opens succeeded; the other 5 were Runtime
health-check or initialization timeouts. So the pool opens fast one at a
time, and a burst is what pre-open has to hide.

\hypertarget{hiding-the-container-launch-pre-open-and-heartbeat}{%
\subsection{Hiding the container launch: pre-open and
heartbeat}\label{hiding-the-container-launch-pre-open-and-heartbeat}}

A fresh session per episode means a cold start per episode: allocate a
microVM, pull the task's image, boot it. The demo hides it with two
mechanisms, both on by default.

\textbf{Pre-open the next one or two batches.} slime hands the pool only
\texttt{CONC\ ×\ engines} episodes of the current batch, and draws the
next batch only when this rollout ends, so the pool alone cannot see far
enough ahead. A \textbf{look-ahead data source}
(\texttt{tb\_lookahead\_source.py}, wired with
\texttt{-\/-data-source-path}) replays slime's deterministic prompt
order: a shuffled dataset walked in order, with a fixed seed per epoch.
With the fully-async driver it slides per group. Drawn groups are
appended to the pool's demand, the next
\texttt{TB\_PREOPEN\_BATCHES\ ×\ RB} prompts become the look-ahead, and
aborted groups are un-reserved. The pool keeps that many upcoming
episodes booted, \texttt{TB\_PREOPEN\_BATCHES\ ×\ RB\ ×\ NS} = 2 × 4 × 8
= \textbf{64}, each on its own new session on its task's Runtime. When a
worker claims one, the next upcoming episode starts booting. A
misprediction is trimmed at the next draw and those episodes open on
demand. \texttt{TB\_PREOPEN\_BATCHES=0} turns it off.

\textbf{Heartbeat across the weight-sync stall.} While the trainer syncs
weights into SGLang, generation pauses and sessions receive no invokes.
An own-Runtime session that stays idle past its timeout is reclaimed. A
\texttt{SessionKeeper} (\texttt{session\_keeper.py}) registers every
open session and pings (\texttt{exec("true")}) any session that has been
idle for longer than
\texttt{WARM\_PING\_FRACTION\ ×\ WARM\_IDLE\_TIMEOUT} (0.5 × 900 s). The
keeper lives \textbf{outside} the harness, so no harness needs its own
keepalive. As a backstop, Runtimes created by
\texttt{deploy/agentcore\_deploy.py} get an idle timeout of
\texttt{RUNTIME\_IDLE\_TIMEOUT} = 3600 s. The cost is that an abandoned
session is billed for up to an hour.

\textbf{The \texttt{terminalbench/warmstart/} probe isolates both
mechanisms} (\texttt{bash\ terminalbench/warmstart/run\_warmstart.sh}:
host-only, no GPU, no model, about 7 min and a few cents). The same M =
24 MockPolicy \texttt{sql} episodes run in batches of 6 (think time 2 s)
against a stdlib Runtime \texttt{acs\_warmstart\_cpu} with a 60 s idle
timeout, and against the Code Interpreter. Recorded run
\texttt{warmstart\_20260925\_100103.*} (2026-09-25, us-west-2, all
rewards 1.0):

\begin{longtable}[]{@{}
  >{\raggedright\arraybackslash}p{(\columnwidth - 10\tabcolsep) * \real{0.1667}}
  >{\raggedright\arraybackslash}p{(\columnwidth - 10\tabcolsep) * \real{0.1667}}
  >{\raggedleft\arraybackslash}p{(\columnwidth - 10\tabcolsep) * \real{0.1667}}
  >{\raggedleft\arraybackslash}p{(\columnwidth - 10\tabcolsep) * \real{0.1667}}
  >{\raggedleft\arraybackslash}p{(\columnwidth - 10\tabcolsep) * \real{0.1667}}
  >{\raggedleft\arraybackslash}p{(\columnwidth - 10\tabcolsep) * \real{0.1667}}@{}}
\toprule\noalign{}
\begin{minipage}[b]{\linewidth}\raggedright
Backend
\end{minipage} & \begin{minipage}[b]{\linewidth}\raggedright
Mode
\end{minipage} & \begin{minipage}[b]{\linewidth}\raggedleft
Wait p50
\end{minipage} & \begin{minipage}[b]{\linewidth}\raggedleft
Episodes that waited
\end{minipage} & \begin{minipage}[b]{\linewidth}\raggedleft
Idle held
\end{minipage} & \begin{minipage}[b]{\linewidth}\raggedleft
Makespan
\end{minipage} \\
\midrule\noalign{}
\endhead
\bottomrule\noalign{}
\endlastfoot
own Runtime & cold & 774 ms & 24/24 & 0 s & 31.9 s \\
own Runtime & pre-open & \textbf{5 ms} & 8/24 & 75 s & 25.5 s (−20
\%) \\
Code Interpreter & cold & 1715 ms & 24/24 & 0 s & 28.7 s \\
Code Interpreter & pre-open & \textbf{0 ms} & 6/24 & 70 s & 24.3 s (−15
\%) \\
\end{longtable}

\begin{longtable}[]{@{}
  >{\raggedright\arraybackslash}p{(\columnwidth - 6\tabcolsep) * \real{0.2500}}
  >{\raggedright\arraybackslash}p{(\columnwidth - 6\tabcolsep) * \real{0.2500}}
  >{\raggedleft\arraybackslash}p{(\columnwidth - 6\tabcolsep) * \real{0.2500}}
  >{\raggedright\arraybackslash}p{(\columnwidth - 6\tabcolsep) * \real{0.2500}}@{}}
\toprule\noalign{}
\begin{minipage}[b]{\linewidth}\raggedright
Backend
\end{minipage} & \begin{minipage}[b]{\linewidth}\raggedright
Heartbeat
\end{minipage} & \begin{minipage}[b]{\linewidth}\raggedleft
Sessions surviving 150 s idle
\end{minipage} & \begin{minipage}[b]{\linewidth}\raggedright
Resume
\end{minipage} \\
\midrule\noalign{}
\endhead
\bottomrule\noalign{}
\endlastfoot
own Runtime & on (14 pings) & \textbf{3/3} & 104 ms, files intact \\
own Runtime & off & 0/3 & \textbf{silently reclaimed}: the same session
id returns a fresh, empty sandbox (948 ms) \\
Code Interpreter & on or off & 0/3 & expires at its session timeout
regardless of pings \\
\end{longtable}

What the probe shows:

\begin{itemize}
\tightlist
\item
  \textbf{Pre-open takes the boot off the critical path; only the first
  batch pays it.} The cost is idle held time, meaning booted sessions
  waiting for work. The Terminal-Bench smoke shows the same trade-off at
  full scale (about 95 \% of pre-opened time idle), which is why
  \texttt{TB\_PREOPEN\_BATCHES=1} is worth trying there.
\item
  \textbf{A heartbeat only helps an own Runtime.} Without it, a
  reclaimed session \emph{looks healthy} but has lost its files. Because
  the failure is silent, the demo keeps the heartbeat on by default.
\item
  \textbf{On the Code Interpreter a heartbeat does nothing.} A session
  ends at the \texttt{sessionTimeoutSeconds} set when it starts, however
  much traffic it gets.
\end{itemize}

\hypertarget{scaling-the-harness-with-inference-and-the-sandbox-the-harness-tier}{%
\subsection{Scaling the harness with inference and the sandbox: the
harness
tier}\label{scaling-the-harness-with-inference-and-the-sandbox-the-harness-tier}}

Rollout throughput is \texttt{min(inference,\ sandbox,\ harness)}
(§``Scaling the loop''). The harness has to scale \textbf{along with}
the other two, or a few coroutines silently become the ceiling. It
should also scale \textbf{with the sandbox rather than inside the
trainer}: the mini\_claude loop runs in the episode's own session, so N
sessions carry N loops, as DSec does {[}31{]}.
\texttt{TerminalBenchPool} subclasses the shared \texttt{HarnessPool}
(\texttt{harness\_pool.py}) and swaps only the sandbox tier. The three
tiers are sized independently and joined per episode, and no worker owns
an engine or a session:

\begin{itemize}
\tightlist
\item
  \textbf{harness}: \texttt{W\ =\ HARNESS\_WORKERS} dispatch workers in
  the trainer (default \texttt{RB\ ×\ NS} = 32). A free worker pulls the
  next episode off one shared queue.
\item
  \textbf{sandbox}: a fresh session per episode on its task's Runtime
  (pre-opened, heartbeated).
\item
  \textbf{inference}: the engine with the fewest in-flight episodes
  (ties rotate), sticky for all the episode's turns. This replaced the
  earlier four workers, each pinned to one engine.
\end{itemize}

\texttt{CONC} defaults to \texttt{⌈2W\ /\ E⌉}, so W running and W queued
episodes are admitted. The sandbox is taken before the engine, so an
episode waiting for a session does not count as engine load. Each
episode returns one bare \texttt{Sample}, the shape the fully-async
driver needs. Each episode's \texttt{{[}demo{]}} record shows which tier
binds: \texttt{qdepth} \textgreater{} 0 (harness),
\texttt{sandbox\_wait\_ms} \textgreater{} 0 (sandbox), or a climbing
\texttt{engine\_inflight} (inference). A session that cannot host the
loop (no \texttt{python3}, or the adapter unreachable) raises
\texttt{HarnessPlacementError}. There is no in-trainer fallback.

\hypertarget{run-it-step-by-step}{%
\subsection{Run it step by step}\label{run-it-step-by-step}}

Run the numbered scripts from
\texttt{code/demos/demo\_terminalbench\_kernelbench/} \textbf{on the
host}; the container steps \texttt{docker\ exec} into
\texttt{agentcore\_p4} themselves.

\begin{Shaded}
\begin{Highlighting}[]
\FunctionTok{bash}\NormalTok{ 2\_build\_image.sh      }\CommentTok{\# fetch the corpus (once) + build/push one arm64 image per task (resumable, \textasciitilde{}2 h cold)}
\FunctionTok{bash}\NormalTok{ 3\_push\_ecr.sh         }\CommentTok{\# report only (step 2 pushed)}
\FunctionTok{bash}\NormalTok{ 4\_deploy\_runtime.sh   }\CommentTok{\# one VPC Runtime per image {-}\textgreater{} terminalbench/tasks/arn\_by\_image.json}
\FunctionTok{bash}\NormalTok{ 1\_prepare\_data.sh     }\CommentTok{\# rows for the deployed tasks only (train/eval split by task)}
\FunctionTok{bash}\NormalTok{ 5\_run.sh              }\CommentTok{\# 10{-}step smoke {-}\textgreater{} logs/\textless{}tag\textgreater{}.\{log,csv,eval.csv,png\}}
\FunctionTok{bash}\NormalTok{ 6\_cleanup.sh          }\CommentTok{\# delete every acs\_tb\_* Runtime (ECR images kept)}
\end{Highlighting}
\end{Shaded}

Setting \texttt{TB\_LIMIT=2} (or \texttt{TB\_TASKS=a,b}) on steps 2 and
4 validates the pipeline on a couple of tasks first, and
\texttt{validate\_oracle.py} checks each task's oracle on its Runtime.
The long run is \texttt{convergence/20\_cpu\_terminalbench.sh} (under
\texttt{code/demos/}). Run it \textbf{inside \texttt{agentcore\_p4}},
after steps 2-4 and 1. It sets \texttt{NR=201} (200 train steps),
\texttt{EVAL\_INTERVAL=20}, \texttt{EVAL\_NS=4} and
\texttt{SAVE\_INTERVAL=100} (about 56 GB per saved checkpoint). Never
run two \texttt{run.sh} at once: the second one's Ray start kills the
first one's job. Training starts from the base Megatron checkpoint
\texttt{Qwen3-4B\_torch\_dist} in \texttt{code/p4\_slime\_train/ckpt/}
(override with \texttt{CKPT\_BASE}).

\begin{longtable}[]{@{}
  >{\raggedright\arraybackslash}p{(\columnwidth - 4\tabcolsep) * \real{0.3333}}
  >{\raggedright\arraybackslash}p{(\columnwidth - 4\tabcolsep) * \real{0.3333}}
  >{\raggedright\arraybackslash}p{(\columnwidth - 4\tabcolsep) * \real{0.3333}}@{}}
\toprule\noalign{}
\begin{minipage}[b]{\linewidth}\raggedright
Knob
\end{minipage} & \begin{minipage}[b]{\linewidth}\raggedright
Default
\end{minipage} & \begin{minipage}[b]{\linewidth}\raggedright
Meaning
\end{minipage} \\
\midrule\noalign{}
\endhead
\bottomrule\noalign{}
\endlastfoot
\texttt{NR} / \texttt{EVAL\_INTERVAL} & 11 / 10 (\texttt{5\_run.sh}) &
rollouts (NR−1 train steps) / eval cadence \\
\texttt{RB} / \texttt{NS} & 4 / 8 & prompts per step / GRPO group
size \\
\texttt{TB\_PREOPEN\_BATCHES} & 2 & batches of sessions booted ahead (0
= open on demand) \\
\texttt{HEARTBEAT} & 1 & keep idle sessions alive across trainer
stalls \\
\texttt{EPISODE\_MAX\_TURNS} / \texttt{EPISODE\_MAX\_TOKENS} & 8 / 1024
& agent turns / tokens per turn \\
\texttt{TB\_EXEC\_TIMEOUT} & 180 & per bash-tool command, seconds \\
\texttt{TB\_BUILD\_CONC} / \texttt{TB\_REBUILD} & 16 / 0 & parallel
image builds / force rebuild \\
\end{longtable}

\hypertarget{the-terminal-bench-demo-in-depth-pipeline-metrics-and-measured-results}{%
\section{The Terminal-Bench demo in depth --- pipeline, metrics, and
measured
results}\label{the-terminal-bench-demo-in-depth-pipeline-metrics-and-measured-results}}

This section is the deep dive behind §``The Terminal-Bench demo'': the
pipeline and its configs, the metric set, how to size the two episode
knobs, the measured smokes, and the plan to convergence. The numbers are
read from this project's own run artifacts: the per-step CSVs and the
\texttt{{[}sbx-op{]}}, \texttt{{[}sglang-gen{]}} and \texttt{{[}demo{]}}
records in the logs. They come from \textbf{short runs} (10-step
smokes), so treat them as the shape and order of magnitude of the loop,
not as converged results. Where a metric is only \emph{instrumented}
(emitted by the code but not yet populated in a run), this is said
explicitly.

\hypertarget{the-pipeline-one-episode-one-step}{%
\subsection{The pipeline --- one episode, one
step}\label{the-pipeline-one-episode-one-step}}

\begin{figure}[htbp]
\centering
\resizebox{0.98\textwidth}{!}{%
\begin{tikzpicture}[
  font=\small,
  box/.style={draw, rounded corners, minimum height=1.0cm, minimum width=2.0cm, align=center, thick, fill=white},
  gpu/.style={box, fill=blue!8},
  cpu/.style={box, fill=black!4},
  focus/.style={box, fill=orange!18},
  lbl/.style={font=\scriptsize, align=center, fill=white, inner sep=1.5pt},
  ptitle/.style={font=\scriptsize\bfseries, text=black!55},
  panel/.style={draw, rounded corners, dashed, thick, inner sep=0.30cm},
  >={Stealth[]},
  node distance=0.55cm,
]
\node[cpu]              (cold) {pre-opened\\\scriptsize microVM session};
\node[cpu, right=of cold](seed) {seed\\\scriptsize \texttt{write\_file}$\times2$};
\node[focus, right=of seed](loop){turn loop\\\scriptsize $\le$\texttt{EPISODE\_MAX\_TURNS}};
\node[cpu, right=of loop](reward){reward\\\scriptsize hidden pytest oracle};
\node[cpu, right=of reward](tok) {token stream\\\scriptsize loss-masked \texttt{Sample}};
\node[gpu, right=of tok] (buf)  {rollout\\buffer};
\node[gpu, right=of buf] (step) {Megatron\\\scriptsize GRPO step};
\node[gpu, right=of step](sync) {weight\\sync};
% flow
\draw[->,thick] (cold)--(seed);
\draw[->,thick] (seed)--(loop);
\draw[->,thick] (loop)--(reward);
\draw[->,thick] (reward)--(tok);
\draw[->,thick] (tok)--(buf);
\draw[->,thick] (buf)--(step);
\draw[->,thick] (step)--(sync);
% feedback: weight sync -> inference used in the loop
\draw[->,thick,dashed] (sync.north) to[out=120,in=60] node[lbl,above]{weight sync $\to$ SGLang (new policy; the staleness source)} (loop.north);
% panels
\begin{scope}[on background layer]
  \node[panel, fill=orange!4, fit=(cold)(seed)(loop)(reward)(tok)] (proll) {};
  \node[panel, fill=blue!3,  fit=(buf)(step)(sync)]                 (ptr)   {};
\end{scope}
\node[ptitle, anchor=north] at ([yshift=-2pt]proll.south) {Rollout --- per episode (agent loop inside the AgentCore sandbox session)};
\node[ptitle, anchor=north] at ([yshift=-2pt]ptr.south) {Training --- per step};
\end{tikzpicture}%
}
\caption{The Terminal-Bench pipeline. An episode takes its pre-opened session on its task's Runtime, seeds the workspace, then runs the \textbf{turn loop} --- each turn a generation through the proxy (\texttt{[sglang-gen]}) followed by a tool \texttt{exec} in the sandbox (\texttt{[sbx-op]}). The task's hidden oracle gives the reward; \texttt{build\_token\_stream} emits a loss-masked \texttt{Sample} into the rollout buffer; Megatron takes a GRPO step and \textbf{weight-syncs} the new policy back to SGLang --- the point at which rollouts generated on an older policy version become \emph{stale}. The two labeled latencies are the two things the scaling decision (§"Scaling the loop") is made from.}
\end{figure}

\textbf{Architecture.} The topology is the one in §``The call path''.
slime's \texttt{generate()} hook hands each episode to the harness tier,
which takes the episode's pre-opened session on its task's Runtime and
then the least-loaded engine. The agent loop runs inside the session.
Each \textbf{model turn} dials back to the proxy, which returns the
exact sampled tokens and logprobs (the TITO property). Each \textbf{tool
call} (\texttt{bash} / \texttt{write\_file} / \texttt{read\_file}) runs
in the session. After the loop stops, the hidden oracle produces the
reward.

\textbf{The four configs: trainer, inference, harness, sandbox.}

\begin{longtable}[]{@{}
  >{\raggedright\arraybackslash}p{(\columnwidth - 4\tabcolsep) * \real{0.3333}}
  >{\raggedright\arraybackslash}p{(\columnwidth - 4\tabcolsep) * \real{0.3333}}
  >{\raggedright\arraybackslash}p{(\columnwidth - 4\tabcolsep) * \real{0.3333}}@{}}
\toprule\noalign{}
\begin{minipage}[b]{\linewidth}\raggedright
Component
\end{minipage} & \begin{minipage}[b]{\linewidth}\raggedright
Where it runs
\end{minipage} & \begin{minipage}[b]{\linewidth}\raggedright
Config
\end{minipage} \\
\midrule\noalign{}
\endhead
\bottomrule\noalign{}
\endlastfoot
\textbf{Trainer} & 4 GPUs on the 8×L40S box & Megatron-LM, GRPO,
fully-async rollout; \texttt{RB=4\ ×\ NS=8} =
\texttt{global\_batch\_size=32}; \texttt{ACTOR\_GPUS=4} \\
\textbf{Inference} & 4 GPUs on the same box & SGLang, 2 engines
(\texttt{GPUS\_PER\_ENGINE=2}), \texttt{qwen25} tool-call parser +
\texttt{qwen3} reasoning parser \\
\textbf{Harness} & mini\_claude inside the session; dispatch tier +
adapters in the trainer & \texttt{HARNESS\_WORKERS=32}, least-loaded
engine per episode, \texttt{CONC\ =\ ⌈2W/E⌉}; ≤ 8 turns × 1024 tokens
per turn \\
\textbf{Sandbox} & 218 AgentCore arm64 microVM Runtimes, VPC + NAT & one
fresh \texttt{runtimeSessionId} per episode on the task's Runtime;
\texttt{TB\_PREOPEN\_BATCHES=2}, \texttt{HEARTBEAT=1},
\texttt{TB\_EXEC\_TIMEOUT=180} \\
\end{longtable}

This is the disaggregation of §``Agentic RL and its reference
architecture'' made concrete. The GPUs go to training and inference, the
harness rides in the sandbox session it drives, and the sandbox is the
isolated remote seam.

\textbf{Step by step (one episode).}

\begin{enumerate}
\def\labelenumi{\arabic{enumi}.}
\tightlist
\item
  Take the episode's pre-opened session, or cold-start one on the task's
  Runtime.
\item
  Seed the task prompt and the stdlib loop runner
  (\texttt{session\_runner.py}).
\item
  Run the turn loop (≤ \texttt{EPISODE\_MAX\_TURNS}): generate through
  the proxy (\texttt{{[}sglang-gen{]}}), then run the tool in the
  session.
\item
  After the loop, inject the oracle (\texttt{tests/} +
  \texttt{run-tests.sh}) and run it.
\item
  Compute the reward: 1.0 if every oracle test passes, else 0.0.
\item
  The adapter's token record becomes a loss-masked \texttt{Sample},
  which goes to the rollout buffer, then a Megatron GRPO step, then a
  weight sync back to SGLang.
\end{enumerate}

\textbf{Lifecycle.} One session = one episode. The first invoke
cold-starts the microVM, and each further op under the same
\texttt{runtimeSessionId} reaches the same warm VM.
\texttt{StopRuntimeSession} ends the session when the episode is graded.
The idle timeout (\textasciitilde15 min by default, 3600 s on Runtimes
created by \texttt{deploy/agentcore\_deploy.py}) and the 8 h max
lifetime are only fallbacks. Under asynchronous training, the
weight-sync pause could trip the idle clock of a pre-opened session; the
heartbeat prevents that.

\textbf{What is measured, and what is only instrumented.}

\begin{longtable}[]{@{}
  >{\raggedright\arraybackslash}p{(\columnwidth - 4\tabcolsep) * \real{0.3333}}
  >{\raggedright\arraybackslash}p{(\columnwidth - 4\tabcolsep) * \real{0.3333}}
  >{\raggedright\arraybackslash}p{(\columnwidth - 4\tabcolsep) * \real{0.3333}}@{}}
\toprule\noalign{}
\begin{minipage}[b]{\linewidth}\raggedright
Metric
\end{minipage} & \begin{minipage}[b]{\linewidth}\raggedright
Meaning
\end{minipage} & \begin{minipage}[b]{\linewidth}\raggedright
Status
\end{minipage} \\
\midrule\noalign{}
\endhead
\bottomrule\noalign{}
\endlastfoot
\texttt{rollout/raw\_reward}, \texttt{eval/terminalbench} & mean oracle
reward per step; held-out solve rate & \textbf{measured} \\
\texttt{{[}demo{]}\ episode\_end} fields & per episode: \texttt{solved},
\texttt{tests\_passed}, \texttt{tests\_total}, \texttt{turns},
\texttt{tool\_calls}, \texttt{stop}, \texttt{agent\_ms},
\texttt{grade\_ms}, \texttt{harness\_where}, \texttt{infra\_error} &
\textbf{measured} \\
harness-tier fields & \texttt{worker}, \texttt{qdepth}, \texttt{engine},
\texttt{engine\_inflight}, \texttt{sandbox\_wait\_ms} &
\textbf{measured} \\
pre-open fields & \texttt{open\_wait\_ms}, warm hits,
\texttt{preopen\_idle\_ms} & \textbf{measured} \\
\texttt{perf/sandbox\_*}, \texttt{perf/sglang\_gen\_*} & warm data-op vs
session-open latency; per-turn generation latency & \textbf{measured} \\
\texttt{perf/step\_time}, \texttt{wait\_time\_ratio},
\texttt{actor\_train\_tok\_per\_s} & step wall-clock; fraction of the
step the trainer waits on rollout; trainer throughput &
\textbf{measured} \\
\texttt{ready\_span\_s}, \texttt{open\_peak\_per\_s},
\texttt{sessions.csv} & batch time-to-ready, peak opens per second vs
the 25/s quota, session lifetime vs the idle limit (after DSec {[}31{]})
& \textbf{measured} (re-parsed from \texttt{{[}sbx-op{]}}) \\
staleness, discard ratio, importance-sampling ratio & off-policy signals
& \textbf{instrumented or configured only}, no value recorded \\
\end{longtable}

\texttt{metrics.py} parses the log into a per-step CSV
(\texttt{\textless{}tag\textgreater{}.csv}), a per-eval CSV
(\texttt{\textless{}tag\textgreater{}.eval.csv}), a \texttt{{[}demo{]}}
summary (\texttt{\textless{}tag\textgreater{}.demo.csv}) and a PNG.
Sandbox latency is split into \textbf{warm data ops} and \textbf{session
opens}, and a cold boot is counted only when an open exceeds
\texttt{COLD\_BOOT\_MS} (5 s). CPU and memory used versus requested,
DSec's overcommit signal, is \textbf{not measured yet}: it needs
per-exec resource usage from inside the sandbox server.

\hypertarget{sizing-the-config-response-length-per-turn-and-turns-per-episode}{%
\subsection{Sizing the config --- response length per turn, and turns
per
episode}\label{sizing-the-config-response-length-per-turn-and-turns-per-episode}}

\textbf{Max response length per turn} (\texttt{EPISODE\_MAX\_TOKENS},
1024; \texttt{-\/-rollout-max-response-len}). Set it from the p95 of the
policy's own per-turn response for the family. \textbf{Too small}
truncates the turn mid-thought: the 2026-09-30 smoke stopped 115 of 449
episodes on \texttt{max\_tokens}. \textbf{Too large} wastes KV cache and
slows decode. \texttt{MAX\_TOKENS\_PER\_GPU=16384} must hold one full
trajectory: 8 × (1024 generated + \textasciitilde600 tool tokens) plus
the prompt. Each tool result is clipped to
\texttt{MINI\_AGENT\_MAX\_TOOL\_CHARS=2000}, which bounds the
trajectory.

\textbf{Max turns per episode} (\texttt{EPISODE\_MAX\_TURNS}, 8). This
cap is the single biggest lever on episode cost. Each extra turn adds
one generation and one sandbox round trip, and it grows the context for
every later turn. More turns raise the ceiling on reward, but they
lengthen episodes, which under asynchronous training means \textbf{more
staleness} and longer-held sessions. Terminal-Bench is a long-horizon
family, and 8 turns is a smoke-sized budget. Both knobs are env vars
read by \texttt{run.sh}, so no redeploy is needed to change them.

\hypertarget{measured-the-2026-09-23-smoke-first-folder}{%
\subsection{Measured --- the 2026-09-23 smoke (first
folder)}\label{measured-the-2026-09-23-smoke-first-folder}}

This run used the first \texttt{demo\_cpu\_terminalbench} folder
(removed 2026-10-06), before the 2026-09-28 merge
(\texttt{\_smoke\_logs/lanes\_tb\_0923/terminalbench.log}). It ran 382
episodes in \textasciitilde70 min, ended \texttt{Job\ succeeded} and had
0 infrastructure errors.

\begin{itemize}
\tightlist
\item
  \textbf{Cold start fully hidden.} \texttt{open\_wait\_ms} p95 was 0,
  the open p50 was 0.66 s, and 100 \% of opens were warm hits.
\item
  \textbf{The price: idle pre-opened time.} \texttt{preopen\_idle\_ms}
  p50 was 7.3 min against an episode p50 of 21 s, so about \textbf{95 \%
  of the pre-opened microVM time was idle}.
\item
  \textbf{Sessions.} Re-parsed, the run had 414 sessions, with a
  lifetime p50 of \textbf{471 s} and a p95 of \textbf{1841 s}.
  \textbf{55} of them outlived 900 s: the heartbeat kept them alive, and
  every open pinged back healthy. Opens peaked at \textbf{36 per
  second}, above the 25/s quota, with \textbf{no throttling} observed.
  This is an observation, not an AWS statement.
\item
  \textbf{Sparse reward for a 4B policy.} 4 of 382 episodes were solved
  and \texttt{eval/terminalbench} was 0.0, though 42 \% of episodes
  passed at least one test.
\end{itemize}

\hypertarget{measured-the-2026-09-30-smoke-current-folder}{%
\subsection{Measured --- the 2026-09-30 smoke (current
folder)}\label{measured-the-2026-09-30-smoke-current-folder}}

The 10-step smoke \texttt{bash\ 5\_run.sh} started at 13:48 UTC on
2026-09-30 (\texttt{logs/run-agentcore-runtime\_20260930\_134858.log}).
\textbf{The run was killed on purpose at 14:16, after 9 of its 10 train
steps}, to free the box for the Terminal-Bench-only cleanup. There was
no final rollout, no eval and no CSV, so the numbers below are read from
the log. \texttt{rollout/raw\_reward} per step was 0, 0.16, 0, 0, 0.25,
0, 0, 0.09, 0, and \texttt{perf/step\_time} was 30-41 s in steady state
(140 s and 136 s on the first step and on a step held up by one 504 s
exec). Over 460 \texttt{{[}demo{]}\ episode\_end} records, \textbf{19
episodes were solved} and 216 passed at least one oracle test, with
\textbf{0 infrastructure errors}. 449 ran their loop in the session
(\texttt{harness\_where=session}). Stop reasons were the turn budget
175, \texttt{end\_turn} 159 and \texttt{max\_tokens} 115, and
\texttt{agent\_ms} had a p50 of 12.2 s and a p95 of 57.8 s, with a turns
p50 of 6. The 61 tracebacks in the log are all
\texttt{asyncio.CancelledError} from the fully-async driver cancelling
surplus groups, which is benign. \texttt{train/pg\_loss} was 0 on most
steps because most GRPO groups had zero reward variance. The log tail
shows slow execs with transient \texttt{RuntimeClientError} retries.
\textbf{A re-run of the 10-step smoke is needed} before convergence run
20.

\hypertarget{measured-the-2026-09-30-re-run-current-code}{%
\subsection{Measured --- the 2026-09-30 re-run (current
code)}\label{measured-the-2026-09-30-re-run-current-code}}

\texttt{bash\ 5\_run.sh} re-run at 15:31 UTC
(\texttt{logs/smoke\_tbonly-agentcore-runtime\_20260930\_153115.\{log,csv,eval.csv,png\}})
completed 10 train steps and the final eval in \textbf{22 min} (the
2026-09-23 smoke took about 70 min). There were 449 episodes (387 train
including the fully-async over-draw, 62 eval) with \textbf{0
\texttt{infra\_error}}. 22 of 387 train episodes and \textbf{1 of 62
eval} episodes were solved (\texttt{eval/terminalbench} = 0.016).
Per-step \texttt{rollout/raw\_reward} was 0, 0.09, 0.09, 0, 0.25, 0, 0,
0.09, 0, 0.16. \texttt{perf/step\_time} p50 was 40 s. No train episode
waited for a boot (pre-open), and one heartbeat ping was needed
(pre-open idle max 194 s). Train episode p50/p90: agent loop 12.1/36.2
s, grading 8.9/20.6 s, whole episode 24.9/71.7 s. Only 4 of 49 GRPO
groups had a nonzero-variance reward. One bug was fixed to get here:
eval hung after its last episode because \texttt{get\_tb\_pool} was a
process-wide singleton bound to the fully-async worker's event loop, so
eval's futures were resolved from the wrong thread. It is now one pool
per event loop.

\hypertarget{experiment-plan-training-to-convergence}{%
\subsection{Experiment plan --- training to
convergence}\label{experiment-plan-training-to-convergence}}

A demo is ``done'' only when it shows \textbf{learning to convergence},
not a smoke run. The convergence bar is:

\begin{enumerate}
\def\labelenumi{\arabic{enumi}.}
\tightlist
\item
  Train enough steps that the \textbf{held-out
  \texttt{eval/terminalbench} rises and then plateaus}, with eval every
  \texttt{EVAL\_INTERVAL} steps.
\item
  \texttt{rollout/raw\_reward} and \texttt{eval/terminalbench} trend up
  and then flatten in the run's PNG.
\item
  \texttt{train/grad\_norm} and \texttt{train/kl\_loss} stay sane: the
  policy moves without diverging.
\item
  Only then apply §``Scaling the loop'': scale the \emph{measured}
  bottleneck and show the \texttt{perf/step\_time} gain.
\end{enumerate}

\begin{longtable}[]{@{}
  >{\raggedright\arraybackslash}p{(\columnwidth - 2\tabcolsep) * \real{0.5000}}
  >{\raggedright\arraybackslash}p{(\columnwidth - 2\tabcolsep) * \real{0.5000}}@{}}
\toprule\noalign{}
\begin{minipage}[b]{\linewidth}\raggedright
Run \texttt{convergence/20\_cpu\_terminalbench.sh}
\end{minipage} & \begin{minipage}[b]{\linewidth}\raggedright
\end{minipage} \\
\midrule\noalign{}
\endhead
\bottomrule\noalign{}
\endlastfoot
Steps / eval & \texttt{NR=201} (200 train steps), eval every 20
(\texttt{EVAL\_NS=4}) \\
Held-out start → target & \texttt{eval/terminalbench} 0.0 → any non-zero
plateau \\
Measured bottleneck & session open (hidden by pre-open) + sparse
reward \\
What to scale & pre-open depth vs idle cost; more
\texttt{HARNESS\_WORKERS} + engines \\
\end{longtable}

Two levers matter most before scale:

\begin{itemize}
\tightlist
\item
  \textbf{Reward signal first.} At 4B the oracle reward is sparse. An
  easy-first curriculum (by the task's \texttt{difficulty} field) or a
  partial-credit reward from \texttt{tests\_passed\ /\ tests\_total}
  would give GRPO non-zero advantages on more groups.
\item
  \textbf{Pre-open cost.} \texttt{TB\_PREOPEN\_BATCHES=1} halves the
  booted-ahead sessions. The 2026-09-23 smoke shows how much idle time
  the default of 2 buys.
\end{itemize}

\hypertarget{history-removed-demos-2026-09-30}{%
\section{History: removed demos
(2026-09-30)}\label{history-removed-demos-2026-09-30}}

On 2026-09-30 the project was cut down to the Terminal-Bench demo. The
rest is still in git history. Each line below gives what the removed
item was for and what it established.

\begin{itemize}
\tightlist
\item
  \textbf{\texttt{demo\_cpu\_tasks}}: nine literature datasets (GSM8K,
  MATH-500, DeepScaleR, MBPP, HumanEval, Gretel SQL, HotpotQA, xLAM,
  MCP), all on one generic stdlib image. Its smokes showed that most
  families need no custom image: \texttt{eval/gsm8k} 0.90,
  \texttt{eval/math500} 0.89, \texttt{eval/mbpp} 0.56. Its
  \texttt{SCALE=1} preset (4 engines : 256 harness workers : 256
  sessions, 592 episodes) was inference-bound, which is how the harness
  tier was proven.
\item
  \textbf{\texttt{demo\_gpu\_single}, \texttt{demo\_cpu\_instance}}: one
  session per capacity-provider Instance (g6/L4 KernelBench; m5). The
  warm-shared KernelBench run scored \texttt{eval/kernelbench} 0.65. It
  also showed that an Instance costs \textasciitilde90 s or more of cold
  start (torch pip-installed to stay under the image cap), and that
  there is no usable GPU-instance snapshot.
\item
  \textbf{\texttt{demo\_gpu\_farm}}: a Ray GPU-sharing farm (fractional
  or whole L4, local or on a remote g6). It showed that GPU sharing
  works, and that a non-MIG L4 wedges under sustained nvcc load.
\item
  \textbf{\texttt{demo\_cpu\_branching}}: the real Claude Code CLI as
  harness, where one episode fans out into sub-agent and compaction
  leaves (mean fan-out \textasciitilde5, \texttt{eval/swe} 0.47). Its
  advantages were 0 because of a grouping bug (the group was keyed by
  \texttt{sample.index} rather than
  \texttt{index\ //\ n\_samples\_per\_prompt}). The bug was fixed, but
  the demo was never re-run.
\item
  \textbf{\texttt{demo\_cpu\_swebench}} (SWE-bench Verified on the same
  pool code): one Runtime per instance on its Epoch AI arm64 image,
  reward from the hidden FAIL\_TO\_PASS + PASS\_TO\_PASS tests. It had
  \textbf{20 arm64 instances} and scored \textbf{reward 0} with
  Qwen3-4B. It was removed on 2026-09-30, and the planned 404-instance
  expansion was dropped. The 2 GB cap was probed on the way: a 2041.7
  MiB push was accepted and a 2061.7 MiB push was rejected. Its EKS/ECS
  twins, the same task on self-managed containers, went with it.
\item
  \textbf{Convergence scripts 01-19} and the shared top-level code: the
  library \texttt{code/agentcore\_slime\_demo} plus
  \texttt{code/deploy}, \texttt{code/scripts} and \texttt{code/tests}.
  The surviving demo vendors its own copy of the library and deploy
  helpers.
\end{itemize}

\hypertarget{open-problems-in-production-agentic-rl-infrastructure-and-algorithms}{%
\section{Open problems in production agentic RL --- infrastructure and
algorithms}\label{open-problems-in-production-agentic-rl-infrastructure-and-algorithms}}

The workflow in this report is deliberately small --- one policy, a
harness that runs inside a per-episode sandbox. Running agentic RL at
\emph{frontier} scale surfaces a set of recurring hard problems that the
open reports (2024--2026) now describe in detail. This section catalogs
them in two groups --- infrastructure/systems and algorithms --- with
the mitigations people actually ship, and notes where this project's
design already meets one and where it remains open.

\hypertarget{infrastructure-and-systems}{%
\subsection{Infrastructure and
systems}\label{infrastructure-and-systems}}

\textbf{Long-tail stragglers and partial rollouts.} Rollout lengths are
heavily right-skewed, so a synchronous step idles the whole cluster
waiting on the last few long trajectories --- and generation is more
than 90 \% of RL runtime {[}20{]}. Mitigations: \emph{partial rollouts}
that cap output and resume the unfinished portion next iteration (Kimi
k1.5 {[}17{]}; Kimi K2 {[}5{]}); \emph{active partial rollouts} that
over-provision and recycle interrupted requests (APRIL, +22--44 \%
rollout throughput, now in slime {[}28{]}); \emph{interruptible
generation} that resumes decoding after a weight update (AReaL
{[}18{]}); \emph{skewness-aware routing} of long prompts (StreamRL
{[}20{]}); and \emph{dynamic sampling} that drops zero-gradient groups
(DAPO {[}14{]}).

\textbf{Weight-sync cost and staleness.} Trainer and inference are
separate backends; after each optimizer step, new weights must reach
every rollout replica. Synchronously this is a serialization bubble; at
1T scale it moves hundreds of GB. Mitigations: \emph{fast weight
transfer} (Moonshot's checkpoint-engine updates all of Kimi-K2's
\textasciitilde1T parameters in \textasciitilde20 s {[}29{]});
\emph{zero-redundancy resharding} between train and generate layouts
(veRL / HybridFlow {[}19{]}); and \emph{decoupled, ``one-step-off''
async RL with a bounded staleness} and a staleness-aware objective
(AReaL, up to \textasciitilde2.77× {[}18{]}; Asynchronous RLHF
{[}21{]}). \emph{This project uses slime's fully-async rollout in the
Terminal-Bench demo, so staleness is the accepted cost --- and it is
exactly the signal that is instrumented-yet-unrecorded here (§``The
Terminal-Bench demo in depth'').}

\textbf{Goodput vs MFU, and placement.} An iteration interleaves a
bandwidth-bound, long-tailed generation phase with a compute-bound
training phase, so one engine's GPUs idle while the other works ---
which is why MFU is misleading for RL and practitioners reason in
\emph{goodput} (useful samples per GPU-second). \emph{This project's
measured MFU of 2--4 \% (the removed MBPP run, §``History: removed demos
(2026-09-30)'') is a concrete instance of that gap.} Mitigations:
\emph{colocate-with-offload} (veRL / HybridFlow, the frontier default
{[}19{]}, used by Kimi K2 {[}5{]}) vs \emph{disaggregated placement} for
heterogeneous or elastic hardware (StreamRL, up to 2.66× {[}20{]}; slime
supports both). \emph{Here the 4-train + 4-inference colocation is the
simple default; disaggregation is the lever once the sandbox or a second
cluster enters the picture.}

\textbf{Scaling execution sandboxes to tens of thousands.} Every step
runs untrusted code in an isolated sandbox, so the sandbox layer --- not
the GPU --- often becomes the bottleneck: constant create/destroy,
strong isolation, real startup cost. Kimi K2 runs more than 10,000
concurrent Kubernetes sandboxes for its coding / SWE RL {[}5{]}.
Mitigations: \emph{microVM isolation} (Firecracker-class, sub-second
boot) and \emph{managed remote code-exec backends} (E2B, Modal).
\emph{This is precisely AgentCore Runtime's per-session microVM --- the
seam this whole report builds on; the project scales it as a pool of 218
per-task Runtime ARNs with a fresh session per episode (§``The
Terminal-Bench demo'').}

\textbf{Cache-aware routing.} Rollout throughput lives on prefix
KV-cache reuse; a round-robin balancer scatters prefix-sharing requests
and collapses the hit rate --- acute in agentic RL, where a group of
trajectories shares one task prefix and multi-turn rollouts grow a long
shared prefix. Mitigations: RadixAttention plus SGLang's cache-aware
router (\textasciitilde1.9× throughput; hit rate \textasciitilde20
\%→\textasciitilde75 \%) {[}27{]}, and session-affinity routing for
multi-turn (slime \texttt{consistent\_hashing} {[}9{]}). \emph{This is
the reasoning behind §``Scaling the loop'': one cache-aware SGLang
router for inference, round-robin only on the sandbox side.}

\textbf{Fault tolerance.} One step couples training + inference + reward
+ thousands of live sandboxes, so the blast radius is large.
Mitigations: \emph{async checkpointing}; \emph{worker heartbeats}
(slime's \texttt{/health\_generate} with restart-and-re-sync,
\texttt{-\/-use-fault-tolerance} {[}9{]}); and \emph{sandbox recycling}
with timeouts and rate limits so one bad tool call does not abort a
rollout. \emph{The project's heartbeat (§``Hiding the container launch:
pre-open and heartbeat'') addresses the main sandbox-side failure mode,
idle reclamation; a GPU-instance snapshot probe for instance recycle was
tried and removed.}

\hypertarget{algorithms}{%
\subsection{Algorithms}\label{algorithms}}

\textbf{The train/inference logprob mismatch and off-policy correction.}
Even with identical parameters, the inference engine and the trainer
assign different token probabilities (differing parallelism, precision,
kernels), so ``on-policy'' RL is silently off-policy --- biasing
gradients and risking entropy collapse, and worse for long responses.
Mitigations: \emph{Truncated Importance Sampling} (clip
π\_learner/π\_sampler, now in veRL) {[}22{]}; \emph{fixing the precision
root cause} by reverting to FP16 {[}23{]}; and \emph{decoupled /
staleness-aware PPO} {[}18{]}. \emph{This project's TITO property (the
proxy records the exact sampled \texttt{token\_ids} + logprobs, no
re-tokenization --- §``Token-in-token-out'') removes the tokenization
half of the mismatch, but not the engine-vs-trainer numeric half; TIS is
the missing piece, and the importance-sampling ratio is
instrumented-but-unrecorded here.}

\textbf{GRPO and its variants.} GRPO drops PPO's critic by normalizing
advantages within a group of K sampled responses {[}13{]}, but has
documented failure modes (entropy collapse, length bias, zero-gradient
groups, MoE instability). Upgrades: \emph{DAPO} (clip-higher
ε\_high≈0.28, dynamic sampling, token-level loss, overlong reward
shaping) {[}14{]}; \emph{Dr.~GRPO} (removes the length and std
normalization that inflate wrong-answer length) {[}16{]}; \emph{GSPO} (a
sequence-level, length-normalized importance ratio that stabilizes MoE
RL and underpins Qwen3) {[}15{]}. \emph{This project trains with GRPO;
these are the drop-in improvements as runs scale up.}

\textbf{Credit assignment over long horizons.} An agentic episode is
many steps but the verifiable reward is sparse and terminal, so
trajectory-level GRPO mis-credits the actions inside a failed or
successful trajectory. Mitigations: \emph{hierarchical / turn-level
grouping} (GiGPO, +9--12 \% over GRPO on ALFWorld / WebShop) {[}26{]};
entropy-aware step attribution; and denser process rewards derived from
outcome labels alone. \emph{The Terminal-Bench demo uses a binary
end-of-episode reward with \texttt{EPISODE\_MAX\_TURNS=8}; long-horizon
families like Terminal-Bench and SWE are where turn-level credit, or
partial credit from the fraction of tests passed, becomes necessary.}

\textbf{Reward hacking and verifier trust.} Any learned reward can be
gamed, so frontier labs prefer rule-based verifiable rewards ---
DeepSeek-R1 explicitly avoids neural reward models to prevent reward
hacking {[}2{]}. But verifiers are not automatically trustworthy:
``Spurious Rewards'' shows RLVR with \emph{random} rewards still
improved a Qwen math model, because a clipping bias amplifies
pretraining behaviors rather than reward correctness --- and the effect
is model-dependent {[}24{]}. \emph{The project's rewards are binary RLVR
checkers run inside the sandbox (no LLM judge), which sidesteps
neural-RM hacking; the caution is that a benchmark gain does not by
itself prove the checker is doing the work.}

\textbf{Entropy collapse.} During long RL, entropy drops early and
exploration dies; performance is effectively ``traded from policy
entropy'' {[}25{]}. Mitigations: clip-higher {[}14{]};
covariance-targeted regularization (Clip-Cov / KL-Cov) that intervenes
only on the few high-covariance tokens driving entropy decay {[}25{]};
and KL control or removal.

\textbf{Curriculum, dynamic sampling, and the discard ratio.}
All-correct or all-wrong groups give zero advantage → wasted compute.
Mitigations: \emph{dynamic sampling} to keep only groups with 0
\textless{} correct \textless{} G (DAPO {[}14{]}) and \emph{pass-rate
difficulty curricula} (AdaRFT, up to 2× less training time {[}30{]}).
\emph{This is exactly where the project's \texttt{discard\_ratio} metric
would come alive --- it is defined (goodput depends on it) but zero in
the current runs.}

\hypertarget{takeaways}{%
\subsection{Takeaways}\label{takeaways}}

Most of these are already partly handled by the stack this report uses:
slime provides async rollout, fault-tolerant SGLang workers, and
cache-aware routing; AgentCore provides the isolated, scalable sandbox.
The open work for a production run concentrates in three places --- (1)
turning on and \emph{recording} the off-policy signals (staleness,
importance-sampling ratio, discard ratio) and adding TIS; (2) moving
from GRPO to a variant (DAPO / GSPO) as scale and response length grow;
and (3) turn-level credit assignment for the long-horizon families.
These are algorithm choices layered on top of the infrastructure the
demo already exercises.

\hypertarget{summary}{%
\section{Summary}\label{summary}}

\begin{itemize}
\tightlist
\item
  \textbf{AgentCore Runtime is a workable RL sandbox for per-task
  environments.} A Runtime is a deployed image, a session is one
  microVM, and a fresh session per episode gives each episode its own
  kernel with no state carried over. The arm64 and 2 GB image limits are
  hard, and they decide which tasks fit.
\item
  \textbf{slime needs no fork.} One generate hook, a sandbox client (one
  \texttt{InvokeAgentRuntime} per operation), and the mini\_claude loop
  running inside the session, dialling the trainer's adapter back over
  the VPC, are the whole integration. Token-in-token-out keeps the
  trained tokens exactly the sampled ones.
\item
  \textbf{Terminal-Bench is the demo.} 218 Runtimes are built from the
  tasks' own Dockerfiles, 166 of them oracle-validated. Rows are routed
  strictly to their Runtime. Look-ahead pre-open takes the cold start
  off the critical path (open wait p95 0 at the cost of
  \textasciitilde95 \% idle pre-opened time), and the heartbeat keeps
  sessions alive across the weight-sync stall. The harness tier scales
  the agent loop with the sandbox.
\item
  \textbf{What is open.} Reward at 4B is sparse (4/382 solved on
  2026-09-23; 23/449 in the clean 2026-09-30 re-run, eval 1/62).
  Convergence run 20, plus a curriculum or partial-credit reward, is the
  next step.
\end{itemize}

\hypertarget{references-papers-models}{%
\section{References --- papers \&
models}\label{references-papers-models}}

\footnotesize

\textbf{Agentic RL \& models.} {[}1{]} Agent Lightning --- Harnessed
Agentic RL. \url{https://arxiv.org/abs/2508.03680} · {[}2{]}
DeepSeek-R1. \url{https://arxiv.org/abs/2501.12948} · {[}3{]} Qwen3
Technical Report. \url{https://arxiv.org/abs/2505.09388} · {[}4{]}
Llama-Nemotron. \url{https://arxiv.org/abs/2505.00949} · {[}5{]} Kimi
K2. \url{https://arxiv.org/abs/2507.20534} · {[}6{]} SWE-bench.
\url{https://arxiv.org/abs/2310.06770} · {[}6b{]} SWE-bench README
(arm64 ``experimental''; \texttt{-\/-task-repo} local Buildx build).
\url{https://github.com/SWE-bench/SWE-bench} · {[}6c{]} Terminal-Bench.
\url{https://github.com/laude-institute/terminal-bench} · {[}7{]}
KernelBench.
\url{https://huggingface.co/datasets/ScalingIntelligence/KernelBench} ·
{[}8{]} MBPP. \url{https://arxiv.org/abs/2108.07732} · {[}9{]} slime.
\url{https://github.com/THUDM/slime} · {[}10{]} AceReason-Nemotron:
Advancing Math and Code Reasoning through RL.
\url{https://arxiv.org/abs/2505.16400} · {[}11{]} LiveCodeBench.
\url{https://arxiv.org/abs/2403.07974} · {[}12{]} NVIDIA Nemotron
Post-Training Dataset v2.
\url{https://huggingface.co/datasets/nvidia/Nemotron-Post-Training-Dataset-v2}

\hypertarget{references-agentcore-docs-12}{%
\section{References --- AgentCore docs
(1/2)}\label{references-agentcore-docs-12}}

\scriptsize

\textbf{Official AgentCore docs.} {[}D1{]} What is AgentCore.
\url{https://docs.aws.amazon.com/bedrock-agentcore/latest/devguide/what-is-bedrock-agentcore.html}
· {[}D2{]} Host with Runtime (isolation, compute types).
\url{https://docs.aws.amazon.com/bedrock-agentcore/latest/devguide/agents-tools-runtime.html}
· {[}D3{]} Runtime how it works (idle 15 m / max 8 h).
\url{https://docs.aws.amazon.com/bedrock-agentcore/latest/devguide/runtime-how-it-works.html}
· {[}D4{]} Lifecycle settings.
\url{https://docs.aws.amazon.com/bedrock-agentcore/latest/devguide/runtime-lifecycle-settings.html}
· {[}D5{]} Instances (GPU).
\url{https://docs.aws.amazon.com/bedrock-agentcore/latest/devguide/runtime-instances-how-it-works.html}
· {[}D6{]} Code Interpreter.
\url{https://docs.aws.amazon.com/bedrock-agentcore/latest/devguide/code-interpreter-tool.html}
· {[}D7{]} Code Interpreter session.
\url{https://docs.aws.amazon.com/bedrock-agentcore/latest/devguide/code-interpreter-session-characteristics.html}

\hypertarget{references-agentcore-docs-22}{%
\section{References --- AgentCore docs
(2/2)}\label{references-agentcore-docs-22}}

\scriptsize

{[}D8{]} InvokeAgentRuntime.
\url{https://docs.aws.amazon.com/bedrock-agentcore/latest/APIReference/API_InvokeAgentRuntime.html}
· {[}D9{]} CreateAgentRuntime.
\url{https://docs.aws.amazon.com/bedrock-agentcore-control/latest/APIReference/API_CreateAgentRuntime.html}
· {[}D10{]} NetworkConfiguration.
\url{https://docs.aws.amazon.com/bedrock-agentcore-control/latest/APIReference/API_NetworkConfiguration.html}
· {[}D11{]} VPC interface endpoints.
\url{https://docs.aws.amazon.com/bedrock-agentcore/latest/devguide/vpc-interface-endpoints.html}
· {[}D12{]} StopRuntimeSession.
\url{https://docs.aws.amazon.com/bedrock-agentcore/latest/APIReference/API_StopRuntimeSession.html}
· {[}D13{]} Observability.
\url{https://docs.aws.amazon.com/bedrock-agentcore/latest/devguide/observability.html}
· {[}D17{]} Runtime service contract (port 8080, /invocations, /ping).
\url{https://docs.aws.amazon.com/bedrock-agentcore/latest/devguide/runtime-service-contract.html}

\hypertarget{references-agentcore-sdk-cli-samples}{%
\section{References --- AgentCore SDK, CLI \&
samples}\label{references-agentcore-sdk-cli-samples}}

\scriptsize

\textbf{AWS code, cross-checked against the docs.} {[}D14{]} AgentCore
samples (\texttt{awslabs/agentcore-samples}).
\url{https://github.com/awslabs/agentcore-samples} · {[}D15{]} AgentCore
Python SDK (\texttt{bedrock-agentcore}; \texttt{BedrockAgentCoreApp} /
\texttt{@app.entrypoint}).
\url{https://github.com/aws/bedrock-agentcore-sdk-python} · {[}D16{]}
AgentCore CLI (\texttt{@aws/agentcore}, recommended)
\url{https://www.npmjs.com/package/@aws/agentcore} · legacy starter
toolkit \url{https://github.com/aws/bedrock-agentcore-starter-toolkit}

\hypertarget{references-production-agentic-rl-open-problems}{%
\section{References --- production agentic RL (open
problems)}\label{references-production-agentic-rl-open-problems}}

\scriptsize

\textbf{RL algorithms \& systems.} {[}13{]} DeepSeekMath (GRPO).
\url{https://arxiv.org/abs/2402.03300} · {[}14{]} DAPO.
\url{https://arxiv.org/abs/2503.14476} · {[}15{]} GSPO (Group Sequence
Policy Optimization). \url{https://arxiv.org/abs/2507.18071} · {[}16{]}
Understanding R1-Zero-Like Training (Dr.~GRPO).
\url{https://arxiv.org/abs/2503.20783} · {[}17{]} Kimi k1.5.
\url{https://arxiv.org/abs/2501.12599} · {[}18{]} AReaL (asynchronous RL
system). \url{https://arxiv.org/abs/2505.24298} · {[}19{]} HybridFlow
(veRL). \url{https://arxiv.org/abs/2409.19256} · {[}20{]} StreamRL
(disaggregated stream generation).
\url{https://arxiv.org/abs/2504.15930} · {[}21{]} Asynchronous RLHF.
\url{https://arxiv.org/abs/2410.18252} · {[}22{]} Truncated Importance
Sampling / train--inference mismatch (Yao et al.).
\url{https://fengyao.notion.site/off-policy-rl} · {[}23{]} Defeating the
Training--Inference Mismatch via FP16.
\url{https://arxiv.org/abs/2510.26788} · {[}24{]} Spurious Rewards.
\url{https://arxiv.org/abs/2506.10947} · {[}25{]} The Entropy Mechanism
of RL. \url{https://arxiv.org/abs/2505.22617} · {[}26{]} GiGPO
(hierarchical group advantages). \url{https://arxiv.org/abs/2505.10978}
· {[}27{]} SGLang / RadixAttention.
\url{https://arxiv.org/abs/2312.07104} --- cache-aware router:
\url{https://lmsys.org/blog/2024-12-04-sglang-v0-4/} · {[}28{]} APRIL
(active partial rollouts). \url{https://arxiv.org/abs/2509.18521} ·
{[}29{]} Moonshot checkpoint-engine.
\url{https://github.com/MoonshotAI/checkpoint-engine} · {[}30{]} AdaRFT
(difficulty curriculum). \url{https://arxiv.org/abs/2504.05520} ·
{[}31{]} DSec: DeepSeek Elastic Compute, a sandbox platform for agentic
RL. \url{https://arxiv.org/abs/2609.22978}

\end{document}
